% Incognita paper P004. Written from templates/paper/paper.tex (AGENTS.md, section 5; docs/writing.md,
% sections 2-2a, the owner's house style). Two builds of one source (AGENTS.md 5.2), after
% `python3 scripts/certify.py bib` (papers/library.bib):
%   reader:    latexmk -pdf -interaction=nonstopmode -halt-on-error -cd papers/<P>/paper.tex
%   blueprint: latexmk -pdf -interaction=nonstopmode -halt-on-error -jobname=paper-blueprint
%                -usepretex='\def\IncognitaBlueprint{}' -cd papers/<P>/paper.tex
% Blueprint conventions: every node has a \label (def:, lem:, prop:, thm:, cor:, cite:); statements
% start with \uses when they use definitions; every proof starts with an exhaustive \uses; a proof away
% from its statement has \proves. What serves only the referees or the owner is in DOSSIER.md. Never
% write \lean, \leanok or \mathlibok in this file.
\documentclass[11pt]{article}
\usepackage[T1]{fontenc}
\usepackage[utf8]{inputenc}
\usepackage{lmodern}
\usepackage[a4paper,margin=2.5cm]{geometry}
\usepackage{amsmath,amssymb,mathtools}
\usepackage{booktabs}
\usepackage{algorithm}
\usepackage{algpseudocode}
\usepackage{enumitem}
\usepackage{xcolor}
\usepackage[round]{natbib}
\usepackage[colorlinks=true,linkcolor=blue!50!black,citecolor=blue!50!black,urlcolor=blue!50!black]{hyperref}
\usepackage{../../templates/paper/blueprint-macros}

% ---------------------------------------------------------------- macros (\newcommand only)
% Each symbol has one meaning in the whole paper.
\newcommand{\E}{\mathbb{E}}
\newcommand{\Prob}{\mathbb{P}}
\newcommand{\R}{\mathbb{R}}
\newcommand{\ind}{\mathbf{1}}
\newcommand{\argmin}{\operatorname*{arg\,min}}
% Navigation (house style): an anchor \dt{s:...} at each definition, a macro per symbol linking to it.
\newcommand{\dt}[1]{\hypertarget{#1}{}}
\newcommand{\sy}[2]{{\hyperlink{#1}{#2}}}
\newcommand{\term}[2]{\hyperlink{#1}{#2}}
\newcommand{\imed}{\hyperref[alg:imed]{\texttt{IMED}}}
\newcommand{\Sm}{\sy{s:S}{\mathcal{S}}}
\newcommand{\Mm}{\sy{s:M}{\mathcal{M}}}
\newcommand{\mean}{\sy{s:mean}{m}}
\newcommand{\arms}{\sy{s:arms}{\alpha}}
\newcommand{\Karms}{\sy{s:K}{K}}
\newcommand{\marm}{\sy{s:ma}{m}}
\newcommand{\ms}{\sy{s:mstar}{m^{\star}}}
\newcommand{\gap}{\sy{s:gap}{\Delta}}
\newcommand{\KL}{\sy{s:KL}{\mathrm{KL}}}
\newcommand{\Kinf}{\sy{s:kinf}{\mathcal{K}_{\inf}}}
\newcommand{\hP}{\sy{s:hP}{\hat{P}}}
\newcommand{\bx}{\sy{s:barx}{\bar{x}}}
\newcommand{\bX}{\sy{s:barx}{\bar{X}}}
\newcommand{\Xs}{\sy{s:X}{X}}
\newcommand{\Useed}{\sy{s:U}{U}}
\newcommand{\Ap}{\sy{s:A}{A}}
\newcommand{\Ob}{\sy{s:O}{O}}
\newcommand{\Np}{\sy{s:N}{N}}
\newcommand{\Reg}{\sy{s:R}{R}}
\newcommand{\hnu}{\sy{s:hnu}{\hat{\nu}}}
\newcommand{\hm}{\sy{s:hm}{\hat{m}}}
\newcommand{\hnut}{\sy{s:hnut}{\hat{\nu}}}
\newcommand{\hmt}{\sy{s:hmt}{\hat{m}}}
\newcommand{\hms}{\sy{s:hms}{\hat{m}^{\star}}}
\newcommand{\Dv}{\sy{s:D}{D}}
\newcommand{\Ix}{\sy{s:I}{I}}
\newcommand{\Is}{\sy{s:Is}{I^{\star}}}
\newcommand{\ws}{\sy{s:w}{w}}
\newcommand{\Gam}{\sy{s:Gamma}{\Gamma}}
\newcommand{\lst}{\sy{s:lst}{\lambda^{\star}}}
\newcommand{\Lt}{\sy{s:Lt}{\tilde{\Lambda}}}
\newcommand{\Ki}{\sy{s:Ki}{\mathcal{K}}}
\newcommand{\Pq}{\sy{s:Pq}{P}}
\newcommand{\kl}{\sy{s:kl}{\mathrm{kl}}}
\newcommand{\lmax}{\sy{s:lmax}{\bar{\lambda}}}
\newcommand{\ui}{\sy{s:ui}{u}}
\newcommand{\MT}{\sy{s:MT}{M}}
\newcommand{\TA}{\sy{s:TA}{\mathsf{A}}}
\newcommand{\TB}{\sy{s:TB}{\mathsf{B}}}
\newcommand{\TC}{\sy{s:TC}{\mathsf{C}}}

\title{Finite variance suffices for the asymptotic optimality of IMED on semi-bounded rewards}
\author{Claude Opus 5.5, Incognita agents\thanks{Written by AI agents of the Incognita project
  (\url{https://github.com/prevostad/Incognita}) and refereed by two AI agents; its certification
  is under the date.}}
\date{2026-11-16}

\begin{document}
\maketitle

\begin{abstract}
We study stochastic multi-armed bandits with semi-bounded rewards: every arm has a law on
$(-\infty,1]$ with a finite mean, and the learner knows no lower bound on the rewards and no
moment. The IMED algorithm was known to attain the Burnetas--Katehakis lower bound in this model
when one optimal arm has a moment generating function near the origin, and it was conjectured that
finite moments of order $2+p$ suffice. We prove that, when the best mean is below $1$, a finite
variance of one optimal arm suffices, with only finite means for the other arms: for every
suboptimal arm $i$, $\limsup_{T}\E N_i(T)/\log T\le 1/\mathcal{K}_{\inf}(\nu_i,m^{\star})$, as a
consequence of a finite-time bound. Hence IMED is consistent and asymptotically optimal on the
semi-bounded laws with a finite moment of order $r$, other than the point mass at $1$, for every
$r\ge 2$. On the scale of power-law lower tails, the exponent two is the threshold: if the lower
tail of the optimal arm is at least $\kappa y^{-q}$ with $1<q<2$, IMED pulls a deterministic
suboptimal arm at least of order $T^{2-q}$ times, so it is not consistent; for tails of order
$y^{-2}$, it misses the lower-bound constant by at least $\kappa/(4e)$.
\end{abstract}

% ================================================================ main text
\section{Introduction}\label{sec:intro}

In the stochastic multi-armed bandit problem, a learner pulls one of $K$ arms at each round and
observes a reward drawn independently from the law of that arm. It wants to pull the suboptimal
arms, those whose mean is not the largest, as rarely as possible. When the learner knows a model
$\mathcal D$ that contains the law of every arm, \citet{LR85} and \citet{BK96} showed that every
algorithm that is consistent on $\mathcal D^K$ pulls each suboptimal arm $i$, in expectation over
$T$ rounds, at least
\begin{equation}\label{eq:intro-lb}
  (1-o(1))\frac{\log T}{\mathcal K_{\inf}(\nu_i,m^\star;\mathcal D)}\quad\text{times},\qquad
  \mathcal K_{\inf}(P,\mu;\mathcal D):=\inf\{\mathrm{KL}(P,Q):Q\in\mathcal D,\ m(Q)>\mu\},
\end{equation}
where $\nu_i$ is the law of arm $i$, $m^\star$ the best mean and $m(Q)$ the mean of $Q$ (we use the
form of \citet[Theorem~1]{GMS19}). When an algorithm is consistent on a model and its counts
satisfy the matching upper bound on every instance, we say that it is asymptotically optimal on that
model; the phrase is only a shorthand for this pair of bounds.

This paper is about the \emph{semi-bounded} model: the laws on $(-\infty,1]$ with a finite mean.
The rewards are bounded above by a known value, normalised to $1$, but no lower bound is known; a
reward may be the opposite of a delay or of a loss that has no known cap
\citep[Section~1.2]{HT15}. \citet{HT15} showed that $\mathcal K_{\inf}$ has a one-dimensional dual on
this model, and introduced the Indexed Minimum Empirical Divergence algorithm (\texttt{IMED}), which
pulls an arm minimising $N_a\mathcal K_{\inf}(\hat\nu_a,\hat m^\star)+\log N_a$, where $N_a$ is the
number of pulls of arm $a$, $\hat\nu_a$ its empirical law and $\hat m^\star$ the largest empirical
mean. They proved that \texttt{IMED} is asymptotically optimal on this model when the best mean is
below $1$ and \emph{one optimal arm has a moment generating function near the origin}
\citep[Theorem~3 and Corollary~4]{HT15}; for a law on $(-\infty,1]$, this is an exponential moment
of the lower tail. They conjectured that finite moments of order $2+p$ should suffice, and left the
question open \citep[Section~9]{HT15}.

We answer this question in both directions. If the best mean $m^\star$ is below $1$ and one optimal
arm has a finite variance, with only finite means for the other arms, then \texttt{IMED} satisfies,
for every suboptimal arm $i$,
\begin{equation}\label{eq:intro-main}
  \limsup_{T\to\infty}\frac{\E N_i(T)}{\log T}\le\frac{1}{\mathcal K_{\inf}(\nu_i,m^\star)},
\end{equation}
as a consequence of a finite-time bound; with the lower bound \eqref{eq:intro-lb}, \texttt{IMED} is
asymptotically optimal on the semi-bounded laws with a finite moment of order $r$, other than the
point mass at $1$, for every $r\ge2$. Conversely, on the scale of power-law lower tails, the
exponent two cannot be lowered: if the lower tail of the optimal arm is at least $\kappa y^{-q}$ with
$1<q<2$, \texttt{IMED} pulls a deterministic suboptimal arm at least of order $T^{2-q}$ times, and for
tails of order $y^{-2}$ the $\liminf$ of $\E N_i(T)/\log T$ exceeds the lower-bound constant by at least
$\kappa/(4e)$.

The analysis of \citet{HT15} bounds the rounds at which every empirical mean is low by a \emph{sum}
over the sample sizes of the optimal arm, whose expectation is infinite without a third moment. We
count the same rounds by a \emph{supremum}, and bound its expectation under a finite variance by a
one-big-jump decomposition of the dual of $\mathcal K_{\inf}$. The negative result comes from the
reference level of \texttt{IMED}, an empirical mean: one large negative reward of the optimal arm
hands the lead to a suboptimal arm for a number of rounds linear in that reward.

\paragraph{Our contributions} We summarise our main contributions below.
\begin{enumerate}[label=(\roman*),leftmargin=2.2em]
  \item \emph{Regret upper bound under a finite variance.} A finite-time bound on the expected
    number of pulls of each suboptimal arm by \texttt{IMED}, when one optimal arm has a finite
    variance and the other arms a finite mean (Theorem~\ref{thm:main}).
  \item \emph{Asymptotic optimality.} \texttt{IMED} is consistent and attains the lower-bound
    constant on the semi-bounded laws with a finite moment of order $r$, other than the point mass
    at $1$, for every $r\ge2$; this contains the class of the conjecture of \citet{HT15}
    (Corollary~\ref{cor:optimal}).
  \item \emph{Lower bounds for heavier lower tails.} With a lower tail $\kappa y^{-q}$, $1<q<2$,
    \texttt{IMED} is not consistent; with a tail of order $y^{-2}$, it misses the lower-bound
    constant (Theorem~\ref{thm:counter}, Corollary~\ref{cor:pareto}).
  \item \emph{Tools.} A pathwise count of the rounds at which every empirical mean is low by a
    supremum over the sample sizes of one arm (Lemma~\ref{lem:sup}); the finiteness of the
    expectation of that supremum under a finite variance (Proposition~\ref{prop:key}); the equality
    of $\mathcal K_{\inf}$ on the semi-bounded laws and on those with a finite moment of order $r$
    (Lemma~\ref{lem:class}).
  \item \emph{Examples.} An explicit instance satisfying the hypotheses of Theorem~\ref{thm:main}
    and not those of \citet[Theorem~3]{HT15}, the degenerate cases, and simulations of the three
    regimes (Section~\ref{sec:examples}).
\end{enumerate}

\paragraph{Outline} Section~\ref{sec:related} discusses related work. Section~\ref{sec:setting}
formalises the setting. Section~\ref{sec:algorithm} recalls \texttt{IMED} and fixes its convention
for empirical leaders. Section~\ref{sec:upper} states the regret upper bound,
Section~\ref{sec:optimal} the asymptotic optimality, and Section~\ref{sec:counter} the lower bounds
for heavier tails. Section~\ref{sec:overview} explains how the proofs work,
Section~\ref{sec:examples} gives examples and simulations, and Section~\ref{sec:open} concludes and
states open questions. The proofs are in the appendices.

% ----------------------------------------------------------------
\section{Related work and context}\label{sec:related}

Table~\ref{tab:comparison} compares the bounds closest to ours. We discuss each line of work below.

\paragraph{Asymptotic optimality in stochastic bandits} The lower bound \eqref{eq:intro-lb} is due to
\citet{LR85} for one-parameter families and to \citet{BK96} for general models. We use the form of
\citet[Theorem~1]{GMS19}, for learners that are uniformly fast convergent on the model, recalled as
Theorem~\ref{cite:GMS19-thm1}; the textbook \citet[Chapter~16]{LS} gives its form for consistent
learners.

\paragraph{Semi-bounded rewards and IMED} For laws with a bounded support, \citet{HT10} introduced
DMED and the dual formula of $\mathcal K_{\inf}$. \citet{HT15} extended the dual formula to the
semi-bounded model \citep[Theorem~2]{HT15}, introduced \texttt{IMED}, and proved a finite-time bound
when one optimal arm $j$ has a moment generating function near $0$ \citep[Theorem~3]{HT15},
recalled as Theorem~\ref{cite:HT15-thm3}. They deduce a remainder $O((\log T)^{10/11})$
\citep[Corollary~4]{HT15}, and a remainder $O(1)$ when every arm has a moment generating function
\citep[Theorem~5]{HT15}.

\paragraph{Heavy-tailed bandits} \citet[Theorem~1]{BCL13} obtain a logarithmic regret for
Robust-UCB with a known bound on a moment of order $1+\varepsilon$, with a constant polynomial in the
ratio of that bound to the gaps, not $1/\mathcal K_{\inf}$. \citet[Corollary~2]{AJK21} prove that
KLinf-UCB is asymptotically optimal on the class of laws on $\R$ with $\E|X|^{1+\varepsilon}\le B$,
for known $\varepsilon$ and $B$. \citet[Section~4.4]{BSH23} prove that MED, a randomised relative of
\texttt{IMED}, is asymptotically optimal on the laws with $\E h(|X|)<B$ (and a little more), for a
convex function $h$ growing faster than $|x|^{1+\eta}$ and a bound $B$, both known to the learner;
a known bound on the variance is an example.

\paragraph{Optimality from generic concentration assumptions} \citet[Theorem~3.1]{PA26} prove that
KLinf-UCB, with an exploration function $\log t+2\log\log t+g(N_a)$, satisfies
$\limsup_T\E N_a(T)/\log T\le1/\mathrm{KL}_{\inf}(\nu_a,\mu_1)$ on every class of laws whose
empirical $\mathrm{KL}_{\inf}$ satisfies two assumptions: a time-uniform bound on its upper
deviations at the mean, with $g(n)=C_1\log(1+n)+C_2$ \citep[Assumption~2.3]{PA26}, and an
exponential bound, quadratic in the deviation, on its lower deviations at a level above the mean
\citep[Assumption~2.4]{PA26}. They check them for bounded supports and for known moment bounds;
semi-bounded laws are not treated. We check them on the semi-bounded laws with a finite mean
(Lemma~\ref{lem:pa-assumptions}, and the paragraph ``Another algorithm'' of
Section~\ref{sec:counter}).

\paragraph{Closest related work} \citet[Theorem~3]{HT15} is the closest result: same model, same
algorithm, same criterion. Its last term depends on Cram\'er's rate function of the lower tail of
the empirical mean of arm $j$, which is positive only under an exponential moment; for a lower tail
of Pareto type the bound is infinite, so it does not give \eqref{eq:intro-main}. Its proof cannot be
adapted by changing a concentration inequality: it bounds the rounds at which every empirical mean is
low by a sum over the sample sizes of the optimal arm, and the expectation of that sum is infinite
unless the optimal arm has a finite third moment (Section~\ref{sec:overview-mgf}).
\citet[Corollary~2]{AJK21} is a different model (two-sided laws, a known moment bound), with a
different constant and an algorithm that uses the bound; it says nothing about \texttt{IMED}, which
knows no moment. \citet[Section~4.4]{BSH23} likewise assume a known moment bound and analyse MED,
not \texttt{IMED}. \citet[Theorem~1]{BCL13} does not give a constant $1/\mathcal K_{\inf}$.
\citet[Theorem~3.1]{PA26} is about KLinf-UCB, which inverts $\mathcal K_{\inf}$ to compute an upper
confidence bound, and its assumptions are not checked in that paper for semi-bounded laws; it says
nothing about \texttt{IMED}.

\paragraph{Gap with the literature} We found no published analysis of \texttt{IMED} that attains the
constant of \eqref{eq:intro-lb} on semi-bounded laws without an exponential moment, and no result
on \texttt{IMED} under polynomial lower tails. For another algorithm, KLinf-UCB, the general theorem
of \citet[Theorem~3.1]{PA26} reduces the question to two concentration assumptions, which no
published work checks on semi-bounded laws; Lemma~\ref{lem:pa-assumptions} checks them with a
finite mean only, by the published method of \citet[Proposition~5]{AJK21} (see also
\citealp[Appendix~A.1]{PA26}) in a one-dimensional, semi-bounded form. The decomposition of the count of a suboptimal arm, the counting lemma for the
minimal index, the dual formula, the bound on the deviations of the suboptimal arm and the grid on
the dual variable come from \citet{HT15}. What is new is to apply their counting lemma once, at a
random level, so that a supremum replaces a sum; the bound on the expectation of that supremum by
big negative jumps and a good part, a split classical for sums of heavy-tailed variables (see also
\citealp{FN71,Nag79}), here inside the dual product of $\mathcal K_{\inf}$, where one big jump costs a
linear factor and the second moment enters in one term only; and the frozen-index argument behind
the lower bounds. Appendix~\ref{sec:attribution} credits each technique.

\begin{table}[t]
\centering
\caption{Bounds for semi-bounded or heavy-tailed rewards. ``Both'': the upper bound matches the
lower bound of the same model at first order. $\mathcal K_{\inf}$ is taken in the model of the row.}
\label{tab:comparison}
\small
\begin{tabular}{@{}p{0.43\textwidth}p{0.43\textwidth}l@{}}
\toprule
Reference: setting and assumptions & Algorithm: bound & Both \\
\midrule
\citet[Thm.~3, Cor.~4]{HT15}: laws on $(-\infty,1]$, $m^\star<1$, an exponential moment for one optimal
arm & \texttt{IMED}: $\E N_i(T)\le\log T/\mathcal K_{\inf}(\nu_i,m^\star)+O((\log T)^{10/11})$ & yes \\
\citet[Thm.~5]{HT15}: same, an exponential moment for every arm & \texttt{IMED}: regret
$\sum_i\Delta_i\log T/\mathcal K_{\inf}(\nu_i,m^\star)+O(1)$ & yes \\
\citet[Cor.~2]{AJK21}: laws on $\R$ with $\E|X|^{1+\varepsilon}\le B$, with $\varepsilon$ and $B$
known & KLinf-UCB with a parameter $\tilde\varepsilon\ge0$:
$\limsup_T\E N_i(T)/\log T\le(1+\tilde\varepsilon)/\mathcal K_{\inf}(\nu_i,m^\star)$ & if $\tilde\varepsilon=0$ \\
\citet[Sec.~4.4]{BSH23}: laws on $\R$ with $\E h(|X|)<B$, with $h$ and $B$ known & MED:
asymptotically optimal on that class & yes \\
\citet[Thm.~3.1]{PA26}: any class satisfying their Assumptions~2.3--2.4 (checked there for bounded
supports and known moment bounds) & KLinf-UCB: $\limsup_T\E N_a(T)/\log T\le1/\mathrm{KL}_{\inf}(\nu_a,\mu_1)$
& yes \\
\citet[Thm.~1]{BCL13}: laws on $\R$, a known bound on $\E|X|^{1+\varepsilon}$ & Robust-UCB:
logarithmic regret, constant not $1/\mathcal K_{\inf}$ & no \\
Theorem~\ref{thm:main}, Corollary~\ref{cor:optimal}: laws on $(-\infty,1]$, $m^\star<1$, a finite
variance for one optimal arm & \texttt{IMED}: $\limsup_T\E N_i(T)/\log T\le1/\mathcal K_{\inf}(\nu_i,m^\star)$, from a
finite-time bound & yes \\
Theorem~\ref{thm:counter}: two arms, lower tail of the optimal arm $\ge\kappa y^{-q}$, $1<q<2$ &
\texttt{IMED}: $\E N_2(T)\ge cT^{2-q}$, not consistent & \\
\bottomrule
\end{tabular}
\end{table}

% ----------------------------------------------------------------
\section{Problem formulation}\label{sec:setting}

In this section we introduce the setting: first the classical stochastic bandit with a model of
laws known to the \term{w:learner}{learner}, then what is specific to semi-bounded rewards.

\subsection{Classical setting}

A \term{w:learner}{learner} pulls one arm at each round and observes a reward drawn from the law of that arm, given
the past. It knows a model $\mathcal D$, a set of laws that contains the law of every arm, and no
horizon. The definitions below are the usual ones, written so that the samples of each arm form an
i.i.d.\ sequence; the only objects specific to this paper are the classes of semi-bounded laws
$\Sm$, $\Sm_r$ and $\Mm_r$, which we introduce with the \term{w:instance}{instances}.

\begin{definition}[Semi-bounded laws and instances]\label{def:model}
\dt{s:S}Let $\Sm$ be the set of probability laws $P$ on $\R$ with $P((-\infty,1])=1$ and
$\E_P|X|<\infty$. \dt{s:mean}The mean of $P$ is $\mean(P)\coloneqq\E_PX$. \dt{s:M}For $r\ge1$, let
\begin{equation*}
  \Sm_r\coloneqq\{P\in\Sm:\E_P|X|^r<\infty\},\qquad\Mm_r\coloneqq\Sm_r\setminus\{\delta_1\},
\end{equation*}
where $\delta_x$ is the Dirac mass at $x$. \dt{s:arms}The arms form a finite set $\arms$.
\dt{s:K}The number of arms is $\Karms\coloneqq|\arms|\ge2$. \dt{w:instance}An \emph{instance} is
$\nu=(\nu_a)_{a\in\arms}\in\Sm^\arms$. \dt{s:ma}The mean of arm $a$ is $\marm_a\coloneqq \mean(\nu_a)$.
\dt{s:mstar}The best mean is $\ms\coloneqq\max_a\marm_a$. \dt{s:gap}The gap of arm $a$ is
$\gap_a\coloneqq\ms-\marm_a$. \dt{w:optimal}An arm is \emph{optimal} if $\gap_a=0$
\dt{w:subopt}and \emph{suboptimal} if $\gap_a>0$. For a set $\mathcal D$ of laws, $\mathcal D^\arms$
is the set of \term{w:instance}{instances} whose laws all belong to $\mathcal D$.
\end{definition}

\begin{definition}[The divergence $\Kinf$ and empirical laws]\label{def:kinf}
\uses{def:model}
\dt{s:KL}$\KL(P,Q)\coloneqq\int\log\frac{dP}{dQ}\,dP$ if $P$ is absolutely continuous with respect to $Q$,
and $+\infty$ otherwise. For a set $\mathcal D$ of laws with finite means, a law $P$ and
$\mu\in\R$,
\begin{equation*}
  \dt{s:kinf}\Kinf(P,\mu;\mathcal D)\coloneqq\inf\{\KL(P,Q):Q\in\mathcal D,\ \mean(Q)>\mu\}\in[0,\infty],
  \qquad\inf\emptyset\coloneqq+\infty,
\end{equation*}
and $\Kinf(P,\mu)\coloneqq\Kinf(P,\mu;\Sm)$. \dt{s:hP}For reals $x_1,\dots,x_s$ ($s\ge1$), let
$\hP[x_{1:s}]\coloneqq\frac1s\sum_{l=1}^s\delta_{x_l}$ be their empirical law. \dt{s:barx}Their mean
is $\bx_s\coloneqq\frac1s\sum_{l=1}^sx_l$.
\end{definition}

\begin{definition}[Runs, learners and regret]\label{def:run}
\uses{def:model, def:kinf}
Fix an \term{w:instance}{instance} $\nu\in\Sm^\arms$. \dt{s:X}On a probability space $(\Omega,\mathcal F,\Prob)$, let
$(\Xs_{a,s})_{a\in\arms,\,s\ge1}$ be independent random variables with $\Xs_{a,s}\sim\nu_a$,
\dt{s:U}and let $\Useed$ be uniform on $[0,1]$ and independent of them. \dt{w:learner}A \emph{learner} is a
sequence of measurable maps $(\pi_t)_{t\ge0}$. \dt{s:A}At round $t=0,1,2,\dots$ it pulls the arm
\begin{equation*}
  \Ap_t\coloneqq\pi_t(\Useed,\Ap_0,\Ob_0,\dots,\Ap_{t-1},\Ob_{t-1})\in\arms
\end{equation*}
\dt{s:O}and observes $\Ob_t\coloneqq \Xs_{\Ap_t,\Np_{\Ap_t}(t)+1}$, where
\begin{equation*}
  \dt{s:N}\Np_a(t)\coloneqq|\{u<t:\Ap_u=a\}|
\end{equation*}
is the number of pulls of arm $a$ before round $t$. The horizon $T\ge1$ consists of the rounds
$0,\dots,T-1$, and $\Np_a(T)$ is the number of pulls of $a$ among them. The regret is
\begin{equation*}
  \dt{s:R}\Reg_T\coloneqq\sum_{a\in\arms}\gap_a\E \Np_a(T).
\end{equation*}
\dt{s:hnu}For an arm $a$ and $s\ge1$, $\hnu_{a,s}\coloneqq\hP[\Xs_{a,1:s}]$, \dt{s:hm}and $\hm_{a,s}$ is its mean.
\end{definition}

Given the past, $\Ob_t$ has the law $\nu_{\Ap_t}$: this is the usual bandit model, written so that
the samples of each arm form an i.i.d.\ sequence. \dt{w:anytime}Every \term{w:learner}{learner} here is
\emph{anytime}: it does not know $T$.

\begin{definition}[Consistency]\label{def:classes}
\uses{def:model, def:run}
Let $\mathcal D\subseteq\Sm$. \dt{w:consistent}A \term{w:learner}{learner} is \emph{consistent} on $\mathcal D$ if $\Reg_T=o(T^p)$ for
every $p>0$ and every $\nu\in\mathcal D^\arms$ \citep[Definition~16.1]{LS}. \dt{w:ufc}It is \emph{uniformly fast
convergent} on $\mathcal D$ if $\E \Np_a(T)=o(T^p)$ for every $p\in(0,1]$, every $\nu\in\mathcal D^\arms$
and every \term{w:subopt}{suboptimal} arm $a$ \citep[Definition~1]{GMS19}.
\end{definition}

A \term{w:consistent}{consistent} \term{w:learner}{learner} is \term{w:ufc}{uniformly fast convergent}, because $\gap_a\E \Np_a(T)\le \Reg_T$ and
$\gap_a>0$ for a \term{w:subopt}{suboptimal} arm.

\subsection{Semi-bounded rewards}

A law of $\Sm$ with mean $1$ is $\delta_1$, so
\begin{equation*}
  \ms<1\quad\text{if and only if no arm has the law }\delta_1;
\end{equation*}
in particular every \term{w:instance}{instance} of $\Mm_r^\arms$ has $\ms<1$.

The inequality $\mean(Q)>\mu$ in Definition~\ref{def:kinf} is strict. For $\mu\ge1$ no law of $\Sm$ has a
mean above $\mu$, so $\Kinf(P,\mu)=+\infty$. For $\mu<1$ and $P\in\Sm$, \citet{HT15} proved the dual
formula
\begin{equation}\label{eq:dual}
  \Kinf(P,\mu)=\max_{0\le\lambda\le1/(1-\mu)}\E_P\big[\log\big(1-\lambda(X-\mu)\big)\big]
\end{equation}
(Lemma~\ref{cite:HT15-thm2}). The properties of $\Kinf$ that we use are collected in
Appendix~\ref{sec:dual}.

\paragraph{Notational convention} Rounds are numbered from $0$, and $\Np_a(T)$ counts the pulls in
the rounds $0,\dots,T-1$. $\log$ is the natural logarithm; $x^\pm\coloneqq\max\{\pm x,0\}$;
$\ind\{E\}=\ind_E$ is the indicator of $E$; $\lfloor x\rfloor$ is the integer part; $\E_P$ and
$\Prob_P$ refer to a generic $X\sim P$. Appendix~\ref{sec:notation} lists the symbols.

% ----------------------------------------------------------------
\section{The algorithm IMED}\label{sec:algorithm}

In this section we recall \imed{}, then fix the one convention we add: the index of an
\term{w:leader}{empirical leader} when an arm has observed only the reward $1$.

\subsection{IMED in the classical setting}

For a model with a known support, \citet{HT10} and \citet[Algorithm~1]{HT15} compare each arm with
the largest empirical mean through the divergence $\Kinf$. After $n$ pulls of arm $a$, with empirical
law $\hat\nu$, and with the largest empirical mean $\hat m^\star$, the index of \citet{HT15} is
\begin{equation*}
  n\,\Kinf(\hat\nu,\hat m^\star)+\log n,
\end{equation*}
and \imed{} pulls an arm whose index is minimal. Unlike an upper confidence bound, this index
needs no inversion of $\Kinf$, and for $\hat m^\star<1$ it is a one-dimensional concave maximisation by
\eqref{eq:dual}.

\subsection{IMED on semi-bounded rewards}

We use the same index, with one convention: the divergence term of an \term{w:leader}{empirical leader} is $0$.

\begin{definition}[IMED]\label{def:imed}
\uses{def:model, def:kinf, def:run}
\imed{} is the \term{w:learner}{learner} of Algorithm~\ref{alg:imed}. At rounds $t<\Karms$ it pulls each arm once. At a
round $t\ge \Karms$, every arm has $\Np_a(t)\ge1$; \dt{s:hmt}let $\hmt_a(t)\coloneqq\hm_{a,\Np_a(t)}$,
\dt{s:hnut}$\hnut_a(t)\coloneqq\hnu_{a,\Np_a(t)}$ \dt{s:hms}and $\hms(t)\coloneqq\max_{a\in\arms}\hmt_a(t)$.
\dt{w:leader}An arm with $\hmt_a(t)=\hms(t)$ is an \emph{empirical leader}. \dt{s:D}Let
\[
  \Dv_a(t)\coloneqq\begin{cases}0&\text{if $a$ is an \term{w:leader}{empirical leader},}\\
  \Kinf(\hnut_a(t),\hms(t))&\text{otherwise.}\end{cases}
\]
\dt{s:I}The index of arm $a$ is
\begin{equation}\label{eq:index}
  \Ix_a(t)\coloneqq \Np_a(t)\Dv_a(t)+\log \Np_a(t)\in[0,\infty],
\end{equation}
\dt{s:Is}and the minimal index is $\Is(t)\coloneqq\min_a\Ix_a(t)$. \imed{} pulls an arm $\Ap_t$ with
$\Ix_{\Ap_t}(t)=\Is(t)$. \dt{w:tie}A \emph{tie rule} chooses among the minimisers, and the order of the
first $\Karms$ rounds, as a measurable function of $\Useed$ and of the past; every result below holds for every
\term{w:tie}{tie rule}.
\end{definition}

Intuitively, the term $\Np_a(t)\Dv_a(t)$ is large when the samples of arm $a$ are unlikely under every
law whose mean exceeds the best empirical mean, and the term $\log \Np_a(t)$ grows with the count and
forces the exploration of the other arms.

When $\hms(t)<1$, the convention changes nothing: \imed{} is then the algorithm of
\citet[Algorithm~1]{HT15}. The case $\hms(t)=1$, and why the convention matters there, is
discussed in Section~\ref{sec:upper-ht15}.

\begin{algorithm}[t]
\caption{\texttt{IMED} (Indexed Minimum Empirical Divergence)}\label{alg:imed}
\begin{algorithmic}[1]
\Require the arm set $\arms$; no horizon, no bound on the rewards other than $1$, no moment
\For{$t=0,1,\dots,\Karms-1$}
  \State pull each arm once, in an order given by the \term{w:tie}{tie rule}
\EndFor
\For{$t=\Karms,\Karms+1,\dots$}
  \State $\hmt_a(t)\gets\hm_{a,\Np_a(t)}$ and $\hnut_a(t)\gets\hnu_{a,\Np_a(t)}$ for every arm $a$;
    $\hms(t)\gets\max_a\hmt_a(t)$
  \State $\Dv_a(t)\gets0$ if $\hmt_a(t)=\hms(t)$, else $\Dv_a(t)\gets\Kinf(\hnut_a(t),\hms(t))$
  \State $\Ix_a(t)\gets \Np_a(t)\Dv_a(t)+\log \Np_a(t)$
  \State pull $\Ap_t\in\argmin_a\Ix_a(t)$, chosen by the \term{w:tie}{tie rule}
\EndFor
\end{algorithmic}
\end{algorithm}


% ----------------------------------------------------------------
\section{Regret upper bound under a finite variance}\label{sec:upper}

In this section we state our finite-time bound, then recall the bound of \citet{HT15} under an
exponential moment and compare the two. Both have the same first-order term; they differ in the term
that counts the rounds at which the \term{w:optimal}{optimal} arm looks bad.

\subsection{The bound under a finite variance}

Our bound uses the rate of a Chernoff bound for the \term{w:subopt}{suboptimal} arm, at a maximiser of the dual
formula~\eqref{eq:dual}, and the expected supremum of an index-like functional of the
\term{w:optimal}{optimal} arm's samples.

\begin{definition}[Rate of a suboptimal arm]\label{def:rate}
\uses{def:model, def:kinf, lem:optimizer}
Let $P\in\Sm$ and $\mean(P)<\mu<1$. The map $\lambda\mapsto\E_P\log(1-\lambda(X-\mu))$ attains its
maximum on $[0,1/(1-\mu)]$ (Lemma~\ref{lem:optimizer}); \dt{s:lst}let $\lst(P,\mu)$ be a maximiser,
fixed once and for all. For $u\in\R$, let
\begin{equation*}
  \dt{s:Lt}\Lt(P,\mu,u)\coloneqq\sup_{\eta\in[-1,0]}\Big(\eta u-\log\E_P\big[\big(1-\lst(P,\mu)(X-\mu)\big)^{\eta}\big]\Big).
\end{equation*}
\dt{s:Ki}For an \term{w:instance}{instance} with $\ms<1$, a \term{w:subopt}{suboptimal} arm $i$ and $\delta>0$, let
$\Ki_i\coloneqq\Kinf(\nu_i,\ms)$ and $\Lt_i(\delta)\coloneqq\Lt(\nu_i,\ms,\Ki_i-\delta/(1-\ms))$.
\end{definition}

\begin{definition}[Deviation functional of the optimal arm]\label{def:dev}
\uses{def:model, def:kinf}
For $\mu\in\R$, $s\ge1$ and a sequence $x=(x_1,x_2,\dots)$ of reals in $(-\infty,1]$, let
\begin{equation*}
  \dt{s:w}\ws_s(x;\mu)\coloneqq s\exp\big(s\Kinf(\hP[x_{1:s}],\mu)\big)\,\ind\{\bx_s\le\mu\}\in[0,\infty],
\end{equation*}
with $\exp(+\infty)\coloneqq+\infty$ and $0\cdot\infty\coloneqq0$. For $P\in\Sm$, let
\begin{equation*}
  \dt{s:Gamma}\Gam(P,\mu)\coloneqq\E\Big[\sup_{s\ge1}\ws_s(X;\mu)\Big]\in[0,\infty],
\end{equation*}
where $X=(X_1,X_2,\dots)$ is an i.i.d.\ sequence of law $P$.
\end{definition}

The map $x_{1:s}\mapsto\Kinf(\hP[x_{1:s}],\mu)$ is Borel (Lemma~\ref{lem:dualfacts}(c), \emph{empirical
dual}), so $\Gam$ is well defined. Note that $\log \ws_s(x;\mu)$ is the index \eqref{eq:index} of an arm
with the samples $x_{1:s}$, computed at the reference level $\mu$ instead of $\hms(t)$.

We obtain the following guarantee for \imed{}.

\begin{theorem}[IMED under a finite variance of one optimal arm]\label{thm:main}
\uses{def:model, def:kinf, def:run, def:imed, def:dev, def:rate}
Let $\nu\in\Sm^\arms$ be an \term{w:instance}{instance} with $\ms<1$, and let $j$ be an \term{w:optimal}{optimal} arm with
$\nu_j\in\Sm_2$. Let $i$ be a \term{w:subopt}{suboptimal} arm. Then $0<\Ki_i<\infty$. Let
$\delta\in(0,(1-\ms)\Ki_i/2)$. Then $\Lt_i(\delta)>0$ and $\Gam(\nu_j,\ms-\delta)<\infty$, and
\imed{}, with any \term{w:tie}{tie rule}, satisfies for every $T\ge1$
\[
  \E \Np_i(T)\le1+\frac{\log T}{\Ki_i-2\delta/(1-\ms)}+\frac{1}{1-e^{-\Lt_i(\delta)}}
  +\Gam(\nu_j,\ms-\delta).
\]
\end{theorem}

(See the proof in Appendix~\ref{sec:proof-main}.)

\paragraph{Nature of the constants} The first three terms are explicit. The rate $\Lt_i(\delta)$
depends only on $(\nu_i,\ms,\delta)$ and uses no moment of $\nu_i$ beyond its mean
(Lemma~\ref{lem:subopt}). The last term depends only on $(\nu_j,\delta)$. Inequality~\eqref{eq:gammabound}
of Appendix~\ref{sec:proof-key} bounds it by an explicit expression in $\delta$, $1-\ms+\delta$, the
first and second moments of $(X-\ms+\delta)^-$ and the second moment of $X-\ms+\delta$ under
$\nu_j$, and an exponent $\gamma\in(1,2]$ with
\begin{equation*}
  \E_{\nu_j}\Big[\Big(\frac{1-X}{1-\ms+\delta}\Big)^\gamma\Big]\le1.
\end{equation*}
Such a $\gamma$ exists (Lemma~\ref{lem:moments}), but it is not explicit in general. We do not track
the dependence on $\delta$ as $\delta\to0$.

\paragraph{The range of $\delta$} The admissible $\delta$ satisfy $\delta<\gap_i/2$, because
\begin{equation*}
  (1-\ms)\Ki_i\le\gap_i
\end{equation*}
(Lemmas~\ref{lem:dualfacts}(b) and~\ref{lem:shift} at the levels $\marm_i$ and $\ms$). Unlike in
Theorem~\ref{cite:HT15-thm3}, where $\delta<\min_{i'}\gap_{i'}/2$, the range involves arm $i$ only: it
is the range on which the second term is finite and positive.

\paragraph{Scope} Theorem~\ref{thm:main} is a first-order result: optimising $\delta$ gives a
remainder $o(\log T)$ without a rate, whereas \citet[Corollary~4]{HT15} gives $O((\log T)^{10/11})$
under an exponential moment. A finite variance is sufficient, not necessary: Theorem~\ref{thm:counter}
shows that it cannot be replaced by a lower moment for power-law tails, but laws with an infinite
variance and a tail lighter than $y^{-2}$ (for example $\Prob(X\le-y)$ of order $1/(y^2\log y)$) are
covered by neither theorem. The assumption $\ms<1$ cannot be dropped: if an \term{w:optimal}{optimal} arm has the law
$\delta_1$, every non-leader has an infinite index. Only \imed{} is analysed.

\subsection{Comparison with the bound under an exponential moment}\label{sec:upper-ht15}

We now recall the finite-time bound of \citet{HT15}, of which Theorem~\ref{thm:main} keeps the first
two terms. \citet{HT15} write $\mathcal A$ for the set of all laws on $(-\infty,1]$, $F$ for a law, $\mathrm
E(F)$ for its mean, $D(F\|G)$ for $\KL(F,G)$, $D_{\inf}(F,\mu;\mathcal A)$ for
$\inf\{D(F\|G):G\in\mathcal A,\ \mathrm E(G)>\mu\}$, $\nu$ for the dual variable $\lambda$,
$L_{\max}(F,\mu)$ for the right-hand side of~\eqref{eq:dual}, $T_i(n)$ for $\Np_i(T)$ (pulls in the
first $n=T$ rounds), $\mu_i$, $\mu^*$ for $\marm_i$, $\ms$, $\mathcal I_{\mathrm{opt}}$ for the set of optimal
arms, $\nu_i^*$ for $\lst(\nu_i,\ms)$, and $\lambda_{i,\mu}$ for $\gamma_{i,\mu}$ below (with $P=\nu_i$, the
supremum of the exponents $\gamma$ with $\phi(\gamma)\le1$ in Lemma~\ref{lem:moments}, \emph{moments
of the dual statistic}: the letter $\gamma$ always denotes such an exponent). A law
of $(-\infty,1]$ with mean above $\mu$ has a finite mean, so $D_{\inf}(P,\mu;\mathcal A)=\Kinf(P,\mu)$
for $P\in\Sm$.


% Verbatim: literature/notes/HT15.md, (4)-(6) and Theorem 3.
\begin{theorem}[{From \citet[Theorem~3]{HT15}}]\label{cite:HT15-thm3}
\uses{def:model, def:kinf, def:run, def:imed, def:rate}
For an arm $a$ and $\mu<1$, let $\gamma_{a,\mu}\coloneqq\sup\{\gamma\in\R\cup\{\infty\}:
\E_{\nu_a}[((1-X)/(1-\mu))^\gamma]\le1\}$ and $\Lambda_a^*(x)\coloneqq\sup_{\eta\in\R}\{\eta x-\log
\E_{\nu_a}[e^{\eta X}]\}$; for a suboptimal arm $i$, let
$\Lt_i^*(x)\coloneqq\sup_{\eta\in\R}\{\eta x-\log\E_{\nu_i}[(1-\lst(\nu_i,\ms)(X-\ms))^\eta]\}$.
Let $\nu\in\Sm^\arms$ with $\ms<1$, and assume that $\E_{\nu_j}[e^{\eta X}]<\infty$ for every $\eta$ in
a neighbourhood of $0$, for some \term{w:optimal}{optimal} arm $j$. Then, for every fixed
$0<\delta<\min_{i':\gap_{i'}>0}\gap_{i'}/2$, every \term{w:subopt}{suboptimal} arm $i$ and every $T\ge1$, \imed{}
satisfies
\begin{align*}
  \E \Np_i(T)\le{}&\frac{\log T}{\Ki_i-\frac{2\delta}{1-\ms}}+\frac{1}{1-e^{-\Lt_i^*(\Ki_i-\frac{\delta}{1-\ms})}}\\
  &+\min_{j'\text{ optimal}}\frac{6e}{(1-1/\gamma_{j',\ms-\delta})\big(1-e^{-(1-1/\gamma_{j',\ms-\delta})\Lambda_{j'}^*(\ms-\delta)}\big)^3}.
\end{align*}
\end{theorem}

\citet{HT15} state it for their Algorithm~1 with an arbitrary \term{w:tie}{tie rule}. \imed{} of
Definition~\ref{def:imed} is that algorithm with one choice when every index of
\citet[Algorithm~1]{HT15} is infinite (the remark below), so this restatement is a special
case of theirs.

\begin{remark}[Relation to the algorithm of \citet{HT15}]
If $\hms(t)<1$, then $\Kinf(\hnut_a(t),\hms(t))=0$ for an \term{w:leader}{empirical leader} $a$
(Lemma~\ref{lem:dualfacts}(b), \emph{zero at or above the mean}), so $\Ix_a(t)$ is the index of
\citet[Algorithm~1]{HT15}. If $\hms(t)=1$ (an arm has observed only the reward $1$), their
formula gives $+\infty$ to every arm, since no law of $\Sm$ has a mean above $1$, and
Definition~\ref{def:imed} pulls an \term{w:leader}{empirical leader}: one admissible choice of their arbitrary tie
rule. The choice matters. The proof of \citet[Theorem~3]{HT15} uses that a leader's index is at most
the logarithm of the round, and a rule that pulls a \term{w:subopt}{suboptimal} arm whenever every index is infinite
pulls it at every round with positive probability if an \term{w:optimal}{optimal} arm has an atom at $1$ (its first
reward may be $1$, after which its empirical law never changes).
\end{remark}

\paragraph{The difference is the last term} The logarithmic term and the term
in $\Lt_i(\delta)$ are those of Theorem~\ref{cite:HT15-thm3}, with the supremum in the rate restricted
to $\eta\in[-1,0]$; the additional $1$ counts the first pull. The difference is the last term. In
Theorem~\ref{cite:HT15-thm3} it is a function of Cram\'er's rate $\Lambda_j^*(\ms-\delta)$ of the lower
tail of the empirical mean of arm $j$, which is positive only if
\begin{equation*}
  \E_{\nu_j}e^{-\eta X}<\infty\quad\text{for some }\eta>0.
\end{equation*}
Here it is $\Gam(\nu_j,\ms-\delta)$, finite as soon as $\nu_j$ has a finite variance
(Proposition~\ref{prop:key}). By Lemma~\ref{lem:sup}, $\Gam(\nu_j,\ms-\delta)$ bounds the expected
number of rounds at which the \term{w:subopt}{suboptimal} arm $i$ is pulled while every empirical mean is at most
$\ms-\delta$: it is the price of an \term{w:optimal}{optimal} arm that looks bad. As $T\to\infty$ with $\delta$ fixed, the
term $\log T/(\Ki_i-2\delta/(1-\ms))$ dominates.

% ----------------------------------------------------------------
\section{Asymptotic optimality}\label{sec:optimal}

In this section we first recall the lower bound for \term{w:ufc}{uniformly fast convergent} \term{w:learner}{learners}, then show
that \imed{} attains it on the semi-bounded laws with a finite moment of order $r\ge2$.

\subsection{Lower bound in the classical setting}

\citet{GMS19} write $\underline\nu$ for an \term{w:instance}{instance}, $\psi$ for a \term{w:learner}{learner}, $N_{\psi,a}(T)$ for
$\Np_a(T)$ (pulls at the rounds $1,\dots,T$, the same count as ours), $\mu^\star$ for $\ms$, $\ln$ for
$\log$ and $\Kinf(\nu_a,x,\mathcal D)$ for $\Kinf(\nu_a,x;\mathcal D)$, defined with the same strict
inequality; uniform fast convergence is \citet[Definition~1]{GMS19}, as in
Definition~\ref{def:classes}.

% Verbatim: literature/notes/GMS19.md, Theorem 1.
\begin{theorem}[{From \citet[Theorem~1]{GMS19}}]\label{cite:GMS19-thm1}
\uses{def:model, def:kinf, def:run, def:classes}
Let $\mathcal D$ be a set of laws on $\R$ with finite means, and let a \term{w:learner}{learner} be uniformly fast
convergent on $\mathcal D$. Then for every $\nu\in\mathcal D^\arms$ and every \term{w:subopt}{suboptimal} arm $a$,
$\liminf_{T\to\infty}\E \Np_a(T)/\log T\ge1/\Kinf(\nu_a,\ms;\mathcal D)$.
\end{theorem}

\subsection{Asymptotic optimality of IMED}

\begin{corollary}[Asymptotic optimality]\label{cor:optimal}
\uses{def:model, def:kinf, def:run, def:classes, def:imed, thm:main}
Consider \imed{} with any \term{w:tie}{tie rule}.
\begin{enumerate}[label=(\alph*)]
\item Under the assumptions of Theorem~\ref{thm:main}, \imed{} satisfies
  \[
    \limsup_{T\to\infty}\frac{\E \Np_i(T)}{\log T}\le\frac{1}{\Kinf(\nu_i,\ms)}.
  \]
\item Let $r\ge2$. Every \term{w:instance}{instance} of $\Mm_r^\arms$ has $\ms<1$, and \imed{} is \term{w:consistent}{consistent} and
  \term{w:ufc}{uniformly fast convergent} on $\Mm_r$. For every $\nu\in\Mm_r^\arms$ and every \term{w:subopt}{suboptimal} arm
  $i$,
  \[
    \lim_{T\to\infty}\frac{\E \Np_i(T)}{\log T}=\frac{1}{\Kinf(\nu_i,\ms)}=\frac{1}{\Kinf(\nu_i,\ms;\Mm_r)},
  \]
  and every \term{w:learner}{learner} that is \term{w:ufc}{uniformly fast convergent} on $\Mm_r$ (in particular every \term{w:consistent}{consistent}
  one) satisfies $\liminf_{T\to\infty}\E \Np_i(T)/\log T\ge1/\Kinf(\nu_i,\ms;\Mm_r)$. Hence
  $\lim_T\Reg_T/\log T=\sum_{i:\gap_i>0}\gap_i/\Kinf(\nu_i,\ms)$ for \imed{}, the smallest possible
  $\liminf_T\Reg_T/\log T$ for a \term{w:learner}{learner} \term{w:ufc}{uniformly fast convergent} on $\Mm_r$.
\end{enumerate}
\end{corollary}

(See the proof in Appendix~\ref{sec:proof-optimal}.)

\paragraph{Comparison with the lower bound} Item (b) pairs the upper bound of Theorem~\ref{thm:main}
with the lower bound of Theorem~\ref{cite:GMS19-thm1} in the same class, the \term{w:learner}{learners} uniformly fast
convergent on $\Mm_r$: the first-order constant $\sum_{i:\gap_i>0}\gap_i/\Kinf(\nu_i,\ms)$ cannot be
improved, and \imed{} attains it on every \term{w:instance}{instance} of $\Mm_r^\arms$. The case $r=2+p$ is the class of
the conjecture of \citet[Section~9]{HT15}. The law $\delta_1$ is excluded because an \term{w:instance}{instance} with an
arm $\delta_1$ has $\ms=1$. The equality of the two divergences in (b) is Lemma~\ref{lem:class}: the
lower-bound constant does not depend on the moment class.

% ----------------------------------------------------------------
\section{Lower bounds for heavier lower tails}\label{sec:counter}

In this section we show that a finite variance cannot be replaced by a lower moment on the scale of
power-law lower tails: with a tail $\kappa y^{-q}$, $1<q<2$, \imed{} is not \term{w:consistent}{consistent}, and with
a tail of order $y^{-2}$ it misses the lower-bound constant.

\begin{theorem}[Heavier lower tails: IMED fails]\label{thm:counter}
\uses{def:model, def:kinf, def:run, def:imed}
Let $\arms=\{1,2\}$, $\nu_2=\delta_0$, and $\nu_1=P\in\Sm$ with $\mean(P)\in(0,1)$, so that arm $1$ is
optimal and arm $2$ \term{w:subopt}{suboptimal}. For $X\sim P$, let $Y\coloneqq-X$. Consider \imed{} with any \term{w:tie}{tie rule}.
\begin{enumerate}[label=(\roman*)]
\item If there are $q\in(1,2)$, $\kappa>0$ and $y_0>0$ with $\Prob(Y\ge y)\ge\kappa y^{-q}$ for
  every $y\ge y_0$, then there is $T_0$, depending on $P$, such that for every $T\ge T_0$,
  \[
    \E \Np_2(T)\ge\frac{\kappa\,2^q}{64\,e^q}\,T^{2-q}.
  \]
\item If there are $0<\kappa\le\kappa'$ and $y_0>0$ with
  $\kappa y^{-2}\le\Prob(Y\ge y)\le\kappa'y^{-2}$ for every $y\ge y_0$, then
  \[
    \liminf_{T\to\infty}\frac{\E \Np_2(T)}{\log T}\ge\frac{1}{\log(1/(1-\mean(P)))}+\frac{\kappa}{4e},
    \qquad\text{where }\ \Kinf(\delta_0,\mean(P))=\log\frac{1}{1-\mean(P)}.
  \]
\end{enumerate}
\end{theorem}

(See the proof in Appendix~\ref{sec:proofs-counter}.) For $Y$ with a finite mean, the hypothesis of
(i) implies that $\E|X|^r=\infty$ for $r\ge q$; it allows $\E|X|^r<\infty$ for every $r<q$.

\begin{corollary}[Pareto lower tails]\label{cor:pareto}
\uses{def:model, def:kinf, def:run, def:classes, def:imed, thm:counter, cite:GMS19-thm1}
For $q>1$ and $0<\varpi<(q-1)/(2q-1)$, let $V$ satisfy $\Prob(V\ge v)=v^{-q}$ for $v\ge1$, and let
\dt{s:Pq}$\Pq_{q,\varpi}$ be the law of $X$ equal to $1$ with probability $1-\varpi$ and to $-V$ with probability
$\varpi$. Then $\Pq_{q,\varpi}\in\Sm$, $\mean(\Pq_{q,\varpi})=1-\varpi(2q-1)/(q-1)\in(0,1)$, and $\E|X|^r<\infty$ if and
only if $r<q$. Consider \imed{} on the \term{w:instance}{instance} $(\nu_1,\nu_2)=(\Pq_{q,\varpi},\delta_0)$, and let
$m\coloneqq \mean(\Pq_{q,\varpi})$.
\begin{enumerate}[label=(\alph*)]
\item If $1<q<2$, then $\Reg_T\ge m\varpi2^qe^{-q}T^{2-q}/64$ for every large $T$. Hence, for every
  $p\in(0,1)$, \imed{} is not \term{w:consistent}{consistent} on $\Sm_{1+p}$ with two arms: take $q\in(1+p,2)$.
\item If $q=2$, then $\liminf_T\Reg_T/\log T\ge m\big(1/\log\frac1{1-m}+\varpi/(4e)\big)$. Let
  $\mathcal C\subseteq\Sm$ contain $\Pq_{2,\varpi}$ and every law supported by $\{0,1\}$ (for example
  $\mathcal C=\bigcap_{r<2}\Sm_r$). Then $\Kinf(\delta_0,m;\mathcal C)=\log\frac{1}{1-m}$, and \imed{}
  satisfies $\liminf_T\E \Np_2(T)/\log T\ge1/\Kinf(\delta_0,m;\mathcal C)+\varpi/(4e)$. So, at this
  \term{w:instance}{instance} of $\mathcal C^{\{1,2\}}$, the $\liminf$ of \imed{} exceeds by at least $\varpi/(4e)$ the
  constant that Theorem~\ref{cite:GMS19-thm1} gives for every \term{w:learner}{learner} \term{w:ufc}{uniformly fast convergent} on
  $\mathcal C$.
\end{enumerate}
\end{corollary}

(See the proof in Appendix~\ref{sec:proofs-counter}.)

\paragraph{The threshold, and what is not claimed} Let the optimal arm have a lower tail
$\Prob(X\le-y)$ of order $y^{-q}$, with $q>1$ for a finite mean. If $q>2$, the variance is finite and
Theorem~\ref{thm:main} applies. If $q=2$, \imed{} does not attain the lower-bound constant
(Theorem~\ref{thm:counter}(ii)); we do not know whether its regret is still of order $\log T$. If
$q<2$, \imed{} is not \term{w:consistent}{consistent}. So, on the scale of power-law tails, the exponent two is the
threshold. This does not make a finite variance necessary in general: the tails between $y^{-2}$ and
the square-integrable ones are open (Open question~\ref{oq:boundary}). Nor does it say that fewer
moments make the problem harder: the lower-bound constant $\log(1/(1-m))$ of
Corollary~\ref{cor:pareto}(b) is the same in every class containing the laws on $\{0,1\}$. The
failure is a property of \imed{}, as the next paragraph indicates.

\paragraph{Another algorithm} \citet[Theorem~3.1]{PA26} prove that KLinf-UCB satisfies
$\limsup_T\E\Np_i(T)/\log T\le1/\mathrm{KL}_{\inf}(\nu_i,\ms)$ on every class of laws that satisfies
their Assumptions~2.3 and~2.4, where $\mathrm{KL}_{\inf}$ is taken over the class, with the
inequality $\mean(Q)\ge\mu$. Lemma~\ref{lem:pa-assumptions} (\emph{concentration with a finite mean
only}) shows that the class $\Mm_1=\Sm\setminus\{\delta_1\}$ satisfies both with a finite mean only:
over $\Mm_1$, their $\mathrm{KL}_{\inf}$ is $\Kinf$ (part (c)); Assumption~2.3 holds with
$g(n)=\log(1+n)+C_2$ for every $C_2>0$ (part (a)); Assumption~2.4 holds at every level in
$(\mean(P),1)$ (part (b)). Assumption~2.4 is stated there at every level above the mean, and fails on
$\Mm_1$ at the levels $\mu\ge1$, where both divergences are infinite (as it does for laws on $[a,b]$
at the levels above $b$, one of their examples); the proof of \citet[Theorem~3.1]{PA26} uses it at
the best mean only, which is below $1$ on $\Mm_1^\arms$. So their theorem indicates that KLinf-UCB,
computed with $\Kinf$, attains the constant $1/\Kinf(\nu_i,\ms)$ on every \term{w:instance}{instance} of
$\Mm_1^\arms$, with no moment condition, at the price of inverting $\Kinf$ at each round. We have not
checked their proof line by line, and we state this as a reduction, not as a theorem;
Open question~\ref{oq:finite-mean} asks for an index policy of the type of \imed{} with the same
property.

\paragraph{The mechanism} \imed{} compares each arm with the largest \emph{empirical mean}.
One negative reward $-y$ of the \term{w:optimal}{optimal} arm, with $y$ larger than the sum of its other rewards, puts
its empirical mean below $0$: the deterministic arm becomes the \term{w:leader}{empirical leader}, with index
$\log \Np_2$, while the index of the \term{w:optimal}{optimal} arm stays frozen at a value of at least $\log(y/e)$
(Lemma~\ref{lem:jump}). The \term{w:subopt}{suboptimal} arm is then pulled until $\Np_2$ reaches $y/e$. Under a tail
$y^{-q}$, a reward of order $T$ occurs among the first $T/4$ pulls of the \term{w:optimal}{optimal} arm with probability
of order $T^{1-q}$, and costs of order $T$ pulls: this gives (i). For $q=2$, the jumps of size
$y\approx s$ at the $s$-th pull add up to an extra term of order $\log T$: this gives (ii).


% ----------------------------------------------------------------
\section{How the proofs work}\label{sec:overview}

In this section we give the architecture of the proofs: the decomposition of \citet{HT15}
(Section~\ref{sec:overview-decomp}), where their argument needs an exponential moment
(Section~\ref{sec:overview-mgf}), the two new steps (Section~\ref{sec:overview-new}), how the pieces give
Theorem~\ref{thm:main} (Section~\ref{sec:overview-combine}) and the mechanism of
Theorem~\ref{thm:counter} (Section~\ref{sec:overview-counter}).

\subsection{The decomposition}\label{sec:overview-decomp}

Fix a \term{w:subopt}{suboptimal} arm $i$, an \term{w:optimal}{optimal} arm $j$ with a finite variance, $\delta$ as in
Theorem~\ref{thm:main}, and the level $\mu\coloneqq\ms-\delta$. After the first $\Karms$ rounds, a pull of
$i$ happens either while every empirical mean is at most $\mu$, or while some empirical mean exceeds
$\mu$:
\begin{equation}\label{eq:decomp}
  \Np_i(T)\le1+\underbrace{\sum_{t=\Karms}^{T-1}\ind\{\Ap_t=i,\ \hms(t)\le\mu\}}_{\TA}
  +\underbrace{\sum_{t=\Karms}^{T-1}\ind\{\Ap_t=i,\ \hms(t)>\mu\}}_{\TB}
\end{equation}
(Lemma~\ref{lem:decomp}). In $\TB$, the reference level of \imed{} is above $\mu$, so a pull of $i$
requires its empirical law to be close to a law of mean at least $\mu$; this costs $\log T/\Ki_i$
pulls at first order and uses no moment (Lemmas~\ref{lem:termB} and~\ref{lem:termC}). The term $\TA$
counts the rounds at which the \term{w:optimal}{optimal} arms look bad. It is the only place where the tails of the
\term{w:optimal}{optimal} arm matter (Lemma~\ref{lem:termA}).

\subsection{Where the exponential moment enters the analysis of \texorpdfstring{\citet{HT15}}{Honda and Takemura (2015)}}\label{sec:overview-mgf}

\citet[(25)--(26)]{HT15} bound $\TA$ by splitting it according to the value $s$ of $\Np_j(t)$ and
applying their counting lemma (Lemma~\ref{lem:sup}(a) below) once for each $s$. This gives the sum
\begin{equation*}
  \sum_{s\ge1}\ws_s(\Xs_{j,\cdot};\mu),\qquad\text{where }\Xs_{j,\cdot}\coloneqq(\Xs_{j,1},\Xs_{j,2},\dots).
\end{equation*}
They bound the expectation of each term by \citet[Proposition~11]{HT15}, in which the probability
that $\hm_{j,s}\le\mu$ is bounded by Cram\'er's bound $e^{-s\Lambda_j^*(\mu)}$; its rate is positive
only under an exponential moment of the lower tail. Replacing Cram\'er's bound by a polynomial one does
not help, because the sum itself has an infinite expectation for many laws with a finite variance.

\begin{lemma}[A summed bound needs three moments]\label{lem:three-moments}
\uses{def:model, def:kinf, def:dev}
Let $P\in\Sm$, $\mu<\mean(P)$, and let $X=(X_1,X_2,\dots)$ be i.i.d.\ of law $P$ with
$\E[(X_1^-)^3]=\infty$. Then
\[
  \sum_{s\ge1}\E\,\ws_s(X;\mu)\ge\sum_{s\ge1}s\,\Prob(\bX_s\le\mu)=\infty.
\]
\end{lemma}

(See the proof in Appendix~\ref{sec:proof-three}.) This is the necessity half of a one-sided form of
the theorem of Baum and Katz on convergence rates in the law of large numbers. For one-sided theorems
of this kind, see also \citet{CL75}, the work that \citet[Section~9]{HT15} cite for the moments of
order $2+p$ of their conjecture, and \citet{BK65} for the two-sided theorem. With
$P=\Pq_{q,\varpi}$ of Corollary~\ref{cor:pareto} and $2<q\le3$, the variance is finite and the summed
bound is infinite.


\subsection{The two new steps}\label{sec:overview-new}

The first step counts the same rounds by a supremum over $s$ instead of a sum.

\begin{lemma}[Counting the rounds at which every empirical mean is low]\label{lem:sup}
\uses{def:model, def:kinf, def:run, def:imed, def:dev}
Let $\nu\in\Sm^\arms$ and consider \imed{} with any \term{w:tie}{tie rule}. Then, for every outcome:
\begin{enumerate}[label=(\alph*)]
\item for every arm $i$ and every $x\in\R$,
  $|\{t\ge \Karms:\Ap_t=i,\ \Is(t)\le x\}|\le e^x$;
\item for all arms $i,j$ and every $\mu\in\R$,
  $|\{t\ge \Karms:\Ap_t=i,\ \hms(t)\le\mu\}|\le\sup_{s\ge1}\ws_s(\Xs_{j,\cdot};\mu)$.
\end{enumerate}
\end{lemma}

(See the proof in Subsection~\ref{sec:proof-lem-sup}.) Part~(a) is \citet[Lemma~13]{HT15}.
\citet{HT15} apply it once for each value $s$ of $\Np_j(t)$, which counts every round once for each $s$
and gives a sum. Part~(b) applies it once, at the random level
\begin{equation*}
  M\coloneqq\sup\big\{s\Kinf(\hnu_{j,s},\mu)+\log s:\ s\ge1,\ \hm_{j,s}\le\mu\big\},
\end{equation*}
which gives the supremum. The second step bounds the expectation of that supremum.

\begin{proposition}[Key deviation bound]\label{prop:key}
\uses{def:model, def:kinf, def:dev}
Let $P\in\Sm_2$ and $\mu<\mean(P)$. Then $\Gam(P,\mu)<\infty$. An explicit upper bound is
\eqref{eq:gammabound} in Appendix~\ref{sec:proof-key}.
\end{proposition}

(See the proof in Appendix~\ref{sec:proof-key}.)

\paragraph{Why the variance enters, and only once} Let $X_1,X_2,\dots$ be i.i.d.\ of law $P$,
$Z_l\coloneqq X_l-\mu$ and $Y_l\coloneqq Z_l^-$. By the dual formula~\eqref{eq:dual},
\begin{equation*}
  \exp\big(s\Kinf(\hP[X_{1:s}],\mu)\big)=\max_{0\le\lambda\le1/(1-\mu)}\prod_{l=1}^s(1-\lambda Z_l).
\end{equation*}
A negative sample $Z_l=-y$ contributes the factor $1+\lambda y$, which is \emph{linear} in $y$:
the Cram\'er bound on the empirical mean would make it exponential. Split the indices $l\le s$ into
the \emph{big jumps}, with $Y_l>\beta s$ for a small constant $\beta$, and the \emph{good part}.
If the good part has a sum of at least $s\delta/2$, where $\delta\coloneqq \mean(P)-\mu$, and there is one
big jump $l^\star$, then the maximum above is at most $2Y_{l^\star}/(s\delta)$
(Lemma~\ref{lem:goodpart}), so $\ws_s\le2Y_{l^\star}/\delta$. A given big jump at index $l$ can
matter only for $s$ with $\beta s<Y_l$ and $s\ge l$, so it contributes to $\sup_s\ws_s$ at most
$\frac2\delta Y_l\ind\{Y_l>\beta l\}$, and
\begin{equation*}
  \E\sum_{l\ge1}\frac2\delta Y_l\ind\{Y_l>\beta l\}\le\frac{2\E Y_1^2}{\beta\delta}.
\end{equation*}
This is the only term in which the second moment is indispensable. The other cases are summed
over $s$: several big jumps (independence gives a factor $(\E[Y_1\ind\{Y_1>\beta s\}])^2$, summable in
$s$ under a finite variance), many big jumps (Chebyshev's inequality and a H\"older bound with an
exponent $\gamma>1$, Lemma~\ref{lem:moments}), and a good part that is too small (Bennett's
inequality, Lemma~\ref{lem:bennett}). Splitting a sum of heavy-tailed variables into its large
summands and a truncated rest is classical in the large deviations of such sums (see also
\citealp{FN71,Nag79}); what is specific here is the dual product, in which one big jump costs a
factor linear in its size.

\paragraph{Relation with the last exit time} Since $\Kinf\ge0$, the functional of
Definition~\ref{def:dev} satisfies $\ws_s(X;\mu)\ge s\,\ind\{\bX_s\le\mu\}$, so
\begin{equation*}
  \sup_{s\ge1}\ws_s(X;\mu)\ge L\coloneqq\sup\{s\ge1:\bX_s\le\mu\}\qquad(\sup\emptyset\coloneqq0),
\end{equation*}
the last time at which the random walk $\sum_{l\le s}(X_l-\mu)$, of positive drift, is at most $0$.
The moments of such last exit times are classical, and a finite second moment of the negative part of
the steps is the condition for $\E L<\infty$ (see also \citealp{CL75,Jan86}); for a law on
$(-\infty,1]$ it is a finite variance. So Proposition~\ref{prop:key} can be read as an extension of
$\E L<\infty$, but not as a corollary of it: $\ws_s$ exceeds $s$ by the factor
$\exp(s\Kinf(\hP[X_{1:s}],\mu))$, which one big jump makes of the order of the jump, and the proof
must control this factor.

\subsection{How the pieces give Theorem~\ref{thm:main}}\label{sec:overview-combine}

The two terms of \eqref{eq:decomp} are controlled separately, each by one lemma of
Appendix~\ref{sec:proof-main}.
\begin{itemize}[leftmargin=*]
  \item $\TA$: by Lemma~\ref{lem:sup}(b), $\TA\le\sup_{s\ge1}\ws_s(\Xs_{j,\cdot};\mu)$ on every run, so
    $\E\TA\le\Gam(\nu_j,\mu)$, finite by Proposition~\ref{prop:key} (Lemma~\ref{lem:termA}).
  \item $\TB$: at a round counted in $\TB$, the index of an \term{w:leader}{empirical leader} is at most $\log T$,
    which forces $\Np_i(t)\Kinf(\hnut_i(t),\mu)\le\log T$. Moving the level from $\mu$ to $\ms$ costs
    $\delta/(1-\ms)$ (Lemma~\ref{lem:shift}), so $\TB$ is at most $\log T/(\Ki_i-2\delta/(1-\ms))$ plus
    a count $\TC$ of the sample sizes at which the empirical divergence of arm $i$ at the level $\ms$ is
    abnormally small (Lemma~\ref{lem:termB}); a Chernoff bound at the dual maximiser bounds $\E\TC$
    without any moment of $\nu_i$ (Lemma~\ref{lem:termC}).
\end{itemize}
Adding the bounds gives Theorem~\ref{thm:main}.

\paragraph{Where the difficulty lies} The term $\TB$ is classical. The difficulty is $\TA$ under a
finite variance only: the summed bound of \citet{HT15} is infinite without a third moment
(Lemma~\ref{lem:three-moments}), and the supremum of Lemma~\ref{lem:sup}(b) is useful only
because its expectation can be bounded directly, by the one-big-jump decomposition of
Proposition~\ref{prop:key}. That proposition is the longest proof of the paper.

\subsection{The mechanism of Theorem~\ref{thm:counter}}\label{sec:overview-counter}

The lower bounds rest on one inequality: after a negative reward $-y$ at the $s$-th pull of the
\term{w:optimal}{optimal} arm, its index at the reference level $0$ is at least $\log(y/e)$, whatever its other
rewards (Lemma~\ref{lem:jump}).

\paragraph{Frozen index} In Theorem~\ref{thm:counter}, suppose that the $s$-th pull of arm $1$
returns $-Y_s$ with $Y_s\ge s-1$. The empirical mean of arm $1$ is then at most $0$, so arm $2$
($\delta_0$) is an \term{w:leader}{empirical leader} with index $\log \Np_2$, and the reference level is $0$. As long
as arm $1$ is not pulled, its index stays at its value, at least $\log(Y_s/e)$ by
Lemma~\ref{lem:jump}. Hence arm $2$ is pulled at every round until $\Np_2\ge Y_s/e$, whatever the tie
rule. Since the $s$-th reward of arm $1$ is independent of the time of the $s$-th pull, the
expected excess of $\Np_2$ is a sum over $s$ of tail integrals of $Y$. For (i), one jump $Y_s\ge eT/2$
among the first $T/4$ pulls of arm $1$ occurs with probability of order $T\Prob(Y\ge eT/2)$ and
costs $T/2$ pulls. For (ii), the jumps $Y_s\ge s-1$ contribute
$\sum_s\int_{s-1}^{T}\Prob(Y\ge y)\,dy\approx\kappa\log T$ up to a constant factor, on top of the
$\log T/\log(1/(1-\mean(P)))$ pulls that \imed{} always spends on a point mass (Lemma~\ref{lem:explore}).
The proofs are in Appendix~\ref{sec:proofs-counter}.

\paragraph{Map of the appendices} Appendix~\ref{sec:cited} recalls the dual formula of
\citet{HT15}. Appendix~\ref{sec:dual} proves the properties of $\Kinf$ that we use (monotonicity, dual
maximiser, level shift, deviations of a \term{w:subopt}{suboptimal} arm, independence of the moment class,
and the concentration used for KLinf-UCB).
Appendix~\ref{sec:proof-main} proves Theorem~\ref{thm:main} (the decomposition, one lemma per term,
and Lemma~\ref{lem:sup}), Appendix~\ref{sec:proof-key} Proposition~\ref{prop:key},
Appendix~\ref{sec:proof-three} Lemma~\ref{lem:three-moments}, Appendix~\ref{sec:proof-optimal} Corollary~\ref{cor:optimal}, and Appendix~\ref{sec:proofs-counter}
Theorem~\ref{thm:counter} and Corollary~\ref{cor:pareto}. Appendix~\ref{sec:attribution} gives the
sources of the techniques and Appendix~\ref{sec:notation} the notation.

% ----------------------------------------------------------------
\section{Examples and simulations}\label{sec:examples}

In this section we give an \term{w:instance}{instance} satisfying all the hypotheses of Theorem~\ref{thm:main} and none
of those of Theorem~\ref{cite:HT15-thm3}, the degenerate cases, and simulations of the three regimes.

\begin{example}[A Pareto lower tail with a finite variance]
Take two arms, $\nu_1=\Pq_{3,1/10}$ of Corollary~\ref{cor:pareto} (the reward is $1$ with probability
$0.9$ and $-V$ with probability $0.1$, where $\Prob(V\ge v)=v^{-3}$ for $v\ge1$) and
$\nu_2=\delta_0$. Then $\marm_1=1-0.1\cdot5/2=0.75$, $\gap_2=0.75$, $\E X^2<\infty$ under $\nu_1$,
and $\E_{\nu_1}e^{-\eta X}=\infty$ for every $\eta>0$. So arm $1$ is the only \term{w:optimal}{optimal} arm and it has
no moment generating function near $0$: \citet[Theorem~3]{HT15} does not apply, while
Theorem~\ref{thm:main} does. By Lemma~\ref{lem:dualfacts}(d), $\Kinf(\delta_0,0.75)=\log4\approx1.386$,
and Corollary~\ref{cor:optimal} gives $\E \Np_2(T)/\log T\to1/\log4\approx0.721$. The \term{w:instance}{instance} belongs
to $\Mm_r^{\{1,2\}}$ for every $r\in[2,3)$.
\end{example}

\paragraph{Degenerate cases}
\begin{itemize}[leftmargin=*]
\item \emph{Bounded supports.} If every arm has a law on $[a,1]$, all moments are finite and
  Theorem~\ref{thm:main} applies with any \term{w:optimal}{optimal} arm; its first order is that of
  \citet[Corollary~4]{HT15}.
\item \emph{Two Bernoulli arms} with means $p_2<p_1<1$. For $p<\mu<1$, the dual~\eqref{eq:dual} of
  the Bernoulli law of mean $p$ is maximised at $\lambda=(\mu-p)/(\mu(1-\mu))$, which gives
  \dt{s:kl}$\Kinf(\cdot,\mu)=\kl(p,\mu)\coloneqq p\log\frac p\mu+(1-p)\log\frac{1-p}{1-\mu}$. So
  Corollary~\ref{cor:optimal} gives $\E \Np_2(T)/\log T\to1/\kl(p_2,p_1)$, the Lai--Robbins constant
  of the Bernoulli model.
\item \emph{Several \term{w:optimal}{optimal} arms.} One \term{w:optimal}{optimal} arm with a finite variance suffices: the others may
  have any law of $\Sm$, for instance a lower tail $y^{-3/2}$. Theorem~\ref{thm:counter} does not
  contradict this, because $\hms(t)\ge\hmt_j(t)$ for the \term{w:optimal}{optimal} arm $j$ with a finite
  variance, so an outlier of another \term{w:optimal}{optimal} arm does not lower the reference level.
\item \emph{A deterministic \term{w:optimal}{optimal} arm.} If $\nu_j=\delta_{\ms}$, then $\bx_s=\ms>\mu$ for every
  sample of arm $j$, so $\Gam(\delta_{\ms},\ms-\delta)=0$.
\item \emph{Atoms at $1$.} If an arm has observed only the reward $1$, \imed{} pulls an empirical
  leader (Definition~\ref{def:imed} and the remark of Section~\ref{sec:upper-ht15}); the \term{w:instance}{instances} of
  Corollary~\ref{cor:pareto} have such an atom.
\end{itemize}

\paragraph{Setup of the simulations} We simulate \imed{} on the instances $(\Pq_{q,1/10},\delta_0)$ of
Corollary~\ref{cor:pareto} for $q\in\{3,2,1.5\}$, with means $\marm_1=0.75$, $0.7$ and $0.6$. Arm $1$
is pulled first, ties are broken in favour of arm $1$, and each run of pulls of arm $2$ during
which the index of arm $1$ is frozen is simulated in one step (this is exact, since only $\Np_2$
changes during such a run). For $T\in\{10^3,10^4\}$ we use $3000$ independent runs, and for
$T=10^5$ another $200$ runs. On the instance $q=1.5$ we also run KLinf-UCB
\citep[Section~3]{PA26}, the algorithm of the paragraph ``Another algorithm'' of
Section~\ref{sec:counter}, with the exploration function $\log t+2\log\log t+\log(1+\Np_a(t))+1$
and $\Kinf$ computed by the dual formula~\eqref{eq:dual}, for $500$ runs up to $T=10^4$.
Table~\ref{tab:sim} reports the results. The scripts (\texttt{imed\_sim.py},
\texttt{klinf\_ucb\_sim.py}) and their output, with the command line and the seed of each
campaign, are in the directory \texttt{papers/P004-imed-finite-variance/notebook} of the
repository \url{https://github.com/prevostad/Incognita}; running
\texttt{python imed\_sim.py q runs T seed} (or \texttt{klinf\_ucb\_sim.py}) with the recorded
arguments reproduces each row.

\begin{table}[t]
\centering
\caption{\texttt{IMED} on $(P_{q,1/10},\delta_0)$: mean of $N_2(T)$ (standard error), median, $90\%$
quantile, maximum, and fraction of runs with $N_2(T)>T/10$; $3000$ runs for $T\le10^4$, $200$ runs
for $T=10^5$. Last two rows: KLinf-UCB on $(P_{1.5,1/10},\delta_0)$, $500$ runs.}\label{tab:sim}
\small
\begin{tabular}{@{}rrrrrrrr@{}}
\toprule
$q$ & $T$ & $\log T/\Kinf(\delta_0,\marm_1)$ & mean (s.e.) & median & $90\%$ & max & fraction \\
\midrule
$3$ & $10^3$ & $4.98$ & $5.15$ $(0.07)$ & $5$ & $5$ & $155$ & $0$ \\
$3$ & $10^4$ & $6.64$ & $6.46$ $(0.07)$ & $6$ & $6$ & $155$ & $0$ \\
$3$ & $10^5$ & $8.30$ & $7.17$ $(0.05)$ & $7$ & $7$ & $14$ & $0$ \\
\midrule
$2$ & $10^3$ & $5.74$ & $7.57$ $(0.27)$ & $5$ & $9$ & $508$ & $0.004$ \\
$2$ & $10^4$ & $7.65$ & $10.00$ $(0.79)$ & $7$ & $9$ & $1954$ & $0.001$ \\
$2$ & $10^5$ & $9.56$ & $9.94$ $(0.93)$ & $8$ & $9$ & $175$ & $0$ \\
\midrule
$1.5$ & $10^3$ & $7.54$ & $27.5$ $(1.5)$ & $7$ & $39$ & $901$ & $0.053$ \\
$1.5$ & $10^4$ & $10.05$ & $82.7$ $(8.9)$ & $9$ & $45$ & $7360$ & $0.016$ \\
$1.5$ & $10^5$ & $12.56$ & $135.6$ $(55.5)$ & $11$ & $78$ & $10189$ & $0.005$ \\
\midrule
\multicolumn{8}{@{}l}{KLinf-UCB} \\
$1.5$ & $10^3$ & $7.54$ & $8.42$ $(0.04)$ & $8$ & $9$ & $11$ & $0$ \\
$1.5$ & $10^4$ & $10.05$ & $14.14$ $(0.03)$ & $14$ & $15$ & $16$ & $0$ \\
\bottomrule
\end{tabular}
\end{table}

\paragraph{Interpretation of Table~\ref{tab:sim}} The three regimes of Theorems~\ref{thm:main}
and~\ref{thm:counter} are visible.
\begin{itemize}[leftmargin=*]
\item $q=3$ (finite variance): the mean of $\Np_2(T)$ is close to $\log T/\Kinf(\delta_0,\marm_1)$, and
  below it for $T\ge10^4$, as Corollary~\ref{cor:optimal} allows at first order.
\item $q=1.5$: the median follows $\log T/\Kinf$, but the mean is far above it, carried by rare runs
  in which one negative reward of arm $1$ hands the lead to arm $2$. From $T=10^3$ to $T=10^4$, the
  mean is multiplied by $3.0$ and the fraction of runs with $\Np_2(T)>T/10$ by $0.30$, against
  $\sqrt{10}\approx3.16$ and $1/\sqrt{10}\approx0.32$ for the orders $T^{2-q}$ and $T^{1-q}$ of the
  mechanism of Theorem~\ref{thm:counter}(i).
\item $q=2$: the mean exceeds $\log T/\Kinf$ by $1.8\pm0.3$ and $2.4\pm0.8$ at $T=10^3$ and $10^4$.
  Theorem~\ref{thm:counter}(ii) gives an asymptotic excess of at least $\frac{\varpi}{4e}\log T$, about
  $0.06$ and $0.09$ at these horizons, much smaller than the observed one.
\item KLinf-UCB at $q=1.5$: no run shows the failure of \imed{}; the maximum of $\Np_2(T)$ is $16$ at
  $T=10^4$, against $7360$ for \imed{}. One negative reward barely moves its upper confidence bound,
  which is not an empirical mean. Its mean is above $\log T/\Kinf$ at these horizons, which the
  terms $2\log\log t+\log(1+\Np_a(t))+1$ of its exploration function, beyond $\log t$, can explain.
\end{itemize}
These are light Monte Carlo experiments: they illustrate the mechanism and prove nothing.

% ----------------------------------------------------------------
\section{Conclusion and open questions}\label{sec:open}

On semi-bounded rewards with a best mean below $1$, a finite variance of one \term{w:optimal}{optimal} arm suffices for
\imed{} to attain the Burnetas--Katehakis constant $1/\Kinf(\nu_i,\ms)$ for every \term{w:subopt}{suboptimal} arm,
with a finite-time bound, and \imed{} is asymptotically optimal on the laws with a finite moment
of order $r$, for every $r\ge2$. On the scale of power-law lower tails the exponent two is the
threshold for \imed{}: below it, \imed{} is not \term{w:consistent}{consistent}, and at it, \imed{} misses the
lower-bound constant. The limitations are the absence of a rate for the second-order term, the gap
between the square-integrable tails and the tails of order $y^{-2}$, and the restriction to
\imed{}, whose reference level, an empirical mean, is what fails under heavy tails (KLinf-UCB does
not have this weakness, by the paragraph ``Another algorithm'' of Section~\ref{sec:counter}); the
questions below start from them.

Each question below is stated with its class and criterion, and says where the method of this paper
stops; each could be the main result of a next paper.

\begin{openquestion}[The exact boundary for IMED]\label{oq:boundary}
\extends{D9}
Let $P\in\Sm$ with $\mean(P)\in(0,1)$, let $\bar F(y)\coloneqq\Prob_P(X<-y)$, and run \imed{} on the
\term{w:instance}{instance} $(P,\delta_0)$; let $d^\star\coloneqq\log(1/(1-\mean(P)))$. (a) Does $\lim_T\E \Np_2(T)/\log T=1/d^\star$
hold whenever $\int_1^Ty\bar F(y)\,dy=o(\log T)$, for example when $\bar F(y)$ is of order
$1/(y^2\log y)$, where the variance is infinite? (b) If $y^2\bar F(y)\to\kappa>0$, does
$\E \Np_2(T)/\log T$ converge, and is its limit $1/d^\star+2\kappa/(e\,\mean(P)^2)$? A heuristic count of the frozen-index
mechanism (a jump $Y$ at the $s$-th pull costs about $Y/(e\,\mean(P))$ pulls of arm $2$ when $Y\gg s$)
gives this constant; Theorem~\ref{thm:counter}(ii) proves only the excess $\kappa/(4e)$. The method
stops in Proposition~\ref{prop:key}, whose one-jump term uses $\E Y^2$: a version truncated at the
horizon would require the other terms, in particular the H\"older step of Lemma~\ref{lem:moments}, to
be redone with constants uniform in $T$.
\end{openquestion}

\begin{openquestion}[An index of the type of IMED that needs only finite means]\label{oq:finite-mean}
\extends{D9}
For an arm $a$ and a round $t\ge\Karms$, let $n\coloneqq\Np_a(t)$ and
\begin{equation*}
  \rho_a(t)\coloneqq\frac1n\sum_{l\le n}\max\{\Xs_{a,l},-\sqrt n\},
\end{equation*}
the empirical mean of the rewards of arm $a$ truncated below at $-\sqrt n$: one outlier moves it by at
most $(1+\sqrt n)/n$. Let \texttt{IMED}-$\rho$ be the \term{w:learner}{learner} of
Definition~\ref{def:imed} in which $\hmt_a(t)$ is replaced by $\rho_a(t)$, both in the reference level
$\hms(t)$ and in the definition of the \term{w:leader}{empirical leaders}. Is \texttt{IMED}-$\rho$
\term{w:ufc}{uniformly fast convergent} on $\Mm_1=\Sm\setminus\{\delta_1\}$, and does it satisfy
$\limsup_T\E \Np_i(T)/\log T\le1/\Kinf(\nu_i,\ms)$ for every $\nu\in\Mm_1^\arms$ and every
\term{w:subopt}{suboptimal} arm $i$? By Lemma~\ref{lem:class} and Theorem~\ref{cite:GMS19-thm1}, this is the
lower-bound constant on $\Mm_1$. Theorem~\ref{thm:counter} rules out \imed{} itself. For KLinf-UCB,
the paragraph ``Another algorithm'' of Section~\ref{sec:counter} reduces the same property to
\citet[Theorem~3.1]{PA26}; KLinf-UCB inverts $\Kinf$: for each arm at each round, it computes
the largest level $\mu$ with $\Np_a(t)\Kinf(\hnut_a(t),\mu)$ below a threshold of order $\log t$. The index of \texttt{IMED}-$\rho$
evaluates $\Kinf$ at one given level per arm, with no inversion, and the question is whether it keeps
the constant. (A reference level defined by an inversion, such as
$\sup\{\mu<1:\Np_a(t)\Kinf(\hnut_a(t),\mu)\le1\}$, would cost one inversion per arm and per round, as
KLinf-UCB does.) A robust estimate of the mean as the reference of a policy of the minimum empirical
divergence type has a precedent: \citet[Remark~1 and Appendix~B.4]{BSH23} allow a truncated mean in
MED, for laws with a known bound on a moment of order $1+\varepsilon$; here no moment bound is known
and only the lower tail is truncated. The method of this paper stops at
Proposition~\ref{prop:key}, which needs a finite variance to bound the rounds at which the reference
level is low. Lemma~\ref{lem:sup}(b) transfers with the event that the truncated mean of the first
$s$ rewards of arm $j$ is at most $\mu$ in place of $\{\hm_{j,s}\le\mu\}$, but on that event the empirical law of the \term{w:optimal}{optimal} arm may still
contain outliers, and a bound on the expectation of the supremum of Lemma~\ref{lem:sup}(b) over that
event is what is missing.
\end{openquestion}

\begin{openquestion}[The second-order term under a finite variance]\label{oq:second-order}
\extends{D9}
Under the assumptions of Theorem~\ref{thm:main}, is
$\E \Np_i(T)-\log T/\Kinf(\nu_i,\ms)=O(1)$, as \citet[Theorem~5]{HT15} proves when every arm has a
moment generating function near $0$? Or at least $O((\log T)^{\omega})$ for an explicit $\omega<1$,
as \citet[Corollary~4]{HT15} gives with $\omega=10/11$ under an exponential moment of one \term{w:optimal}{optimal} arm?
Theorem~\ref{thm:main} gives only $o(\log T)$: optimising $\delta$ requires the dependence of
$\Gam(\nu_j,\ms-\delta)$ on $\delta$, which goes through the exponent $\gamma$ of
Lemma~\ref{lem:moments} (with $\gamma-1$ of order $\delta$) and the parameters of
\eqref{eq:gammabound}, and that of $\Lt_i(\delta)$, neither of which we tracked.
\end{openquestion}

\begin{openquestion}[Non-parametric Thompson sampling on semi-bounded laws]\label{oq:thompson}
\extends{D10}
The non-parametric Thompson sampling of \citet[Section~3]{RH20} draws, for each arm,
a Dirichlet-weighted average of its rewards and of the point $1$, and is defined on
$(-\infty,1]$ as well. Is it asymptotically optimal on $\Mm_2^\arms$, in the sense of
Corollary~\ref{cor:optimal}(b)? Is it inconsistent on \term{w:instance}{instances} like those of
Theorem~\ref{thm:counter}(i), with a lower tail $y^{-q}$, $1<q<2$? Its sampled values of the optimal
arm are also moved by one large negative reward, but by a random weight instead of $1/N$, so the
mechanism of Lemma~\ref{lem:jump} does not transfer directly.
\end{openquestion}


% ================================================================ appendices
\appendix

\section{Useful lemmas}\label{sec:cited}

In this section we present the result from another paper that our proofs use: the dual formula of
$\Kinf$ on semi-bounded laws (Lemma~\ref{cite:HT15-thm2}), used throughout. Two other published
results are recalled in the main text: the finite-time bound of \citet{HT15}
(Theorem~\ref{cite:HT15-thm3}), which we compare with and do not use, and the lower bound of
\citet{GMS19} (Theorem~\ref{cite:GMS19-thm1}), used in Corollaries~\ref{cor:optimal}
and~\ref{cor:pareto}; the notation of \citet{HT15} is listed before Theorem~\ref{cite:HT15-thm3}. We
also use by name a few textbook facts: the inequalities of Markov, Chebyshev, Jensen, H\"older and
Bonferroni, Doob's maximal inequality for nonnegative supermartingales, monotone and dominated
convergence, Fatou's lemma, and the weak and strong laws of large numbers.

\citet{HT15} write $\mathcal A$ for the set of all laws on $(-\infty,1]$, $F$ for a law, $\mathrm
E(F)$ for its mean, $D_{\inf}(F,\mu;\mathcal A)$ for $\inf\{D(F\|G):G\in\mathcal A,\ \mathrm
E(G)>\mu\}$, which is $\Kinf(F,\mu)$ when $F\in\Sm$, $\nu$ for the dual variable $\lambda$, and
$L_{\max}(F,\mu)$ for the right-hand side below.

% Verbatim: literature/notes/HT15.md, Theorem 2.
\begin{lemma}[{From \citet[Theorem~2(ii)]{HT15}}]\label{cite:HT15-thm2}
\uses{def:model, def:kinf}
For every $P\in\Sm$ and every $\mu<1$,
$\Kinf(P,\mu)=\max_{0\le\lambda\le1/(1-\mu)}\E_P[\log(1-\lambda(X-\mu))]$.
\end{lemma}

\section{Properties of \texorpdfstring{$\mathcal{K}_{\inf}$}{Kinf} on semi-bounded laws}\label{sec:dual}

This appendix derives from the dual formula (Lemma~\ref{cite:HT15-thm2}) the properties of $\Kinf$
used in the proofs. They are ordered by use; the table says where each one serves.


\begin{center}
\small
\begin{tabular}{@{}p{0.4\textwidth}p{0.54\textwidth}@{}}
\toprule
Lemma & Used in \\
\midrule
\ref{lem:dualfacts} (elementary properties) & the remark of Section~\ref{sec:upper-ht15}, Lemmas~\ref{lem:sup}, \ref{lem:termB}, \ref{lem:jump}, \ref{lem:class}, \ref{lem:explore}, \ref{lem:moments}, Theorem~\ref{thm:counter}, Corollary~\ref{cor:pareto}, Proposition~\ref{prop:key} \\
\ref{lem:optimizer} (dual maximiser) & Definition~\ref{def:rate}, Theorem~\ref{thm:main}, Lemmas~\ref{lem:termB}, \ref{lem:termC}, \ref{lem:subopt}, \ref{lem:class} \\
\ref{lem:shift} (shifting the level) & Lemma~\ref{lem:termB} \\
\ref{lem:subopt} (deviations of a \term{w:subopt}{suboptimal} arm) & Lemma~\ref{lem:termC} \\
\ref{lem:class} (independence of the moment class) & Corollary~\ref{cor:optimal} \\
\ref{lem:pa-assumptions} (concentration with a finite mean only) & the paragraph ``Another algorithm'' of Section~\ref{sec:counter}, Open question~\ref{oq:finite-mean} \\
\bottomrule
\end{tabular}
\end{center}

\begin{lemma}[Elementary properties]\label{lem:dualfacts}
\uses{def:model, def:kinf}
Let $P\in\Sm$.
\begin{enumerate}[label=(\alph*)]
\item (Monotonicity in the level.) If $\mu'\le\mu$, then $\Kinf(P,\mu')\le\Kinf(P,\mu)$.
\item (Zero at or above the mean.) If $\mean(P)>\mu$, or if $\mean(P)=\mu<1$, then $\Kinf(P,\mu)=0$. If
  $\mu\ge1$, then $\Kinf(P,\mu)=+\infty$.
\item (Empirical dual.) Let $\mu<1$, $s\ge1$ and $x_1,\dots,x_s\le1$. Then
  \[
    \exp\big(s\Kinf(\hP[x_{1:s}],\mu)\big)=\max_{0\le\lambda\le1/(1-\mu)}\prod_{l=1}^s\big(1-\lambda(x_l-\mu)\big).
  \]
  For every $\mu\in\R$, the map $x_{1:s}\mapsto\Kinf(\hP[x_{1:s}],\mu)$ is Borel on $(-\infty,1]^s$.
\item (Point mass.) $\Kinf(\delta_0,\mu)=\log\frac1{1-\mu}$ for $\mu\in[0,1)$.
\end{enumerate}
\end{lemma}

\begin{proof}
\uses{def:model, def:kinf, cite:HT15-thm2}
(a) If $\mu'\le\mu$, every law with mean above $\mu$ has mean above $\mu'$, so the infimum
defining $\Kinf(P,\mu')$ is over a larger set.

(b) If $\mean(P)>\mu$, take $Q=P$: $\KL(P,P)=0$. If $\mean(P)=\mu<1$, let $\varepsilon\in(0,1)$ and
$Q_\varepsilon\coloneqq(1-\varepsilon)P+\varepsilon\delta_1\in\Sm$. Its mean is
$\mu+\varepsilon(1-\mu)>\mu$, and $dP/dQ_\varepsilon\le1/(1-\varepsilon)$, so
$\KL(P,Q_\varepsilon)\le\log\frac1{1-\varepsilon}$, which tends to $0$. If $\mu\ge1$, no law of $\Sm$
has a mean above $\mu$, and $\inf\emptyset=+\infty$.

(c) The empirical law $\hP\coloneqq\hP[x_{1:s}]$ belongs to $\Sm$, and
$\E_{\hP}\log(1-\lambda(X-\mu))=\frac1s\log\prod_l(1-\lambda(x_l-\mu))$, with $\log0\coloneqq-\infty$.
Lemma~\ref{cite:HT15-thm2} and the monotonicity of $\exp$ (with $\exp(-\infty)\coloneqq0$) give the
identity. For measurability: if $\mu\ge1$ the map is the constant $+\infty$ by (b). If $\mu<1$,
the product is a polynomial in $\lambda$ whose coefficients are polynomials in $x_{1:s}$, so its
maximum over $[0,1/(1-\mu)]$ equals its supremum over the countable set of rational $\lambda$ in that
interval together with $1/(1-\mu)$; this is a Borel function of $x_{1:s}$, and so is its logarithm
divided by $s$.

(d) By Lemma~\ref{cite:HT15-thm2}, $\Kinf(\delta_0,\mu)=\max_{0\le\lambda\le1/(1-\mu)}\log(1+\lambda\mu)$.
For $\mu\ge0$ the function is nondecreasing in $\lambda$, so the maximum is at $1/(1-\mu)$, where it
equals $\log(1+\mu/(1-\mu))=\log\frac1{1-\mu}$.
\end{proof}

\begin{lemma}[The dual maximiser]\label{lem:optimizer}
\uses{def:model, def:kinf}
Let $P\in\Sm$ and $\mean(P)<\mu<1$; let \dt{s:lmax}$\lmax\coloneqq1/(1-\mu)$ and
$L(\lambda)\coloneqq\E_P\log(1-\lambda(X-\mu))$ for $\lambda\in[0,\lmax]$. Then:
\begin{enumerate}[label=(\alph*)]
\item $L$ is well defined with values in $[-\infty,\infty)$, finite on $[0,\lmax)$, concave, and
  continuous on $[0,\lmax]$ as a map into $[-\infty,\infty)$; it attains its maximum, and
  $\max L=\Kinf(P,\mu)\in(0,\infty)$;
\item every maximiser $\lambda^\star$ satisfies $\lambda^\star>0$ and
  $\E_P\big[1/(1-\lambda^\star(X-\mu))\big]\le1$, with $1/0\coloneqq+\infty$; in particular
  $1-\lambda^\star(X-\mu)>0$ almost surely. If $\lambda^\star<\lmax$, the expectation equals $1$.
\item for every maximiser $\lambda^\star$, the variable $V\coloneqq\log(1-\lambda^\star(X-\mu))$, for
  $X\sim P$, satisfies $\E e^{-V}\le1$, is integrable, and $\E V=\Kinf(P,\mu)$.
\end{enumerate}
\end{lemma}

\emph{Role.} It defines $\lst(P,\mu)$ in Definition~\ref{def:rate}, gives
$0<\Ki_i<\infty$ in the proof of Theorem~\ref{thm:main} (Appendix~\ref{sec:proof-main}), and gives the integrability that
Lemmas~\ref{lem:subopt} and~\ref{lem:class} need.

\begin{proof}
\uses{def:model, def:kinf, cite:HT15-thm2}
\emph{Setting.} Fix $P$ and $\mu$; write $Z\coloneqq X-\mu$ for $X\sim P$, so $Z\le1-\mu=1/\lmax$
and $\E Z=\mean(P)-\mu<0$. \emph{Goal.} (a), (b) and (c).

\emph{Step 1: values and concavity.} For $\lambda\in[0,\lmax]$ and $z\le1/\lmax$, the factor
$1-\lambda z$ is nonnegative, and $\log(1-\lambda z)\le\log(1+\lmax|z|)\le\lmax|z|$, which is
integrable. So $L(\lambda)\in[-\infty,\infty)$ is well defined. For $\lambda<\lmax$,
$1-\lambda z\ge1-\lambda/\lmax>0$, so the integrand is bounded below and $L(\lambda)$ is finite. Each
$\lambda\mapsto\log(1-\lambda z)$ is concave, hence so is $L$.

\emph{Step 2: continuity and attainment.} As $\lambda\uparrow\lmax$, the integrand increases on
$\{Z\le0\}$ (where it is nonnegative) and decreases on $\{Z>0\}$ (where it is nonpositive and equal
to $0$ at $\lambda=0$). Monotone convergence on each part gives $L(\lambda)\to L(\lmax)$. On the open
interval $(0,\lmax)$, $L$ is concave and finite, hence continuous. At $\lambda=0$: for
$\lambda\in(0,\lmax/2]$ and $z\le1/\lmax$,
\begin{equation*}
  |\log(1-\lambda z)|\le2\lambda|z|,
\end{equation*}
since if $z\le0$, $0\le\log(1+\lambda|z|)\le\lambda|z|$, and if $z>0$, $\lambda z\le1/2$ and
$0\le-\log(1-\lambda z)\le\lambda z/(1-\lambda z)\le2\lambda z$. Hence
$|L(\lambda)|\le2\lambda\E|Z|\to0=L(0)$ as $\lambda\downarrow0$. So $L$ is continuous on the compact
$[0,\lmax]$ into $[-\infty,\infty)$, bounded above by $\E\log(1+\lmax|Z|)<\infty$, and attains its
maximum. By Lemma~\ref{cite:HT15-thm2}, $\max L=\Kinf(P,\mu)<\infty$.

\emph{Step 3: the maximum is positive and $\lambda^\star>0$.} By the bound of Step 2,
$|\log(1-\lambda z)|/\lambda\le2|z|$ for $\lambda\in(0,\lmax/2]$ and $z\le1/\lmax$. By
dominated convergence, $L(\lambda)/\lambda\to-\E Z=\mu-\mean(P)>0$ as $\lambda\downarrow0$. Hence
$L(\lambda)>0=L(0)$ for small $\lambda>0$: the maximum is positive, and $0$ is not a maximiser.

\emph{Step 4: the derivative.} For $\lambda$ in a compact subinterval $[\lambda_1,\lambda_2]$ of
$(0,\lmax)$, $|\partial_\lambda\log(1-\lambda z)|=|z|/(1-\lambda z)$ is at most $1/\lambda_1$ if
$z\le0$ and at most $1/(\lmax(1-\lambda_2/\lmax))$ if $z>0$. Differentiation under the expectation
gives, for $\lambda\in(0,\lmax)$,
\[
  L'(\lambda)=\E\Big[\frac{-Z}{1-\lambda Z}\Big]=\frac1\lambda\Big(1-\E\Big[\frac{1}{1-\lambda Z}\Big]\Big).
\]

\emph{Step 5: the first-order condition.} Let $\lambda^\star$ be a maximiser; $\lambda^\star>0$ by
Step 3. If $\lambda^\star<\lmax$, then $L'(\lambda^\star)=0$, so $\E[1/(1-\lambda^\star Z)]=1$. If
$\lambda^\star=\lmax$, the concave function $L$ is nondecreasing on $[0,\lmax]$, so $L'\ge0$ on
$(0,\lmax)$, that is $\E[1/(1-\lambda Z)]\le1$ for every $\lambda<\lmax$. The integrands are
nonnegative and converge to $1/(1-\lmax Z)\in[0,\infty]$ as $\lambda\uparrow\lmax$, so Fatou's lemma
gives $\E[1/(1-\lmax Z)]\le1$. In both cases the expectation is finite, so $1-\lambda^\star Z>0$
almost surely.

\emph{Step 6: (c).} Let $V\coloneqq\log(1-\lambda^\star Z)$. By (b), $\E e^{-V}=\E[1/(1-\lambda^\star Z)]\le1$.
Since $\lambda^\star\le\lmax$ and $Z\le1/\lmax$, $V^+\le\log(1+\lambda^\star|Z|)\le\lambda^\star|Z|$, which is
integrable because $P$ has a finite mean; and $V^-\le e^{V^-}\le1+e^{-V}$, which is integrable. So
$V$ is integrable and $\E V=L(\lambda^\star)=\max L=\Kinf(P,\mu)$ by (a).
\end{proof}

\begin{lemma}[Shifting the level]\label{lem:shift}
\uses{def:model, def:kinf}
Let $Q\in\Sm$ and $\mu'<\mu<1$. Then
$\Kinf(Q,\mu)\le\Kinf(Q,\mu')+\log\frac{1-\mu'}{1-\mu}\le\Kinf(Q,\mu')+\frac{\mu-\mu'}{1-\mu}$.
\end{lemma}

\emph{Role.} In Lemma~\ref{lem:termB}, it moves the level of the \term{w:subopt}{suboptimal} arm from
$\ms-\delta$ to $\ms$; it replaces \citet[Lemma~7]{HT15} by a direct comparison of the duals.

\begin{proof}
\uses{def:model, def:kinf, cite:HT15-thm2}
\emph{Setting.} Fix $Q$, $\mu'<\mu<1$, and let $\rho\coloneqq(1-\mu)/(1-\mu')\in(0,1)$.
\emph{Goal.} Compare the duals at the two levels.

Let $\lambda\in[0,1/(1-\mu)]$ and $\lambda'\coloneqq\rho\lambda\in[0,1/(1-\mu')]$. For $x\le1$, let
$u\coloneqq1-x\ge0$ and $A\coloneqq1-\lambda(1-\mu)\in[0,1]$. Then
\[
  1-\lambda(x-\mu)=A+\lambda u,\qquad1-\lambda'(x-\mu')=1-\lambda(1-\mu)+\rho\lambda u=A+\rho\lambda u.
\]
Since $\rho A\le A$, we get $\rho(A+\lambda u)\le A+\rho\lambda u$, that is
\[
  \log(1-\lambda(x-\mu))\le\log(1-\lambda'(x-\mu'))+\log\frac1\rho.
\]
Integrating against $Q$ and using Lemma~\ref{cite:HT15-thm2} at the level $\mu'$,
$\E_Q\log(1-\lambda(X-\mu))\le\Kinf(Q,\mu')+\log\frac1\rho$. Maximising over $\lambda$ and using
Lemma~\ref{cite:HT15-thm2} at the level $\mu$ gives the first inequality. The second is
$\log(1+v)\le v$ with $v\coloneqq(\mu-\mu')/(1-\mu)=1/\rho-1$.
\end{proof}

\begin{lemma}[Deviations of a suboptimal arm need no moment]\label{lem:subopt}
\uses{def:model, def:kinf, def:rate}
Let $P\in\Sm$ with $\mean(P)<\mu<1$, $s\ge1$, and let $\hP_s$ be the empirical law of $X_1,\dots,X_s$,
i.i.d.\ of law $P$. Then, for every $u\in\R$,
$\Prob\big(\Kinf(\hP_s,\mu)\le u\big)\le e^{-s\Lt(P,\mu,u)}$. Moreover $\Lt(P,\mu,u)>0$ for every
$u<\Kinf(P,\mu)$.
\end{lemma}

\emph{Role.} The probability bound of Lemma~\ref{lem:termC}, in the proof of Theorem~\ref{thm:main}. It is
\citet[Proposition~10]{HT15}, with the Chernoff parameter restricted to $[-1,0]$, proved here from
the first-order condition of Lemma~\ref{lem:optimizer} only.

\begin{proof}
\uses{def:model, def:kinf, def:rate, cite:HT15-thm2, lem:optimizer}
\emph{Setting.} Fix $P$, $\mu$, $s$; let $\lambda^\star\coloneqq\lst(P,\mu)$,
$V_l\coloneqq\log(1-\lambda^\star(X_l-\mu))\in[-\infty,\infty)$ and, for $\eta\in[-1,0]$,
$\psi(\eta)\coloneqq\log\E e^{\eta V_1}$. \emph{Goal.} The Chernoff bound and the positivity of the
rate.

\emph{Step 1: the moments of $e^{-V}$.} By Lemma~\ref{lem:optimizer}(b), $\E e^{-V_1}\le1$; in
particular $V_1>-\infty$ almost surely and $\psi(0)=0$. For $\eta\in[-1,0]$, Jensen's inequality for
the concave map $y\mapsto y^{-\eta}$ gives
\[
  \E e^{\eta V_1}=\E\big[(e^{-V_1})^{-\eta}\big]\le(\E e^{-V_1})^{-\eta}\le1,
\]
so $\psi\le0$ on $[-1,0]$.

\emph{Step 2: the Chernoff bound.} By Lemma~\ref{cite:HT15-thm2} applied to $\hP_s$, every $\lambda$
gives a lower bound, so $\Kinf(\hP_s,\mu)\ge\frac1s\sum_{l\le s}V_l$. For $\eta\in[-1,0)$, Markov's
inequality and independence give
\[
  \Prob\Big(\frac1s\sum_{l\le s}V_l\le u\Big)=\Prob\Big(e^{\eta\sum_lV_l}\ge e^{\eta su}\Big)\le e^{-\eta su}\,
  e^{s\psi(\eta)}=e^{-s(\eta u-\psi(\eta))}.
\]
For $\eta=0$ the bound $e^{0}=1$ is trivial. Taking the infimum over $\eta\in[-1,0]$ gives the
first claim.

\emph{Step 3: $\E V_1=\Kinf(P,\mu)$ and $V_1$ is integrable.} This is Lemma~\ref{lem:optimizer}(c).
Moreover $\E e^{V_1^-}\le1+\E e^{-V_1}\le2$.

\emph{Step 4: positivity.} For $\eta\in[-1/2,0)$, $|e^{\eta V_1}-1|/|\eta|\le|V_1|e^{V_1^-/2}$, by
$|e^x-1|\le|x|e^{x^+}$. The right-hand side is integrable: it equals $V_1^+$ where $V_1\ge0$, and is at
most $(2/e)e^{V_1^-}$ where $V_1<0$. By dominated convergence,
$(\E e^{\eta V_1}-1)/\eta\to\E V_1$ as $\eta\uparrow0$, hence $\psi(\eta)/\eta\to\E V_1$. Let
$u<\Kinf(P,\mu)=\E V_1$. Then
\[
  \eta u-\psi(\eta)=\eta\Big(u-\frac{\psi(\eta)}{\eta}\Big),
\]
and for $\eta<0$ close to $0$ the factor $u-\psi(\eta)/\eta$ is negative. Hence $\eta u-\psi(\eta)>0$
for such $\eta$, and $\Lt(P,\mu,u)>0$.
\end{proof}

\begin{lemma}[The constant does not depend on the moment class]\label{lem:class}
\uses{def:model, def:kinf}
Let $r\ge1$, $P\in\Mm_r$ and $\mu<1$. Then $\Kinf(P,\mu;\Mm_r)=\Kinf(P,\mu)$.
\end{lemma}

\emph{Role.} In Corollary~\ref{cor:optimal}(b), it identifies the lower-bound constant on $\Mm_r$
with the upper bound of Theorem~\ref{thm:main}.

\begin{proof}
\uses{def:model, def:kinf, cite:HT15-thm2, lem:dualfacts, lem:optimizer}
\emph{Idea.} Since $\Mm_r\subseteq\Sm$, only $\le$ needs a proof. We exhibit laws of $\Mm_r$ with
mean above $\mu$ whose divergence from $P$ approaches $\Kinf(P,\mu)$: the optimal law $G^\star$ of
the proof of \citet[Theorem~2]{HT15}, mixed with a small mass at $1$. Its density with respect to
$P$ is at most $1$ below $\mu$, so it keeps the moment of order $r$.

\emph{Setting.} Fix $r$, $P$, $\mu$. \emph{Goal.} For every $\varepsilon\in(0,1)$, a law
$G_\varepsilon\in\Mm_r$ with $\mean(G_\varepsilon)>\mu$ and
$\KL(P,G_\varepsilon)\le\Kinf(P,\mu)+\log\frac1{1-\varepsilon}$.

\emph{Step 1: the case $\mean(P)\ge\mu$.} Here $\Kinf(P,\mu)=0$ by Lemma~\ref{lem:dualfacts}(b), since
$\mean(P)<1$ ($P\ne\delta_1$). Let $G_\varepsilon\coloneqq(1-\varepsilon)P+\varepsilon\delta_1$. It is
a law on $(-\infty,1]$ with a finite moment of order $r$ and mean
$\mean(P)+\varepsilon(1-\mean(P))>\mu$; it is not $\delta_1$; and $\KL(P,G_\varepsilon)\le\log\frac1{1-\varepsilon}$
because $dP/dG_\varepsilon\le1/(1-\varepsilon)$.

\emph{Step 2: the law $G^\star$ when $\mean(P)<\mu$.} Let $\lambda^\star$ be a maximiser from
Lemma~\ref{lem:optimizer}, $g(x)\coloneqq1/(1-\lambda^\star(x-\mu))\in(0,\infty]$ for $x\le1$, and
$c\coloneqq\E_Pg(X)\le1$ (Lemma~\ref{lem:optimizer}(b)). Let
\[
  G^\star\coloneqq g\cdot P+(1-c)\,\delta_1,
\]
a probability measure on $(-\infty,1]$. If $c<1$, then $\lambda^\star=\lmax$ by
Lemma~\ref{lem:optimizer}(b), and $P(\{1\})=0$ because $g(1)=+\infty$ there while $\E_Pg<\infty$.

\emph{Step 3: the mean of $G^\star$ is $\mu$.} Where $g$ is finite,
$(x-\mu)g(x)=(g(x)-1)/\lambda^\star$, because $g(x)(1-\lambda^\star(x-\mu))=1$. This function is
$P$-integrable (it is bounded by $1/\lambda^\star$ below $\mu$ and by $(1-\mu)g$ above $\mu$), so
\[
  \int(x-\mu)\,dG^\star(x)=\frac{c-1}{\lambda^\star}+(1-c)(1-\mu)=(1-c)\Big(1-\mu-\frac1{\lambda^\star}\Big)=0,
\]
since either $c=1$, or $\lambda^\star=\lmax=1/(1-\mu)$.

\emph{Step 4: the divergence.} The measure $(1-c)\delta_1$ is carried by a $P$-null set when $c<1$
(Step 2), and $g>0$, so $P$ is absolutely continuous with respect to $G^\star$ with
$dP/dG^\star=1/g$ $P$-almost surely. Hence, by Lemma~\ref{lem:optimizer}(a),
\[
  \KL(P,G^\star)=\E_P\log\frac1{g(X)}=\E_P\log(1-\lambda^\star(X-\mu))=\Kinf(P,\mu).
\]

\emph{Step 5: the moment of order $r$.} On $\{x\le\mu\}$, $g(x)=1/(1+\lambda^\star(\mu-x))\le1$. On
$\{\mu<x\le1\}$, $|x|\le\max\{1,|\mu|\}$. Since $\int_{\{x>\mu\}}g\,dP\le c$,
\[
  \int|x|^r\,dG^\star\le\E_P|X|^r+\max\{1,|\mu|^r\}\,c+(1-c)\le\E_P|X|^r+\max\{1,|\mu|^r\}<\infty.
\]

\emph{Step 6: mixing with $\delta_1$.} Let $G_\varepsilon\coloneqq(1-\varepsilon)G^\star+\varepsilon\delta_1$.
It is a law on $(-\infty,1]$ with a finite moment of order $r$, its mean is
$\mu+\varepsilon(1-\mu)>\mu$, so it is not $\delta_1$ and belongs to $\Mm_r$, and
$dP/dG_\varepsilon\le(1-\varepsilon)^{-1}dP/dG^\star$, so
$\KL(P,G_\varepsilon)\le\Kinf(P,\mu)+\log\frac1{1-\varepsilon}$ by Step 4.

\emph{Step 7: conclusion.} In both cases, $\Kinf(P,\mu;\Mm_r)\le\Kinf(P,\mu)+\log\frac1{1-\varepsilon}$
for every $\varepsilon\in(0,1)$; letting $\varepsilon\to0$ gives the claim.
\end{proof}

\begin{lemma}[Concentration with a finite mean only]\label{lem:pa-assumptions}
\uses{def:model, def:kinf, def:rate}
Let $P\in\Mm_1$, let $X_1,X_2,\dots$ be i.i.d.\ of law $P$, and let $\hP_n\coloneqq\hP[X_{1:n}]$.
\begin{enumerate}[label=(\alph*)]
\item (Upper deviations at the mean, uniformly in time.) For every $x\in\R$,
  \[
    \Prob\big(\exists n\ge1:\ n\Kinf(\hP_n,\mean(P))-\log(n+1)\ge x\big)\le e^{-x}.
  \]
\item (Lower deviations above the mean.) Let $\mean(P)<\mu<1$, $\lambda^\star\coloneqq\lst(P,\mu)$,
  $V\coloneqq\log(1-\lambda^\star(X_1-\mu))$ and $C\coloneqq\E\big[V^2e^{V^-/2}\big]$. Then
  $0<C<\infty$, and for every $n\ge1$ and every $d\in(0,C/2]$,
  \[
    \Prob\big(\Kinf(\hP_n,\mu)\le\Kinf(P,\mu)-d\big)\le\exp\Big(-\frac{n\,d^2}{2C}\Big).
  \]
\item (The non-strict inequality.) For every $Q\in\Sm$ and every $\mu\in\R$,
  \[
    \Kinf(Q,\mu)=\inf\{\KL(Q,G):\ G\in\Mm_1,\ \mean(G)\ge\mu\}.
  \]
\end{enumerate}
\end{lemma}

\emph{Role.} In the paragraph ``Another algorithm'' of Section~\ref{sec:counter}: (c) identifies the
divergence of \citet{PA26} over the class $\Mm_1$ with $\Kinf$, and (a) and (b) are their
Assumptions~2.3 (with $g(n)=\log(1+n)+C_2$, $C_2>0$) and~2.4 (with $c_\nu=1/(2C)$, $d_0=C/2$, at the
levels below $1$). It is not used in the proofs of the other results.

\begin{proof}
\uses{def:model, def:kinf, def:rate, lem:dualfacts, lem:optimizer, lem:subopt}
\emph{Idea.} By the dual formula, $\exp(n\Kinf(\hP_n,\mean(P)))$ is the maximum over the dual variable
of products that are nonnegative martingales of mean $1$ under a finite mean only. Their average over
the dual variable is again such a martingale, and it is at least the maximum divided by $n+1$, so
Doob's maximal inequality gives (a). The products are the wealths of the constant-rebalanced
portfolios on two assets of price relatives $1$ and $R_l$ (Step 1), their average is the universal
portfolio with the uniform prior, and the factor $n+1$ is the bound of \citet[Theorem~1]{BK99} for
two assets (see also \citealp{CO96}, to whom they credit it, and \citealp{Cov91}); Step 3 follows
their proof.
That such products are nonnegative martingales of mean $1$ for unbounded observations, as soon as
each factor is nonnegative, is a property of betting martingales: see also
\citet[Proposition~3 in arXiv v7]{WSR24}, and \citet[Section~2]{Min22}, who notes that these
constructions give a one-sided confidence sequence for the mean of observations unbounded on one
side with a bounded mean. The argument as a whole is the method of
\citet[Proposition~5, proof in Appendix~B.4.1]{AJK21}, for laws on $\R$ with a known bound on a
moment of order $1+\varepsilon$: a uniform mixture over the two dual variables, a factor
$2\log(1+n)+1$ between the logarithms of the maximum and of the mixture (their Lemma~17), and
Ville's inequality. \citet[Appendix~A.1]{PA26} state that the same argument, with the dual of
\citet{HT10}, gives Assumption~2.3 for laws on a bounded interval, with $g(n)=\log(1+n)+1$. What (a)
adds is the time-uniform bound on $n\Kinf(\hP_n,\mean(P))$ for semi-bounded laws: only an upper
bound on the rewards and a finite mean are needed, which is the form that Assumption~2.3 of
\citet{PA26} takes on the class $\Mm_1$.
Part (b) is the Chernoff bound
of Lemma~\ref{lem:subopt}, whose rate is quadratic near the mean. Part (c) moves mass to $1$ or away
from it.

\emph{Setting.} Let $m\coloneqq\mean(P)<1$ (since $P\ne\delta_1$) and
$R_l\coloneqq(1-X_l)/(1-m)\ge0$, so $\E R_l=1$. \emph{Goal.} (a), (b), (c).

\emph{Step 1: (a), the products.} For $\theta\in[0,1]$ and $\lambda\coloneqq\theta/(1-m)$,
\[
  1-\lambda(X_l-m)=(1-\theta)+\theta R_l\ge0 .
\]
Let $W_n(\theta)\coloneqq\prod_{l\le n}\big((1-\theta)+\theta R_l\big)$. By
Lemma~\ref{lem:dualfacts}(c) (\emph{empirical dual}) at the level $m<1$,
\[
  \exp\big(n\Kinf(\hP_n,m)\big)=\max_{\theta\in[0,1]}W_n(\theta).
\]

\emph{Step 2: (a), the mixture is a martingale.} Let $M_0\coloneqq1$ and
$M_n\coloneqq\int_0^1W_n(\theta)\,d\theta$; $(\theta,\omega)\mapsto W_n(\theta)$ is a polynomial in
$\theta$ with measurable coefficients. Let $\mathcal F_n\coloneqq\sigma(X_1,\dots,X_n)$ and
$A\in\mathcal F_{n-1}$. For every $\theta$, $W_n(\theta)=W_{n-1}(\theta)\big((1-\theta)+\theta R_n\big)$
with $R_n$ independent of $\mathcal F_{n-1}$ and $\E[(1-\theta)+\theta R_n]=1$, so
$\E[W_n(\theta)\ind_A]=\E[W_{n-1}(\theta)\ind_A]$. By Tonelli's theorem,
\[
  \E[M_n\ind_A]=\int_0^1\E[W_n(\theta)\ind_A]\,d\theta=\int_0^1\E[W_{n-1}(\theta)\ind_A]\,d\theta=\E[M_{n-1}\ind_A].
\]
So $(M_n)_{n\ge0}$ is a nonnegative martingale with $M_0=1$.

\emph{Step 3: (a), maximum against mixture.} Fix the outcome and $n$, and let $\theta^\star$ maximise the
continuous function $W_n$ on $[0,1]$. For $\theta\in[\theta^\star,1]$, write
$\theta=(1-\alpha)\theta^\star+\alpha\cdot1$ with $\alpha\coloneqq(\theta-\theta^\star)/(1-\theta^\star)\in[0,1]$
(when $\theta^\star<1$); for $\theta\in[0,\theta^\star]$, write $\theta=(1-\alpha)\theta^\star+\alpha\cdot0$
with $\alpha\coloneqq(\theta^\star-\theta)/\theta^\star\in[0,1]$ (when $\theta^\star>0$). Each factor of
$W_n$ is affine in $\theta$ and nonnegative at $\theta=0$ and $\theta=1$, so
\[
  (1-\theta)+\theta R_l\ge(1-\alpha)\big((1-\theta^\star)+\theta^\star R_l\big),
  \qquad\text{hence}\qquad W_n(\theta)\ge(1-\alpha)^nW_n(\theta^\star).
\]
Integrating over the two intervals,
\begin{align*}
  M_n&\ge W_n(\theta^\star)\Big(\int_{\theta^\star}^1\Big(1-\frac{\theta-\theta^\star}{1-\theta^\star}\Big)^nd\theta
  +\int_0^{\theta^\star}\Big(1-\frac{\theta^\star-\theta}{\theta^\star}\Big)^nd\theta\Big)\\
  &=W_n(\theta^\star)\Big(\frac{1-\theta^\star}{n+1}+\frac{\theta^\star}{n+1}\Big)=\frac{W_n(\theta^\star)}{n+1},
\end{align*}
where an integral over an interval of length $0$ is $0$.

\emph{Step 4: (a), conclusion.} If $x<0$ the bound exceeds $1$. Let $x\ge0$. By Steps 1 and 3,
$n\Kinf(\hP_n,m)-\log(n+1)\ge x$ implies $M_n\ge e^x$. Doob's maximal inequality for the nonnegative
martingale $(M_n)$ gives $\Prob(\max_{n\le N}M_n\ge e^x)\le e^{-x}\E M_0=e^{-x}$ for every $N$, and
letting $N\to\infty$ gives (a).

\emph{Step 5: (b), the constant $C$.} By Lemma~\ref{lem:optimizer}(c),
$\E e^{-V}\le1$, $V$ is integrable and $\E V=\Kinf(P,\mu)$, which is positive by
Lemma~\ref{lem:optimizer}(a). On $\{V\ge0\}$, $V\le\log(1+\lambda^\star|X_1-\mu|)\le(\lambda^\star|X_1-\mu|)^{1/2}$,
because $\log(1+t^2)\le t$ for $t\ge0$; so $V^2\le\lambda^\star|X_1-\mu|$, which is integrable. On
$\{V<0\}$, $V^2e^{V^-/2}\le(16/e^2)e^{V^-}=(16/e^2)e^{-V}$, because $v^2e^{-v/2}\le16/e^2$ for $v\ge0$.
Hence $C\le\lambda^\star\E|X_1-\mu|+16/e^2<\infty$. If $C=0$, then $V=0$ almost surely and
$\Kinf(P,\mu)=\E V=0$, a contradiction; so $C>0$.

\emph{Step 6: (b), a quadratic bound on the log-moment.} For $y\in\R$, Taylor's formula gives
$e^y\le1+y+\frac{y^2}{2}e^{y^+}$. For $\eta\in[-1/2,0]$ and $y=\eta V$, $(\eta V)^+=|\eta|V^-\le V^-/2$, so
\[
  \E e^{\eta V}\le1+\eta\Kinf(P,\mu)+\frac{\eta^2C}{2},
  \qquad\text{hence}\qquad
  \log\E e^{\eta V}\le\eta\Kinf(P,\mu)+\frac{\eta^2C}{2},
\]
by $\log(1+z)\le z$.

\emph{Step 7: (b), conclusion.} Let $d\in(0,C/2]$, $u\coloneqq\Kinf(P,\mu)-d$ and
$\eta\coloneqq-d/C\in[-1/2,0)$. In the rate of Definition~\ref{def:rate}, $\E_P[(1-\lst(P,\mu)(X-\mu))^\eta]=\E e^{\eta V}$, so by Step 6
\[
  \Lt(P,\mu,u)\ge\eta u-\log\E e^{\eta V}\ge-\eta d-\frac{\eta^2C}{2}=\frac{d^2}{C}-\frac{d^2}{2C}=\frac{d^2}{2C}.
\]
Lemma~\ref{lem:subopt} with $s=n$ gives $\Prob(\Kinf(\hP_n,\mu)\le u)\le e^{-n\Lt(P,\mu,u)}\le e^{-nd^2/(2C)}$.

\emph{Step 8: (c).} Let $\Kinf^{\ge}(Q,\mu)$ be the right-hand side. If $\mu\ge1$, both sides are
$+\infty$: no law of $\Mm_1$ has a mean at least $1$, and Lemma~\ref{lem:dualfacts}(b) applies. Let
$\mu<1$. First, $\Kinf(Q,\mu)\le\Kinf^{\ge}(Q,\mu)$: for $G\in\Mm_1$ with $\mean(G)\ge\mu$ and
$\varepsilon\in(0,1)$, the law $G_\varepsilon\coloneqq(1-\varepsilon)G+\varepsilon\delta_1$ belongs to
$\Sm$, has the mean $\mean(G)+\varepsilon(1-\mean(G))>\mu$ because $\mean(G)<1$, and satisfies
$\KL(Q,G_\varepsilon)\le\KL(Q,G)+\log\frac1{1-\varepsilon}$ because $dG_\varepsilon\ge(1-\varepsilon)dG$;
let $\varepsilon\to0$. Second, $\Kinf^{\ge}(Q,\mu)\le\Kinf(Q,\mu)$: let $G\in\Sm$ with $\mean(G)>\mu$
and $\KL(Q,G)<\infty$. If $G\ne\delta_1$, then $G\in\Mm_1$ and $\mean(G)\ge\mu$. If $G=\delta_1$, then
$Q=\delta_1$, and for $\varepsilon\in(0,\min\{1,1-\mu\})$ the law $(1-\varepsilon)\delta_1+\varepsilon\delta_0\in\Mm_1$
has the mean $1-\varepsilon\ge\mu$ and the divergence $\log\frac1{1-\varepsilon}$ from $Q$, which tends to
$0=\KL(Q,G)$. Taking the infimum over $G$ gives the inequality.
\end{proof}

% ----------------------------------------------------------------
\section{Proof of Theorem~\ref{thm:main}}\label{sec:proof-main}

Throughout this appendix, $\nu\in\Sm^\arms$ is an \term{w:instance}{instance} with $\ms<1$, $j$ an \term{w:optimal}{optimal} arm with
$\nu_j\in\Sm_2$, $i$ a \term{w:subopt}{suboptimal} arm, $\delta\in(0,(1-\ms)\Ki_i/2)$, $T\ge1$, and \imed{} is run with any
\term{w:tie}{tie rule}. We use the level
\begin{equation*}
  \mu\coloneqq\ms-\delta.
\end{equation*}
It is below $\marm_j=\ms$, so that the \term{w:optimal}{optimal} arm $j$ has a mean above $\mu$ and
Proposition~\ref{prop:key} applies to it; and it is close to $\ms$, so that moving the level of arm $i$
from $\mu$ to $\ms$ costs only $\delta/(1-\ms)$ (Lemma~\ref{lem:shift}). Since $\marm_i<\ms<1$,
Lemma~\ref{lem:optimizer}(a) (\emph{the dual maximiser}) gives $0<\Ki_i<\infty$. \dt{s:ui}The level of
the deviations of arm $i$ is
\begin{equation}\label{eq:ui}
  \ui_i\coloneqq\Ki_i-\frac{\delta}{1-\ms}.
\end{equation}
\dt{s:MT}The number of pulls of arm $i$ paid at first order is
\begin{equation}\label{eq:MT}
  \MT_T\coloneqq\frac{\log T}{\Ki_i-2\delta/(1-\ms)}.
\end{equation}
Since $0<\delta<(1-\ms)\Ki_i/2$, we have $\Ki_i/2<\ui_i<\Ki_i$, $\Ki_i-2\delta/(1-\ms)>0$ and $\MT_T\ge0$.
We will often use the fundamental inequality of index policies: for every arm $a$ and every $s\ge1$,
since $\Np_a(t)=s$ and $\Ap_t=a$ imply $\Np_a(t+1)=s+1$,
\begin{equation}\label{eq:fundamental}
  \sum_{t\ge0}\ind\{\Ap_t=a,\ \Np_a(t)=s\}\le1.
\end{equation}

To prove Theorem~\ref{thm:main} we use the following decomposition.

\begin{lemma}[Decomposition]\label{lem:decomp}
\uses{def:model, def:run, def:imed}
Let $\nu\in\Sm^\arms$, let $i$ be an arm, $\mu\in\R$ and $T\ge1$, and consider \imed{} with any \term{w:tie}{tie rule}.
Then on every run,
\dt{s:TA}\dt{s:TB}
\begin{equation*}
  \Np_i(T)\le1+\underbrace{\sum_{t=\Karms}^{T-1}\ind\{\Ap_t=i,\ \hms(t)\le\mu\}}_{\TA}
  +\underbrace{\sum_{t=\Karms}^{T-1}\ind\{\Ap_t=i,\ \hms(t)>\mu\}}_{\TB}.
\end{equation*}
\end{lemma}

(See the proof in Subsection~\ref{sec:proof-lem-decomp}.)

From this decomposition we use the following lemmas to control each term.

\begin{lemma}[The rounds at which every empirical mean is low]\label{lem:termA}
\uses{def:model, def:kinf, def:run, def:imed, def:dev, lem:decomp}
Let $\nu\in\Sm^\arms$, let $i,j$ be arms, $\mu\in\R$ and $T\ge1$, and consider \imed{} with any \term{w:tie}{tie rule}.
Then, with $\TA$ as in Lemma~\ref{lem:decomp}, on every run
\begin{equation*}
  \TA\le\sup_{s\ge1}\ws_s(\Xs_{j,\cdot};\mu),\qquad\text{hence}\qquad\E\TA\le\Gam(\nu_j,\mu).
\end{equation*}
If moreover $\nu_j\in\Sm_2$ and $\mu<\marm_j$, then $\Gam(\nu_j,\mu)<\infty$.
\end{lemma}

(See the proof in Subsection~\ref{sec:proof-lem-termA}.) This is the new step:
\citet[(25)--(26)]{HT15} bound the same term by a sum over the sample sizes of arm $j$, whose
expectation is infinite without a third moment (Lemma~\ref{lem:three-moments});
Lemma~\ref{lem:sup}(b) counts it by a supremum, and Proposition~\ref{prop:key} bounds the expectation
of that supremum under a finite variance.

\begin{lemma}[The rounds at which some empirical mean is high]\label{lem:termB}
\uses{def:model, def:kinf, def:run, def:imed, def:rate, lem:decomp}
Let $\nu\in\Sm^\arms$ with $\ms<1$, let $i$ be a \term{w:subopt}{suboptimal} arm, $\delta\in(0,(1-\ms)\Ki_i/2)$,
$\mu\coloneqq\ms-\delta$ and $T\ge1$, and let $\ui_i$ and $\MT_T$ be given by \eqref{eq:ui}
and~\eqref{eq:MT}. Consider \imed{} with any \term{w:tie}{tie rule}. Then, with $\TB$ as in Lemma~\ref{lem:decomp}, on
every run
\dt{s:TC}
\begin{equation*}
  \TB\le\sum_{s\ge1}\ind\big\{s\Kinf(\hnu_{i,s},\mu)\le\log T\big\}
  \le \MT_T+\underbrace{\sum_{s\ge1}\ind\big\{\Kinf(\hnu_{i,s},\ms)\le \ui_i\big\}}_{\TC}.
\end{equation*}
\end{lemma}

(See the proof in Subsection~\ref{sec:proof-lem-termB}.) This is the classical argument for the
\term{w:subopt}{suboptimal} arm, as in the proof of \citet[Theorem~3]{HT15}; the only change is the level shift of
Lemma~\ref{lem:shift}, a direct comparison of the duals that replaces \citet[Lemma~7]{HT15}.

\begin{lemma}[The deviations of the suboptimal arm]\label{lem:termC}
\uses{def:model, def:kinf, def:run, def:rate, lem:termB}
Under the hypotheses of Lemma~\ref{lem:termB}, $\Lt_i(\delta)>0$ and
\begin{equation*}
  \E\TC\le\frac{1}{1-e^{-\Lt_i(\delta)}}.
\end{equation*}
\end{lemma}

(See the proof in Subsection~\ref{sec:proof-lem-termC}.) This is \citet[Proposition~10]{HT15}, as in
the classical proof; no contribution here, except that Lemma~\ref{lem:subopt} restricts the Chernoff
parameter to $[-1,0]$ and needs no moment of $\nu_i$ beyond its mean.

From these lemmas the proof of the theorem is short.

\begin{proof}[Proof of Theorem~\ref{thm:main}]
\proves{thm:main}
\uses{def:model, def:kinf, def:run, def:imed, def:dev, def:rate, lem:optimizer, lem:decomp, lem:termA, lem:termB, lem:termC}
\emph{The constants.} Since $\marm_i<\ms<1$, Lemma~\ref{lem:optimizer}(a) gives $0<\Ki_i<\infty$. By
Lemma~\ref{lem:termC}, $\Lt_i(\delta)>0$. Since $\nu_j\in\Sm_2$ and $\mu=\ms-\delta<\ms=\marm_j$,
Lemma~\ref{lem:termA} gives $\Gam(\nu_j,\ms-\delta)<\infty$.

\emph{The bound.} Taking expectations in Lemma~\ref{lem:decomp} at the level $\mu=\ms-\delta$, then
using Lemmas~\ref{lem:termA}, \ref{lem:termB} and~\ref{lem:termC} in turn, we get
\begin{align*}
  \E \Np_i(T)&\le1+\E\TA+\E\TB\\
  &\le1+\Gam(\nu_j,\ms-\delta)+\MT_T+\E\TC\\
  &\le1+\frac{\log T}{\Ki_i-2\delta/(1-\ms)}+\frac{1}{1-e^{-\Lt_i(\delta)}}+\Gam(\nu_j,\ms-\delta),
\end{align*}
where the last line uses the definition~\eqref{eq:MT} of $\MT_T$. This is the bound of
Theorem~\ref{thm:main}.
\end{proof}

\subsection{Proof of Lemma~\ref{lem:decomp}}\label{sec:proof-lem-decomp}

\begin{proof}
\proves{lem:decomp}
\uses{def:model, def:run, def:imed}
By Definition~\ref{def:imed}, \imed{} pulls each arm once in the rounds $0,\dots,\Karms-1$. If $T\le \Karms$, arm
$i$ is pulled at most once before round $T$ and both sums are empty, so $\Np_i(T)\le1$. Otherwise,
since $\hms(t)$ is defined for every $t\ge \Karms$, we split each indicator according to the event
$\{\hms(t)\le\mu\}$:
\begin{align*}
  \Np_i(T)&=\sum_{t=0}^{\Karms-1}\ind\{\Ap_t=i\}+\sum_{t=\Karms}^{T-1}\ind\{\Ap_t=i\}\\
  &=1+\sum_{t=\Karms}^{T-1}\ind\{\Ap_t=i,\ \hms(t)\le\mu\}+\sum_{t=\Karms}^{T-1}\ind\{\Ap_t=i,\ \hms(t)>\mu\}.
\end{align*}
This is the claim.
\end{proof}

\subsection{Proof of Lemma~\ref{lem:sup}}\label{sec:proof-lem-sup}

\begin{proof}
\proves{lem:sup}
\uses{def:model, def:kinf, def:run, def:imed, def:dev, lem:dualfacts}
\emph{Setting.} Fix an outcome, arms $i,j$ and reals $x,\mu$. \emph{Goal.} The two counting
bounds.

\emph{Step 1: part (a).} Let $t\ge \Karms$ be a round with $\Ap_t=i$. Since $\Dv_i(t)\ge0$,
\[
  \Is(t)=\Ix_i(t)\ge\log \Np_i(t).
\]
The counts $\Np_i(t)$ at the successive rounds $t$ with $\Ap_t=i$ are distinct positive integers.
Hence the rounds counted in (a) have distinct values $\Np_i(t)\in\{n\ge1:\log n\le x\}$, a set with
at most $e^x$ elements. This proves (a).

\emph{Step 2: a random level.} Let
\[
  M\coloneqq\sup\big\{s\Kinf(\hnu_{j,s},\mu)+\log s:\ s\ge1,\ \hm_{j,s}\le\mu\big\}\in[-\infty,\infty],
\]
with $\sup\emptyset\coloneqq-\infty$. By Definition~\ref{def:dev}, $e^M=\sup_{s\ge1}\ws_s(\Xs_{j,\cdot};\mu)$, with
$e^{-\infty}\coloneqq0$ and $e^{+\infty}\coloneqq+\infty$.

\emph{Step 3: every counted round has $\Is(t)\le M$.} Let $t\ge \Karms$ with $\Ap_t=i$ and
$\hms(t)\le\mu$, and let $s\coloneqq \Np_j(t)\ge1$. Then $\hm_{j,s}=\hmt_j(t)\le\hms(t)\le\mu$,
so $s$ enters the supremum defining $M$. Moreover $\Dv_j(t)\le\Kinf(\hnu_{j,s},\hms(t))$: the
two are equal if $j$ is not an \term{w:leader}{empirical leader}, and $\Dv_j(t)=0$ otherwise. Therefore
\begin{align*}
  \Is(t)\le \Ix_j(t)&=s\Dv_j(t)+\log s\\
  &\le s\Kinf(\hnu_{j,s},\hms(t))+\log s&&\text{(the comparison of $\Dv_j(t)$ above)}\\
  &\le s\Kinf(\hnu_{j,s},\mu)+\log s\le M&&\text{(Lemma~\ref{lem:dualfacts}(a), since $\hms(t)\le\mu$).}
\end{align*}
Hence every round counted in (b) is counted in (a) with $x=M$.

\emph{Step 4: conclusion.} If $M=-\infty$, Step 3 shows that no round is counted in (b). If
$M=+\infty$, (b) is trivial. Otherwise, part (a) holds pathwise for every real $x$, hence for the
random value $x=M$, and the count in (b) is at most $e^M=\sup_s\ws_s(\Xs_{j,\cdot};\mu)$.
\end{proof}

\subsection{Proof of Lemma~\ref{lem:termA}}\label{sec:proof-lem-termA}

\begin{proof}
\proves{lem:termA}
\uses{def:model, def:kinf, def:run, def:imed, def:dev, lem:decomp, lem:sup, prop:key}
Let us look at $\TA$. Each of its rounds $t$ satisfies $t\ge \Karms$, $\Ap_t=i$ and $\hms(t)\le\mu$, so
\begin{align*}
  \TA&=\sum_{t=\Karms}^{T-1}\ind\{\Ap_t=i,\ \hms(t)\le\mu\}\\
  &\le|\{t\ge \Karms:\Ap_t=i,\ \hms(t)\le\mu\}|\\
  &\le\sup_{s\ge1}\ws_s(\Xs_{j,\cdot};\mu)&&\text{(Lemma~\ref{lem:sup}(b))}.
\end{align*}
Now the sequence $\Xs_{j,\cdot}=(\Xs_{j,1},\Xs_{j,2},\dots)$ is i.i.d.\ of law $\nu_j$
(Definition~\ref{def:run}), so by Definition~\ref{def:dev} the expectation of the right-hand side is
$\Gam(\nu_j,\mu)$. This gives $\E\TA\le\Gam(\nu_j,\mu)$. Finally, if $\nu_j\in\Sm_2$ and
$\mu<\marm_j=\mean(\nu_j)$, Proposition~\ref{prop:key} (\emph{the key deviation bound}) gives
$\Gam(\nu_j,\mu)<\infty$.
\end{proof}

\subsection{Proof of Lemma~\ref{lem:termB}}\label{sec:proof-lem-termB}

\begin{proof}
\proves{lem:termB}
\uses{def:model, def:kinf, def:run, def:imed, def:rate, lem:decomp, lem:dualfacts, lem:optimizer, lem:shift}
\emph{Idea.} At a round counted in $\TB$, the index of an \term{w:leader}{empirical leader} is at most $\log T$, which
forces the empirical divergence of arm $i$ at the level $\mu$ below $\log T/\Np_i(t)$; the counts of
distinct rounds are distinct, and moving the level from $\mu$ to $\ms$ costs $\delta/(1-\ms)$.

\emph{Step 1: each counted round has a small divergence.} Let $\Karms\le t\le T-1$ with $\Ap_t=i$ and
$\hms(t)>\mu$, and let $s\coloneqq \Np_i(t)\ge1$. Let $a$ be an \term{w:leader}{empirical leader} at round $t$. By
Definition~\ref{def:imed} its divergence term is $\Dv_a(t)=0$, so its index is
$\Ix_a(t)=\log \Np_a(t)\le\log t\le\log T$, and $\Ix_i(t)=\Is(t)\le\log T$. We show that
$s\Kinf(\hnu_{i,s},\mu)\le\log T$. If $i$ is an \term{w:leader}{empirical leader}, then $\hm_{i,s}=\hms(t)>\mu$, so
$\Kinf(\hnu_{i,s},\mu)=0$ by Lemma~\ref{lem:dualfacts}(b) (\emph{zero at or above the mean}). Otherwise
$\Dv_i(t)=\Kinf(\hnu_{i,s},\hms(t))\ge\Kinf(\hnu_{i,s},\mu)$ by Lemma~\ref{lem:dualfacts}(a)
(\emph{monotonicity in the level}), and, since $\log s\ge0$,
\begin{equation*}
  s\Kinf(\hnu_{i,s},\mu)\le s\Dv_i(t)\le \Ix_i(t)\le\log T.
\end{equation*}

\emph{Step 2: the first inequality.} By Step 1, and since $\hnut_i(t)=\hnu_{i,\Np_i(t)}$ with $\Np_i(t)\ge1$
for $t\ge \Karms$, we get for every round $t\in\{\Karms,\dots,T-1\}$
\begin{align*}
  \ind\{\Ap_t=i,\ \hms(t)>\mu\}
  &\le\ind\big\{\Ap_t=i,\ \Np_i(t)\Kinf(\hnu_{i,\Np_i(t)},\mu)\le\log T\big\}\\
  &=\sum_{s\ge1}\ind\{\Ap_t=i,\ \Np_i(t)=s\}\,\ind\big\{s\Kinf(\hnu_{i,s},\mu)\le\log T\big\}.
\end{align*}
Summing over $t$ and using the fundamental inequality~\eqref{eq:fundamental},
\begin{align*}
  \TB&\le\sum_{s\ge1}\ind\big\{s\Kinf(\hnu_{i,s},\mu)\le\log T\big\}\sum_{t=\Karms}^{T-1}\ind\{\Ap_t=i,\ \Np_i(t)=s\}\\
  &\le\sum_{s\ge1}\ind\big\{s\Kinf(\hnu_{i,s},\mu)\le\log T\big\}.
\end{align*}

\emph{Step 3: move the level to $\ms$.} By Lemma~\ref{lem:shift} (\emph{shifting the level}) with the
levels $\mu<\ms<1$, for every $s\ge1$,
\begin{equation*}
  \Kinf(\hnu_{i,s},\ms)\le\Kinf(\hnu_{i,s},\mu)+\frac{\delta}{1-\ms}.
\end{equation*}
Since $\Ki_i-2\delta/(1-\ms)>0$, an integer $s\ge \MT_T$ satisfies $s(\Ki_i-2\delta/(1-\ms))\ge\log T$, that
is $\log T/s\le\Ki_i-2\delta/(1-\ms)$. Then we split according to the event $\{s<\MT_T\}$:
\begin{align*}
  \ind\big\{s\Kinf(\hnu_{i,s},\mu)\le\log T\big\}
  &\le\ind\{s<\MT_T\}+\ind\Big\{s\ge \MT_T,\ \Kinf(\hnu_{i,s},\mu)\le\frac{\log T}{s}\Big\}\\
  &\le\ind\{s<\MT_T\}+\ind\Big\{s\ge \MT_T,\ \Kinf(\hnu_{i,s},\ms)\le\frac{\log T}{s}+\frac{\delta}{1-\ms}\Big\}\\
  &\le\ind\{s<\MT_T\}+\ind\Big\{\Kinf(\hnu_{i,s},\ms)\le\Ki_i-\frac{2\delta}{1-\ms}+\frac{\delta}{1-\ms}\Big\}\\
  &=\ind\{s<\MT_T\}+\ind\big\{\Kinf(\hnu_{i,s},\ms)\le \ui_i\big\},
\end{align*}
by the definition~\eqref{eq:ui} of $\ui_i$.

\emph{Step 4: conclusion.} There are at most $\MT_T$ integers $s\ge1$ with $s<\MT_T$. Summing the last
display over $s\ge1$ and using Step 2 gives $\TB\le \MT_T+\TC$.
\end{proof}

\subsection{Proof of Lemma~\ref{lem:termC}}\label{sec:proof-lem-termC}

\begin{proof}
\proves{lem:termC}
\uses{def:model, def:kinf, def:run, def:rate, lem:termB, lem:optimizer, lem:subopt}
Since $\marm_i<\ms<1$, Lemma~\ref{lem:optimizer}(a) gives $0<\Ki_i<\infty$, and $\ui_i<\Ki_i$ because
$\delta>0$. The rewards $\Xs_{i,1},\Xs_{i,2},\dots$ are independent of law $\nu_i$, and $\hnu_{i,s}$ is their
empirical law after $s$ samples (Definition~\ref{def:run}). Lemma~\ref{lem:subopt} (\emph{deviations of
a \term{w:subopt}{suboptimal} arm need no moment}) with $P=\nu_i$, the level $\ms\in(\marm_i,1)$ and $u=\ui_i<\Ki_i$ gives
$\Lt(\nu_i,\ms,\ui_i)>0$, that is $\Lt_i(\delta)>0$ by Definition~\ref{def:rate}, and for every $s\ge1$
\begin{equation*}
  \Prob\big(\Kinf(\hnu_{i,s},\ms)\le \ui_i\big)\le e^{-s\Lt_i(\delta)}.
\end{equation*}
Finally, by monotone convergence,
\begin{equation*}
  \E\TC=\sum_{s\ge1}\Prob\big(\Kinf(\hnu_{i,s},\ms)\le \ui_i\big)\le\sum_{s\ge1}e^{-s\Lt_i(\delta)}
  =\frac{e^{-\Lt_i(\delta)}}{1-e^{-\Lt_i(\delta)}}\le\frac{1}{1-e^{-\Lt_i(\delta)}}.
\end{equation*}
\end{proof}

% ----------------------------------------------------------------
\section{Proof of Proposition~\ref{prop:key}}\label{sec:proof-key}

In this appendix we prove that $\Gam(P,\mu)<\infty$ for $P\in\Sm_2$ and $\mu<\mean(P)$, with the explicit
bound~\eqref{eq:gammabound}.

\paragraph{Skeleton} To prove Proposition~\ref{prop:key} we use the decomposition \eqref{eq:split}
below. For each sample size $s$, the indices $l\le s$ are split into big jumps
($Y_l>\beta s$) and a good part, and $\Omega$ into five events according to the number $K_s$ of big
jumps and the sum $G_s$ of the good part. The supremum over $s$ is bounded by the supremum of the
one-jump terms plus the sum over $s$ of the other terms:

\begin{center}
\small
\begin{tabular}{@{}lp{0.3\textwidth}p{0.25\textwidth}p{0.3\textwidth}@{}}
\toprule
Event & Meaning & Bound on its contribution & Tool \\
\midrule
$E5_s$ & no big jump, good part $\ge s\delta/2$ & $0$ & the empirical mean exceeds $\mu$ \\
$E4_s$ & one big jump, good part $\ge s\delta/2$ & $2w/(\beta\delta)$ for the supremum & Lemma~\ref{lem:goodpart}; second moment \\
$E3_s$ & $2$ to $k_0-1$ big jumps, good part $\ge s\delta/2$ & $C_2w^2/\beta^2$ for the sum & Lemma~\ref{lem:goodpart}; independence \\
$E1_s$ & at least $k_0$ big jumps & polynomial in $s_3$ & Chebyshev, Lemma~\ref{lem:moments} \\
$E2_s$ & fewer than $k_0$ big jumps, good part $<s\delta/2$ & polynomial in $s_1$ & Lemma~\ref{lem:bennett}, Lemma~\ref{lem:moments} \\
\bottomrule
\end{tabular}
\end{center}

\paragraph{Lemmas, by the question they answer} \emph{How large can the dual statistic be?}
Lemma~\ref{lem:grid} reduces the maximum over $\lambda$ to a grid, and Lemma~\ref{lem:moments}
bounds its moments of order $\gamma\ge1$. \emph{How large is it when the good part is large?}
Lemma~\ref{lem:goodpart}. \emph{How likely is a small good part?} Bennett's inequality,
Lemma~\ref{lem:bennett}.

\begin{lemma}[Bennett's inequality, in the form used]\label{lem:bennett}
Let $\xi_1,\dots,\xi_s$ be independent centred random variables with $\xi_l\le b$ almost surely for
some $b>0$, and $\sum_l\E\xi_l^2\le V$ for some $V>0$. Then, for every $x>0$,
$\Prob\big(\sum_l\xi_l\ge x\big)\le\big(eV/(bx)\big)^{x/b}$.
\end{lemma}

\emph{Role.} The bound on $\Prob(E2_s)$ in Step 6 of the proof of Proposition~\ref{prop:key}.

\begin{proof}
\emph{Setting.} Let $\varphi(v)\coloneqq(e^v-1-v)/v^2$ for $v\ne0$ and $\varphi(0)\coloneqq1/2$.
\emph{Goal.} A Chernoff bound with the variance.

By Taylor's formula with integral remainder, $\varphi(v)=\int_0^1(1-\theta)e^{v\theta}\,d\theta$,
which is nondecreasing in $v$. For $\eta>0$ and $\xi\le b$, this gives
$e^{\eta\xi}=1+\eta\xi+\eta^2\xi^2\varphi(\eta\xi)\le1+\eta\xi+\eta^2\xi^2\varphi(\eta b)$. Taking
expectations, with $\E\xi_l=0$ and $1+v\le e^v$,
\[
  \E e^{\eta\xi_l}\le1+\eta^2\E\xi_l^2\,\varphi(\eta b)\le\exp\Big(\frac{\E\xi_l^2}{b^2}\big(e^{\eta b}-1-\eta b\big)\Big).
\]
By independence and Markov's inequality,
\[
  \Prob\Big(\sum_l\xi_l\ge x\Big)\le\exp\Big(\frac V{b^2}\big(e^{\eta b}-1-\eta b\big)-\eta x\Big).
\]
Choose $\eta b\coloneqq\log(1+bx/V)$. The exponent becomes
$\frac xb-\big(\frac V{b^2}+\frac xb\big)\log\big(1+\frac{bx}V\big)$, which is at most
$\frac xb-\frac xb\log\frac{bx}{V}=\frac xb\log\frac{eV}{bx}$, because $\log(1+v)\ge\log v$ and
$\frac V{b^2}\log(1+\frac{bx}{V})\ge0$. This is the claim.
\end{proof}

\begin{lemma}[A grid for the dual variable]\label{lem:grid}
Let $c>0$, $s\ge1$, $z_1,\dots,z_s\le c$, $Q(\lambda)\coloneqq\prod_{l=1}^s(1-\lambda z_l)$ for
$\lambda\in[0,1/c]$, and $\lambda_k\coloneqq k/(2cs)$ for $k=0,\dots,2s$. Then
$\max_{\lambda\in[0,1/c]}Q(\lambda)\le e\max_{0\le k\le2s}Q(\lambda_k)$.
\end{lemma}

\emph{Role.} Lemma~\ref{lem:moments}(a); it is the grid of the proof of \citet[Proposition~11]{HT15}
with mesh $1/(2cs)$.

\begin{proof}
\emph{Setting.} Let $\ell(\lambda)\coloneqq s^{-1}\log Q(\lambda)$, with values in
$[-\infty,\infty)$ and $\log0\coloneqq-\infty$. \emph{Goal.} $s\ell(\lambda)\le\max_ks\ell(\lambda_k)+1$ for every
$\lambda\in[0,1/c]$.

\emph{Step 1: concavity and slope.} $\ell$ is concave on $[0,1/c]$ (a mean of concave functions
with values in $[-\infty,\infty)$), finite and differentiable on $[0,1/c)$, where every factor is at least
$1-\lambda c>0$, with $\ell'(\lambda)=-\frac1s\sum_lz_l/(1-\lambda z_l)$. At $\lambda_s=1/(2c)$, each
term $-z/(1-z/(2c))=-2cz/(2c-z)$ lies in $[-2c,2c)$ for $z\le c$: it is below $2c$ if $z<0$ and at
least $-2c\cdot c/(2c-c)=-2c$ if $0\le z\le c$. So $-2c\le\ell'(\lambda_s)<2c$.

\emph{Step 2: $\lambda\le\lambda_s$.} Let $\lambda\in[\lambda_k,\lambda_{k+1}]$ with $k+1\le s$. Since
$\ell'$ is nonincreasing, $\ell'(\lambda_{k+1})\ge\ell'(\lambda_s)\ge-2c$, and the tangent at
$\lambda_{k+1}$ gives
\[
  \ell(\lambda)\le\ell(\lambda_{k+1})+\ell'(\lambda_{k+1})(\lambda-\lambda_{k+1})\le\ell(\lambda_{k+1})+2c\cdot\frac1{2cs}
  =\ell(\lambda_{k+1})+\frac1s.
\]

\emph{Step 3: $\lambda\ge\lambda_s$.} Let $\lambda\in[\lambda_{k-1},\lambda_k]$ with $s+1\le k\le2s$.
Then $\lambda_{k-1}<1/c$, so $\ell$ is differentiable there, $\ell'(\lambda_{k-1})\le\ell'(\lambda_s)<2c$,
and the tangent at $\lambda_{k-1}$ (valid on $[0,1/c]$ by concavity) gives
$\ell(\lambda)\le\ell(\lambda_{k-1})+2c\cdot\frac1{2cs}=\ell(\lambda_{k-1})+\frac1s$.

\emph{Step 4: conclusion.} The intervals of Steps 2 and 3 cover $[0,1/c]$, so
$s\ell(\lambda)\le\max_ks\ell(\lambda_k)+1$ for every $\lambda$; exponentiating gives the claim.
\end{proof}

\begin{lemma}[Moments of the dual statistic]\label{lem:moments}
\uses{def:model, def:kinf}
Let $P\in\Sm$ and $\mu<\mean(P)$; let $\hP_s$ be the empirical law of $X_1,\dots,X_s$, i.i.d.\ of law $P$,
and $\phi(\gamma)\coloneqq\E_P[((1-X)/(1-\mu))^\gamma]\in[0,\infty]$ for $\gamma\ge1$.
\begin{enumerate}[label=(\alph*)]
\item For every $s\ge1$ and $\gamma\ge1$,
  $\E\exp(\gamma s\Kinf(\hP_s,\mu))\le e^\gamma(2s+1)\max\{1,\phi(\gamma)\}^s$. In particular
  $\E\exp(s\Kinf(\hP_s,\mu))\le e(2s+1)$.
\item If $\E_P|X|^{\gamma_0}<\infty$ for some $\gamma_0>1$, there is $\gamma\in(1,\gamma_0]$ with
  $\phi(\gamma)\le1$. For such $\gamma$, every $s\ge1$ and every event $E$,
  \[
    \E\big[s\exp(s\Kinf(\hP_s,\mu))\ind_E\big]\le e\,s\,(2s+1)^{1/\gamma}\,\Prob(E)^{1-1/\gamma}.
  \]
\end{enumerate}
\end{lemma}

\emph{Role.} The crude bounds for small $s$ and the H\"older bounds of Steps 5 and 6 of the proof of
Proposition~\ref{prop:key}. Part (b) is the inequality $\lambda_\mu>1$ of
\citet[Proposition~11(i)]{HT15}, which only needs a moment of order above $1$.

\begin{proof}
\uses{def:model, def:kinf, lem:dualfacts, lem:grid}
\emph{Setting.} Let $Z_l\coloneqq X_l-\mu\le c\coloneqq1-\mu$ and $\lmax\coloneqq1/c$.
\emph{Goal.} (a) and (b).

\emph{Step 1: (a).} By Lemma~\ref{lem:dualfacts}(c) and Lemma~\ref{lem:grid},
\[
  \exp(\gamma s\Kinf(\hP_s,\mu))\le e^\gamma\max_{0\le k\le2s}\prod_l(1-\lambda_kZ_l)^\gamma\le
  e^\gamma\sum_{k=0}^{2s}\prod_l(1-\lambda_kZ_l)^\gamma.
\]
By independence, $\E\prod_l(1-\lambda Z_l)^\gamma=f(\lambda)^s$ with
$f(\lambda)\coloneqq\E(1-\lambda Z_1)^\gamma$. For $\gamma\ge1$, $\lambda\mapsto(1-\lambda z)^\gamma$ is
convex on $[0,\lmax]$ (an increasing convex function of a nonnegative affine function), so $f$ is
convex and $f(\lambda)\le\max\{f(0),f(\lmax)\}=\max\{1,\phi(\gamma)\}$, since
$1-\lmax Z=(1-X)/(1-\mu)$. Summing over the $2s+1$ grid points gives (a). For $\gamma=1$,
$\phi(1)=(1-\mean(P))/(1-\mu)<1$.

\emph{Step 2: the exponent of (b).} For $u\ge0$ and $\gamma\in[1,\gamma_0]$,
$u^\gamma\le1+u^{\gamma_0}$; with $u=1-X\le1+|X|$, the integrand of $\phi$ is dominated on
$[1,\gamma_0]$ by an integrable function. By dominated convergence $\phi$ is finite and continuous on
$[1,\gamma_0]$. Since $\phi(1)<1$, there is $\gamma\in(1,\gamma_0]$ with $\phi(\gamma)\le1$.

\emph{Step 3: the H\"older bound.} For such $\gamma$, (a) gives $\E\exp(\gamma s\Kinf(\hP_s,\mu))\le e^\gamma(2s+1)$.
By H\"older's inequality with the exponents $\gamma$ and $\gamma/(\gamma-1)$,
\[
  \E\big[s\exp(s\Kinf(\hP_s,\mu))\ind_E\big]\le s\big(e^\gamma(2s+1)\big)^{1/\gamma}\Prob(E)^{1-1/\gamma}
  =e\,s\,(2s+1)^{1/\gamma}\Prob(E)^{1-1/\gamma}.
\]
\end{proof}

\begin{lemma}[The dual statistic when the good part is large]\label{lem:goodpart}
Let $c>0$, $s\ge1$, $z_1,\dots,z_s\le c$, and let $B\subseteq\{1,\dots,s\}$ be nonempty with
$z_l<0$ for $l\in B$; let $k\coloneqq|B|$, and assume that $\sum_{l\notin B}z_l\ge h$ for some
$h>0$. Then
\[
  \max_{\lambda\in[0,1/c]}\prod_{l=1}^s(1-\lambda z_l)\le\prod_{l\in B}\max\Big\{1,\frac{k|z_l|}{h}\Big\}.
\]
\end{lemma}

\emph{Role.} Steps 3 and 4 of the proof of Proposition~\ref{prop:key}, with $h=s\delta/2$: one big
jump $y$ multiplies the statistic by at most $2y/(s\delta)$.

\begin{proof}
\emph{Step 1: one variable.} For $y>0$ and $g>0$,
$\sup_{\lambda\ge0}\big(\log(1+\lambda y)-\lambda g\big)\le\log\max\{1,y/g\}$. Indeed the derivative
$y/(1+\lambda y)-g$ is decreasing in $\lambda$. If $y\le g$, it is at most $0$ and the supremum is
the value $0$ at $\lambda=0$. If $y>g$, the maximum is at $\lambda=1/g-1/y$, with value
$\log(y/g)-1+g/y\le\log(y/g)$.

\emph{Step 2: $\lambda<1/c$.} Every factor with $l\notin B$ satisfies $1-\lambda z_l\ge1-\lambda c>0$
and $\log(1-\lambda z_l)\le-\lambda z_l$. Hence
\[
  \log\prod_{l=1}^s(1-\lambda z_l)\le\sum_{l\in B}\log(1+\lambda|z_l|)-\lambda h
  =\sum_{l\in B}\Big(\log(1+\lambda|z_l|)-\lambda\frac hk\Big)\le\sum_{l\in B}\log\max\Big\{1,\frac{k|z_l|}{h}\Big\},
\]
by Step 1 with $g=h/k$.

\emph{Step 3: $\lambda=1/c$.} The product is a polynomial in $\lambda$, hence continuous on
$[0,1/c]$, and the bound of Step 2 extends to $\lambda=1/c$.
\end{proof}

\begin{proof}[Proof of Proposition~\ref{prop:key}]
\proves{prop:key}
\uses{def:model, def:kinf, def:dev, lem:dualfacts, lem:bennett, lem:moments, lem:goodpart}
\emph{Idea.} Split each sample of size $s$ into big negative jumps and a good part. With one big
jump and a large good part, $\ws_s$ is at most a constant times that jump, and the supremum over $s$
of these terms has expectation at most a constant times the second moment. Every other event has a
probability small enough in $s$ to be summed after a H\"older bound.

\emph{Setting.} Fix $P\in\Sm_2$ and $\mu<m\coloneqq \mean(P)$. Let $X_1,X_2,\dots$ be i.i.d.\ of law $P$,
$Z_l\coloneqq X_l-\mu$, $Y_l\coloneqq Z_l^-$, $\delta\coloneqq m-\mu>0$, $c\coloneqq1-\mu>0$,
$v\coloneqq\E Z_1^2$, $w\coloneqq\E Y_1^2$, $y_1\coloneqq\E Y_1$, $S_s\coloneqq\sum_{l\le s}Z_l$, and
$W_s\coloneqq \ws_s(X;\mu)$. By Lemma~\ref{lem:dualfacts}(c),
\[
  W_s=s\max_{0\le\lambda\le1/c}\prod_{l\le s}(1-\lambda Z_l)\,\ind\{S_s\le0\}.
\]
\emph{Goal.} $\Gam(P,\mu)=\E\sup_{s\ge1}W_s\le$ the right-hand side of~\eqref{eq:gammabound}.

Note that $v\ge(\E Z_1)^2=\delta^2>0$.

\emph{Step 1: parameters.} Lemma~\ref{lem:moments}(b) with $\gamma_0=2$ gives $\gamma\in(1,2]$ with
$\E_P[((1-X)/c)^\gamma]\le1$. Let $\theta\coloneqq1-1/\gamma\in(0,1/2]$, $k_0\coloneqq\lceil4/\theta\rceil\ge8$,
$\beta\coloneqq\min\{1,\delta/(4k_0)\}$ and, for $s\ge1$, $e_s\coloneqq\E[Y_1\ind\{Y_1>\beta s\}]$. Then
\begin{equation}\label{eq:es}
  e_s\le y_1,\qquad e_s\le\frac{w}{\beta s},\qquad\sum_{s\ge1}e_s=\E\big[Y_1\,|\{s\ge1:\beta s<Y_1\}|\big]\le\frac{w}{\beta}.
\end{equation}

\emph{Step 2: big jumps, good part, and a partition.} For $s\ge1$, let
$B_s\coloneqq\{l\le s:Y_l>\beta s\}$ (the big jumps), $K_s\coloneqq|B_s|$ and
$G_s\coloneqq\sum_{l\le s,\,l\notin B_s}Z_l$ (the good part). Since $l\in B_s$ if and only if
$Z_l<-\beta s$, we have $G_s=\sum_{l\le s}Z_l\ind\{Z_l\ge-\beta s\}$. The following events partition
$\Omega$:
\begin{align*}
  E1_s&\coloneqq\{K_s\ge k_0\}, & E2_s&\coloneqq\{K_s<k_0,\ G_s<s\delta/2\},\\
  E3_s&\coloneqq\{2\le K_s<k_0,\ G_s\ge s\delta/2\}, & E4_s&\coloneqq\{K_s=1,\ G_s\ge s\delta/2\},\\
  E5_s&\coloneqq\{K_s=0,\ G_s\ge s\delta/2\}.
\end{align*}
On $E5_s$, $S_s=G_s>0$, so $W_s=0$. Since all terms are nonnegative,
\begin{equation}\label{eq:split}
  \sup_{s\ge1}W_s\le\sup_{s\ge1}W_s\ind_{E4_s}+\sum_{s\ge1}W_s\big(\ind_{E1_s}+\ind_{E2_s}+\ind_{E3_s}\big).
\end{equation}

\emph{Step 3: one big jump, where the second moment is used.} On $E4_s$, let $B_s=\{l^\star\}$. If
$W_s>0$, then $S_s=G_s-Y_{l^\star}\le0$, so $Y_{l^\star}\ge G_s\ge s\delta/2$. Lemma~\ref{lem:goodpart}
with $B=\{l^\star\}$ and $h=s\delta/2$ then gives
\[
  W_s\le s\max\Big\{1,\frac{2Y_{l^\star}}{s\delta}\Big\}=\frac{2Y_{l^\star}}{\delta}.
\]
Moreover $Y_{l^\star}>\beta s\ge\beta l^\star$, since $l^\star\le s$. Hence, for every $s$,
$W_s\ind_{E4_s}\le\frac2\delta\sum_{l\ge1}Y_l\ind\{Y_l>\beta l\}$, and by~\eqref{eq:es}
\[
  \E\sup_{s\ge1}W_s\ind_{E4_s}\le\frac2\delta\sum_{l\ge1}e_l\le\frac{2w}{\beta\delta}.
\]

\emph{Step 4: two to $k_0-1$ big jumps.} Let
\[
  C_1\coloneqq\frac1\beta+\frac{2k_0}{\delta}\qquad\text{and}\qquad
  C_2\coloneqq\frac{C_1^2}2+\sum_{k=3}^{k_0-1}\frac{C_1^ky_1^{k-2}}{k!}.
\]
On $E3_s$ with $K_s=k$, each $l\in B_s$ has $Y_l>\beta s$, so
$\max\{1,2kY_l/(s\delta)\}\le C_1Y_l/s$. Lemma~\ref{lem:goodpart} with $h=s\delta/2$ gives
$W_s\le s\prod_{l\in B_s}(C_1Y_l/s)$ on $E3_s$. Bounding the indicator of $\{B_s=B\}$ by
$\prod_{l\in B}\ind\{Y_l>\beta s\}$ and summing over the subsets $B$ of size $k$, independence gives
\[
  \E[W_s\ind_{E3_s}]\le s\sum_{k=2}^{k_0-1}\binom sk\Big(\frac{C_1e_s}{s}\Big)^k
  \le s\sum_{k=2}^{k_0-1}\frac{C_1^ke_s^k}{k!}\le C_2\,s\,e_s^2,
\]
using $\binom sk\le s^k/k!$ and $e_s\le y_1$. By~\eqref{eq:es}, $s\,e_s^2\le(w/\beta)e_s$, so
$\sum_{s\ge1}\E[W_s\ind_{E3_s}]\le C_2w^2/\beta^2$.

\emph{Step 5: many big jumps.} Let $s_3\coloneqq w/\beta^2$. By Chebyshev's inequality,
$\Prob(Y_1>\beta s)\le w/(\beta s)^2$, so
\[
  \Prob(E1_s)\le\binom s{k_0}\Prob(Y_1>\beta s)^{k_0}\le\Big(\frac{s_3}{s}\Big)^{k_0}.
\]
For every $s$, Lemma~\ref{lem:moments}(a) gives $\E[W_s\ind_{E1_s}]\le\E W_s\le e\,s(2s+1)$. For
$s\ge s_3$, Lemma~\ref{lem:moments}(b), $k_0\theta\ge4$ and $(2s+1)^{1/\gamma}\le3s$ give
\[
  \E[W_s\ind_{E1_s}]\le e\,s(2s+1)^{1/\gamma}\Big(\frac{s_3}{s}\Big)^{k_0\theta}\le3e\,s^2\Big(\frac{s_3}{s}\Big)^4=\frac{3e\,s_3^4}{s^2}.
\]
Since $\sum_{s\ge\max\{1,s_3\}}s^{-2}\le2/\max\{1,s_3\}$, the sum over $s\ge s_3$ is at most $6e\,s_3^3$
(if $s_3<1$, use $s_3^4\le s_3^3$). Hence
$\sum_s\E[W_s\ind_{E1_s}]\le e\sum_{1\le s<s_3}s(2s+1)+6e\,s_3^3$.

\emph{Step 6: a small good part.} Write $G_s=\sum_{l\le s}Z_l^{(s)}$ with
$Z_l^{(s)}\coloneqq Z_l\ind\{Z_l\ge-\beta s\}$; these are i.i.d.\ in $l$, with
$\E Z_1^{(s)}=\delta+e_s\ge\delta$ and $\mathrm{Var}\,Z_1^{(s)}\le\E Z_1^2=v$. Let
$\xi_l\coloneqq\E Z_1^{(s)}-Z_l^{(s)}$; they are centred, and $\xi_l\le c+\beta s\eqqcolon b_s$
because $Z_l^{(s)}\ge-\beta s$ and $\E Z_1^{(s)}\le\E Z_1^+\le c$. On $\{G_s<s\delta/2\}$,
$\sum_l\xi_l=s\E Z_1^{(s)}-G_s>s\delta/2$. Lemma~\ref{lem:bennett} with $V=sv$ and $x=s\delta/2$ gives
\[
  \Prob(E2_s)\le\Prob\big(G_s<s\delta/2\big)\le\Big(\frac{2ev}{\delta b_s}\Big)^{s\delta/(2b_s)}.
\]
Let $s_1\coloneqq\max\{c/\beta,\ 2ev/(\delta\beta)\}$. For $s\ge s_1$, $\beta s\le b_s\le2\beta s$, so the
base is at most $2ev/(\delta\beta s)\le s_1/s\le1$ and the exponent is at least
$\delta/(4\beta)\ge k_0$. Hence $\Prob(E2_s)\le(s_1/s)^{k_0}$, and, as in Step 5,
$\E[W_s\ind_{E2_s}]\le3e\,s_1^4/s^2$ for $s\ge s_1$ and $\E[W_s\ind_{E2_s}]\le e\,s(2s+1)$ for every
$s$. So $\sum_s\E[W_s\ind_{E2_s}]\le e\sum_{1\le s<s_1}s(2s+1)+6e\,s_1^3$.

\emph{Step 7: conclusion.} Taking expectations in~\eqref{eq:split} and adding Steps 3--6,
\begin{equation}\label{eq:gammabound}
  \Gam(P,\mu)\le\frac{2w}{\beta\delta}+\frac{C_2w^2}{\beta^2}+e\sum_{1\le s<s_1}s(2s+1)+6e\,s_1^3
  +e\sum_{1\le s<s_3}s(2s+1)+6e\,s_3^3<\infty,
\end{equation}
with $\delta=\mean(P)-\mu$, $c=1-\mu$, $v=\E_P(X-\mu)^2$, $w=\E_P((X-\mu)^-)^2$, $y_1=\E_P(X-\mu)^-$, the
exponent $\gamma$ of Step 1, $k_0=\lceil4\gamma/(\gamma-1)\rceil$, $\beta=\min\{1,\delta/(4k_0)\}$,
$C_1$ and $C_2$ of Step 4, $s_1=\max\{c/\beta,2ev/(\delta\beta)\}$ and $s_3=w/\beta^2$.
\end{proof}

% ----------------------------------------------------------------
\section{Proof of Lemma~\ref{lem:three-moments}}\label{sec:proof-three}

\begin{proof}[Proof of Lemma~\ref{lem:three-moments}]
\proves{lem:three-moments}
\uses{def:model, def:kinf, def:dev}
\emph{Idea.} One summand below $-2s\delta$ among the first $s$, with the others near their mean,
puts the empirical mean below $\mu$; so $\Prob(\bX_s\le\mu)$ is at least of order
$s\,\Prob(X_1^-\gtrsim s)$, and $\sum_s s\cdot s\,\Prob(X_1^-\gtrsim s)$ diverges without a third moment.

The first inequality holds because $\Kinf\ge0$. Let $\delta\coloneqq \mean(P)-\mu$ and
$Z_l\coloneqq X_l-\mu$. Fix $s\ge2$ and let $p_s\coloneqq\Prob(Z_1\le-2s\delta)$. For $l\le s$, the event
$A_l\coloneqq\{Z_l\le-2s\delta\}\cap\{\sum_{r\le s,r\ne l}Z_r\le2s\delta\}$ is contained in
$\{\sum_{r\le s}Z_r\le0\}=\{\bX_s\le\mu\}$. Let $q_s\coloneqq\Prob(\sum_{r=2}^sZ_r\le2s\delta)$; it tends
to $1$ by the weak law of large numbers, since the sum has mean $(s-1)\delta$. By the
Bonferroni inequality and independence,
\[
  \Prob(\bX_s\le\mu)\ge\sum_{l\le s}\Prob(A_l)-\sum_{l<l'\le s}\Prob(A_l\cap A_{l'})
  \ge sp_sq_s-\frac{s^2p_s^2}{2}.
\]
Since $\E|Z_1|<\infty$, $sp_s\le\E[|Z_1|\ind\{|Z_1|\ge2s\delta\}]/(2\delta)\to0$. Hence
$\Prob(\bX_s\le\mu)\ge sp_s/4$ for every large $s$. Finally
$\sum_ss^2p_s=\sum_ss^2\Prob(Z_1^-\ge2s\delta)=\infty$ by comparison with
$\int_0^\infty x^2\Prob(Z_1^-\ge2\delta(x+1))\,dx$, because $\E[(Z_1^-)^3]=\infty$.
\end{proof}

\section{Proof of Corollary~\ref{cor:optimal}}\label{sec:proof-optimal}

In this appendix we deduce Corollary~\ref{cor:optimal} from Theorem~\ref{thm:main}, the lower bound of
Theorem~\ref{cite:GMS19-thm1} and Lemma~\ref{lem:class} (\emph{the constant does not depend on the moment
class}).

\begin{proof}[Proof of Corollary~\ref{cor:optimal}]
\proves{cor:optimal}
\uses{def:model, def:kinf, def:run, def:classes, def:imed, thm:main, cite:GMS19-thm1, lem:class}
(a) Fix $\delta$ as in Theorem~\ref{thm:main}. Dividing its bound by $\log T$ and letting
$T\to\infty$ gives $\limsup_T\E \Np_i(T)/\log T\le1/(\Ki_i-2\delta/(1-\ms))$; letting $\delta\to0$
gives (a).

(b) Let $\nu\in\Mm_r^\arms$. No arm has the law $\delta_1$, so $\ms<1$ (Section~\ref{sec:setting}).
Every \term{w:optimal}{optimal} arm $j$ has $\nu_j\in\Sm_r\subseteq\Sm_2$, since $|x|^2\le1+|x|^r$ for $r\ge2$. With a
fixed $\delta$, Theorem~\ref{thm:main} gives $\E \Np_i(T)=O(\log T)$ for every \term{w:subopt}{suboptimal} arm $i$, so
$\Reg_T=O(\log T)$: \imed{} is \term{w:consistent}{consistent} and \term{w:ufc}{uniformly fast convergent} on $\Mm_r$.
Theorem~\ref{cite:GMS19-thm1} with $\mathcal D=\Mm_r$ gives, for every \term{w:learner}{learner} uniformly fast
convergent on $\Mm_r$, $\liminf_T\E \Np_i(T)/\log T\ge1/\Kinf(\nu_i,\ms;\Mm_r)$. By
Lemma~\ref{lem:class} (with $P=\nu_i\in\Mm_r$ and $\ms<1$),
$\Kinf(\nu_i,\ms;\Mm_r)=\Kinf(\nu_i,\ms)$. Together with (a), the limit of $\E \Np_i(T)/\log T$
exists for \imed{} and equals $1/\Kinf(\nu_i,\ms)$. Multiplying by $\gap_i$ and summing over the
\term{w:subopt}{suboptimal} arms gives the statement on $\Reg_T$.
\end{proof}

% ----------------------------------------------------------------
\section{Proof of Theorem~\ref{thm:counter} and Corollary~\ref{cor:pareto}}\label{sec:proofs-counter}

In this appendix we prove the lower bounds of Section~\ref{sec:counter}. The proof of
Theorem~\ref{thm:counter} uses three lemmas: the index of the \term{w:optimal}{optimal} arm after a negative jump
(Lemma~\ref{lem:jump}), \imed{} explores a point mass logarithmically often
(Lemma~\ref{lem:explore}), and the maximum of nonnegative numbers is at least their sum minus the sum
of their pairwise minima (Lemma~\ref{lem:pairs}). Its skeleton: the frozen-index inequality
$\Np_2(T)\ge\max_sV_s$ (Step 2), a dichotomy on how often arm $1$ is pulled in the first half of the
rounds (Step 3), then part (i) by Bonferroni's inequality (Step 4) and part (ii) by
Lemmas~\ref{lem:explore} and~\ref{lem:pairs} (Step 5).

\begin{lemma}[The index after a negative jump]\label{lem:jump}
\uses{def:kinf}
Let $s\ge1$, let $x_1,\dots,x_{s-1}\le1$ and $y\ge0$, and let $\hat Q$ be the empirical law of
$(x_1,\dots,x_{s-1},-y)$. Then $s\Kinf(\hat Q,0)+\log s\ge\log y-1$, with $\log0\coloneqq-\infty$.
\end{lemma}

\emph{Role.} The frozen index in Step 2 of the proof of Theorem~\ref{thm:counter}.

\begin{proof}
\uses{def:kinf, lem:dualfacts}
\emph{Setting.} Fix $s$, $x_{1:s-1}$ and $y>0$ (the case $y=0$ is trivial). \emph{Goal.} A lower
bound on $s\Kinf(\hat Q,0)$, from one value of the dual variable in~\eqref{eq:dual}, whose range
is $[0,1]$ at the level $0$.

\emph{Step 1: $s=1$.} By Lemma~\ref{lem:dualfacts}(c),
$\Kinf(\delta_{-y},0)=\max_{0\le\lambda\le1}\log(1+\lambda y)=\log(1+y)\ge\log y$.

\emph{Step 2: $s\ge2$.} Take $\lambda\coloneqq1/s\le1/2$ in Lemma~\ref{lem:dualfacts}(c). For
$z\in[0,1/2]$, $\log(1-z)\ge-z-z^2$, because the derivative of $\log(1-z)+z+z^2$ is
$z(1-2z)/(1-z)\ge0$. Hence each $x_r\in[0,1]$ gives
$\log(1-x_r/s)\ge-1/s-1/s^2$, and each $x_r<0$ gives $\log(1-x_r/s)>0$. Therefore
\begin{align*}
  s\Kinf(\hat Q,0)&\ge\sum_{r=1}^{s-1}\log\Big(1-\frac{x_r}{s}\Big)+\log\Big(1+\frac ys\Big)
    &&\text{(the dual at $\lambda=1/s$)}\\
  &\ge-(s-1)\Big(\frac1s+\frac1{s^2}\Big)+\log\Big(1+\frac ys\Big)&&\text{(the two cases above)}\\
  &=-\Big(1-\frac1{s^2}\Big)+\log\Big(1+\frac ys\Big)\ge\log\Big(1+\frac ys\Big)-1.
\end{align*}
Adding $\log s$ and using $\log s+\log(1+y/s)=\log(s+y)\ge\log y$ gives the claim.
\end{proof}

\begin{lemma}[A point mass is explored logarithmically often]\label{lem:explore}
\uses{def:model, def:kinf, def:run, def:imed}
In the setting of Theorem~\ref{thm:counter}, with any $P\in\Sm$ such that $\mean(P)\in(0,1)$ and no
assumption on its tail, almost surely
\[
  \liminf_{T\to\infty}\frac{\Np_2(T)}{\log T}\ge\frac{1}{\log(1/(1-\mean(P)))}.
\]
\end{lemma}

\emph{Role.} The term $1/\log\frac1{1-\mean(P)}$ of Theorem~\ref{thm:counter}(ii).

\begin{proof}
\uses{def:model, def:kinf, def:run, def:imed, lem:dualfacts}
\emph{Setting.} Let $d(x)\coloneqq\Kinf(\delta_0,x)$ for $x\ge0$; by Lemma~\ref{lem:dualfacts}(b)
and (d), $d(x)=\log\frac1{1-x}$ on $[0,1)$ and $d(x)=+\infty$ on $[1,\infty)$, so $d$ is
nondecreasing, and continuous on $[0,1)$. Let $d^\star\coloneqq d(\mean(P))\in(0,\infty)$.
\emph{Goal.} $\liminf_T\Np_2(T)/\log T\ge1/d^\star$ almost surely.

\emph{Step 1: the index of arm $2$.} At a round $t\ge2$, $\hms(t)=\max\{\hmt_1(t),0\}$. If
$\hmt_1(t)\le0$, arm $2$ is an \term{w:leader}{empirical leader} and $\Dv_2(t)=0=d(0)=d(\hms(t))$. Otherwise
$\Dv_2(t)=\Kinf(\delta_0,\hms(t))=d(\hms(t))$. So in all cases
$\Ix_2(t)=\Np_2(t)\,d(\hms(t))+\log \Np_2(t)$.

\emph{Step 2: the last pull of arm $1$.} Fix $T\ge4$. If arm $1$ is pulled at no round of
$\{2,\dots,T-1\}$, then $\Np_2(T)=T-1$. Otherwise let $\ell$ be the last such round. Then
$\Np_1(\ell)=\Np_1(T)-1\ge1$, $\Np_2(\ell)\le \Np_2(T)$, and $\Ix_1(\ell)\le \Ix_2(\ell)$ since $\Ap_\ell=1$. Since
$\Ix_1(\ell)\ge\log \Np_1(\ell)$, and since $n\mapsto nd+\log n$ is nondecreasing for $d\ge0$,
\[
  \log(\Np_1(T)-1)\le \Np_2(T)\,d(\hms(\ell))+\log \Np_2(T).
\]

\emph{Step 3: when $\Np_2(T)\le T/2$.} Then $\Np_1(T)\ge T/2$, so $\Np_1(\ell)\ge T/2-1$ and
$\hms(\ell)\le\bar m_T\coloneqq\sup_{s\ge T/2-1}\max\{\hm_{1,s},0\}$. As $d$ is nondecreasing,
\[
  \log(T/2-1)\le \Np_2(T)\,d(\bar m_T)+\log \Np_2(T).
\]

\emph{Step 4: conclusion.} By the strong law of large numbers, $\hm_{1,s}\to \mean(P)$ almost surely, so
$\bar m_T\to \mean(P)\in(0,1)$; hence $\bar m_T\in(0,1)$ for large $T$ (in particular $d(\bar m_T)$ is
finite, which excludes the case where every sample of arm $1$ equals $1$), and
$d(\bar m_T)\to d^\star$. For such $T$, Steps 2 and 3 give: either $\Np_2(T)\ge\min\{T/2,(\log T)^2\}$,
or $\log \Np_2(T)\le2\log\log T$ and then
\[
  \Np_2(T)\ge\frac{\log(T/2-1)-2\log\log T}{d(\bar m_T)}.
\]
Dividing by $\log T$ and letting $T\to\infty$ gives the claim.
\end{proof}

\begin{lemma}[Maximum from pairs]\label{lem:pairs}
For $n\ge1$ and nonnegative reals $u_1,\dots,u_n$,
$\max_ku_k\ge\sum_ku_k-\sum_{k<k'}\min\{u_k,u_{k'}\}$.
\end{lemma}

\emph{Role.} Step 5 of the proof of Theorem~\ref{thm:counter}.

\begin{proof}
Sort the numbers as $u_{(1)}\ge u_{(2)}\ge\dots\ge u_{(n)}$. The pairs whose smaller element is
$u_{(k)}$ are $k-1$ in number, so
$\sum_{k<k'}\min\{u_k,u_{k'}\}=\sum_{k=1}^n(k-1)u_{(k)}\ge\sum_{k\ge2}u_{(k)}=\sum_ku_k-\max_ku_k$.
\end{proof}

\begin{proof}[Proof of Theorem~\ref{thm:counter}]
\proves{thm:counter}
\uses{def:model, def:kinf, def:run, def:imed, lem:dualfacts, lem:jump, lem:explore, lem:pairs}
\emph{Idea.} After a large negative reward of arm $1$, its index is frozen at a value of at least
$\log(Y_s/e)$ while arm $2$, the \term{w:leader}{empirical leader}, is pulled. The reward is independent of the
time at which it is drawn, so the expected number of pulls of arm $2$ is bounded below by tail
integrals of $Y$.

\emph{Setting.} Fix $P$, a \term{w:tie}{tie rule} and $T\ge8$. Let $h\coloneqq\lfloor T/2\rfloor$ and
$g\coloneqq\lfloor T/4\rfloor\ge T/8$. For $s\ge1$, let $\tau_s$ be the round of the $s$-th pull of
arm $1$ ($\tau_s\coloneqq\infty$ if there is none) and $Y_s\coloneqq-\Xs_{1,s}$. \emph{Goal.} (i) and (ii).

\emph{Step 1: independence.} For every round $u$, the event $\{\tau_s\le u\}$ is a function of $\Useed$
and of $\Xs_{1,1},\dots,\Xs_{1,s-1}$. Indeed, replace $\Xs_{1,s'}$ by $0$ for every $s'\ge s$: the run of
\imed{} with the same $\Useed$ is unchanged up to and including round $\tau_s$, because before that round
the \term{w:learner}{learner} has observed only $\Xs_{1,1},\dots,\Xs_{1,s-1}$ from arm $1$ and the deterministic rewards
of arm $2$. Hence $\{\tau_s\le u\}$ is independent of $Y_s$.

\emph{Step 2: frozen index.} We show that, on $\{\tau_s<h,\ Y_s\ge s-1\}$,
\begin{equation}\label{eq:frozen}
  \Np_2(T)\ge\min\{Y_s/e,\ T/2\}.
\end{equation}
Let $t>\tau_s$ be a round such that arm $1$ has not been pulled in the rounds $\tau_s+1,\dots,t-1$.
If $t=1$, then $\tau_s=0$ and round $1$ pulls arm $2$ by initialisation. If $t\ge2$, then
$\Np_1(t)=s$ and
\[
  \hmt_1(t)=\hm_{1,s}=\frac1s\Big(\sum_{r<s}\Xs_{1,r}-Y_s\Big)\le\frac{s-1-Y_s}{s}\le0=\hmt_2(t),
\]
so $\hms(t)=0$ and arm $2$ is an \term{w:leader}{empirical leader} with $\Ix_2(t)=\log \Np_2(t)$. The empirical law
of arm $1$ is that of $(\Xs_{1,1},\dots,\Xs_{1,s-1},-Y_s)$ with $Y_s\ge0$. By
Lemma~\ref{lem:dualfacts}(b), $\Dv_1(t)=\Kinf(\hnu_{1,s},0)$ whether or not arm $1$ is also a leader.
By Lemma~\ref{lem:jump},
\[
  \Ix_1(t)=s\Kinf(\hnu_{1,s},0)+\log s\ge\log(Y_s/e).
\]
Hence, if $\Np_2(t)<Y_s/e$, then $\Ix_2(t)<\Ix_1(t)$ and $\Ap_t=2$ whatever the \term{w:tie}{tie rule}. Let $t_0$ be the
first round after $\tau_s$ at which arm $1$ is pulled. If $t_0<T$, then $t_0\ge2$ and
$\Np_2(t_0)\ge Y_s/e$, so $\Np_2(T)\ge Y_s/e$. If $t_0\ge T$, every round of
$\{\tau_s+1,\dots,T-1\}$ pulls arm $2$, so $\Np_2(T)\ge T-1-\tau_s\ge T-h\ge T/2$. This
proves~\eqref{eq:frozen}. Consequently, with
\[
  V_s\coloneqq\ind\{\tau_s<h,\ Y_s\ge s-1\}\min\{Y_s/e,\ T/2\}\qquad\text{and}\qquad B_T\coloneqq\max_{1\le s\le g}V_s,
\]
we have $\Np_2(T)\ge B_T$.

\emph{Step 3: a dichotomy.} Either $\Prob(\Np_2(h)>h/2)>1/2$, and then
\[
  \E \Np_2(T)\ge\E \Np_2(h)>\frac h4\ge\frac{T-1}{8};
\]
or $\Prob(\Np_1(h)\ge h/2)\ge1/2$, since $\Np_1(h)+\Np_2(h)=h$. In the second case, for every $s\le g$
(so $s\le h/2$), $\{\Np_1(h)\ge h/2\}\subseteq\{\Np_1(h)\ge s\}=\{\tau_s<h\}$, so
$\Prob(\tau_s<h)\ge1/2$.

\emph{Step 4: part (i).} Let $p_T\coloneqq\Prob(Y\ge eT/2)$ and, for $s\le g$,
$E_s\coloneqq\{\tau_s<h,\ Y_s\ge eT/2\}$. Since $s-1<T/4\le eT/2$, the event $E_s$ is contained in
$\{\tau_s<h,\ Y_s\ge s-1\}$, and on it~\eqref{eq:frozen} gives $\Np_2(T)\ge T/2$. Place ourselves in
the second case of Step 3. By Step 1, $\Prob(E_s)=\Prob(\tau_s<h)\,p_T\ge p_T/2$, and
$\Prob(E_s\cap E_{s'})\le\Prob(Y_s\ge eT/2,\ Y_{s'}\ge eT/2)=p_T^2$ for $s\ne s'$. By the Bonferroni
inequality,
\[
  \Prob\Big(\bigcup_{s\le g}E_s\Big)\ge\frac{gp_T}{2}-\frac{g^2p_T^2}{2}\ge\frac{gp_T}{4}\qquad\text{when }gp_T\le\frac12.
\]
Since $Y^+$ is integrable, $Tp_T\le\frac2e\E[Y^+\ind\{Y\ge eT/2\}]\to0$, so $gp_T\le1/2$ for large
$T$. Then
\[
  \E \Np_2(T)\ge\frac T2\cdot\frac{gp_T}{4}\ge\frac{T^2p_T}{64}\ge\frac{\kappa\,2^q}{64\,e^q}T^{2-q},
\]
using $g\ge T/8$ and $p_T\ge\kappa(eT/2)^{-q}$ when $eT/2\ge y_0$. In the first case of Step 3,
$\E \Np_2(T)\ge(T-1)/8$, which exceeds $\kappa2^qe^{-q}T^{2-q}/64$ for large $T$ because $2-q<1$.
So (i) holds for every $T$ larger than a $T_0$ that depends on $\kappa$, $q$, $y_0$ and the rate at
which $Tp_T\to0$.

\emph{Step 5: part (ii).} Fix $\varepsilon\in(0,1)$ and let $d^\star\coloneqq\log\frac1{1-\mean(P)}$,
$a_T\coloneqq(1-\varepsilon)\log T/d^\star$ and $G_T\coloneqq\{\Np_2(T)\ge a_T\}$.

\emph{(5a) A pathwise split.} We claim
\[
  \Np_2(T)\ge a_T\ind_{G_T}+(B_T-a_T)^+.
\]
On $G_T$ with $B_T\le a_T$ the right-hand side is $a_T\le \Np_2(T)$; on $G_T$ with $B_T>a_T$ it is
$B_T\le \Np_2(T)$; on $G_T^c$, $B_T\le \Np_2(T)<a_T$, so it is $0$. By Lemma~\ref{lem:explore}, almost
surely $\Np_2(T)\ge a_T$ for every large $T$, so $\Prob(G_T)\to1$.

\emph{(5b) The excess, by Lemma~\ref{lem:pairs}.} Let $s_0\coloneqq\lceil2ea_T\rceil+1$ and
$U_s\coloneqq(V_s-a_T)^+$. For large $T$, $3\le s_0\le g$ and $s_0-1\ge y_0$. Since
$(B_T-a_T)^+\ge\max_{s_0\le s\le g}U_s$, Lemma~\ref{lem:pairs} gives
\[
  \E(B_T-a_T)^+\ge\sum_{s=s_0}^g\E U_s-\sum_{s_0\le s<s'\le g}\E\min\{U_s,U_{s'}\}.
\]

\emph{(5c) The sum of the $\E U_s$.} Place ourselves in the second case of Step 3, and let
$s_0\le s\le g$. On $\{Y_s\ge s-1\}$, $Y_s/e\ge(s_0-1)/e\ge2a_T$, and $T/2\ge2a_T$ for large $T$,
so $V_s\ge2a_T$ wherever $V_s>0$ and $U_s\ge V_s/2\ge\frac1{2e}\min\{Y_s,T/2\}\ind\{\tau_s<h,\ Y_s\ge s-1\}$.
By Step 1 and Step 3,
\[
  \E U_s\ge\frac1{4e}\,\E\big[\min\{Y,T/2\}\ind\{Y\ge s-1\}\big].
\]
For $0\le v\le N$, $\min\{Y,N\}\ind\{Y\ge v\}\ge(\min\{Y,N\}-v)^+=\int_v^N\ind\{Y>y\}\,dy$, so, with
$v=s-1\ge y_0$ and $N=T/2\ge s-1$,
\[
  \E\big[\min\{Y,T/2\}\ind\{Y\ge s-1\}\big]\ge\int_{s-1}^{T/2}\Prob(Y\ge y)\,dy
  \ge\kappa\Big(\frac{1}{s-1}-\frac2T\Big).
\]
Summing, with $\sum_{s=s_0}^g\frac1{s-1}\ge\log\frac{g}{s_0-1}$ and $\sum_{s\le g}\frac2T\le\frac12$,
\[
  \sum_{s=s_0}^g\E U_s\ge\frac{\kappa}{4e}\Big(\log\frac{g}{s_0-1}-\frac12\Big).
\]

\emph{(5d) The pairs.} For $s<s'$, $\min\{U_s,U_{s'}\}\le U_s\ind\{U_{s'}>0\}$, with
$U_s\le Y_s\ind\{Y_s\ge s-1\}$ and $\{U_{s'}>0\}\subseteq\{Y_{s'}\ge s'-1\}$. The variables $Y_s$ and
$Y_{s'}$ are independent. For $v\ge y_0$,
$\E[Y\ind\{Y\ge v\}]=v\Prob(Y\ge v)+\int_v^\infty\Prob(Y>y)\,dy\le2\kappa'/v$ and
$\Prob(Y\ge v)\le\kappa'/v^2$. Hence
\[
  \sum_{s_0\le s<s'}\E\min\{U_s,U_{s'}\}\le\sum_{s\ge s_0}\frac{2\kappa'}{s-1}\sum_{s'>s}\frac{\kappa'}{(s'-1)^2}
  \le\sum_{s\ge s_0}\frac{2\kappa'^2}{(s-1)^2}\le\frac{2\kappa'^2}{s_0-2},
\]
which tends to $0$.

\emph{(5e) Conclusion.} In the second case of Step 3, (5a)--(5d) give
\[
  \E \Np_2(T)\ge a_T\Prob(G_T)+\frac{\kappa}{4e}\Big(\log\frac{g}{s_0-1}-\frac12\Big)-\frac{2\kappa'^2}{s_0-2}.
\]
Since $g\ge T/8$ and $s_0=O(\log T)$, the right-hand side is
$(1-\varepsilon)\frac{\log T}{d^\star}(1-o(1))+\frac{\kappa}{4e}\log T-O(\log\log T)$. In the first case
of Step 3, $\E \Np_2(T)\ge(T-1)/8$, which is larger for large $T$. Hence
$\liminf_T\E \Np_2(T)/\log T\ge(1-\varepsilon)/d^\star+\kappa/(4e)$; let $\varepsilon\to0$. Finally,
$d^\star=\Kinf(\delta_0,\mean(P))$ by Lemma~\ref{lem:dualfacts}(d).
\end{proof}

\begin{proof}[Proof of Corollary~\ref{cor:pareto}]
\proves{cor:pareto}
\uses{def:model, def:kinf, def:run, def:classes, def:imed, thm:counter, cite:GMS19-thm1, lem:dualfacts}
\emph{Step 1: the law.} $\Pq_{q,\varpi}$ is a law on $(-\infty,1]$. Since $\E V=q/(q-1)$, its mean is
$1-\varpi-\varpi q/(q-1)=1-\varpi(2q-1)/(q-1)$, which is in $(0,1)$ exactly when $0<\varpi<(q-1)/(2q-1)$. For
$y\ge1$, $\Prob(-X\ge y)=\varpi\Prob(V\ge y)=\varpi y^{-q}$, and $\E|X|^r<\infty$ if and only if
$\E V^r<\infty$, that is $r<q$.

\emph{Step 2: (a).} Theorem~\ref{thm:counter}(i) applies with $\kappa=\varpi$ and $y_0=1$. The regret
is $\Reg_T=\gap_2\E \Np_2(T)=m\,\E \Np_2(T)$. Given $p\in(0,1)$, choose $q\in(1+p,2)$: both
$\Pq_{q,\varpi}$ and $\delta_0$ belong to $\Sm_{1+p}$, and $\Reg_T\ge cT^{2-q}$ for large $T$ with $c>0$, so
$\Reg_T\ne o(T^{2-q})$ and \imed{} is not \term{w:consistent}{consistent} on $\Sm_{1+p}$.

\emph{Step 3: (b).} Theorem~\ref{thm:counter}(ii) applies with $\kappa=\kappa'=\varpi$ and $y_0=1$, and
gives the bound on $\liminf_T\E \Np_2(T)/\log T$, hence on $\Reg_T=m\,\E \Np_2(T)$. It remains to compute
$\Kinf(\delta_0,m;\mathcal C)$. Since $\mathcal C\subseteq\Sm$, it is at least
$\Kinf(\delta_0,m)=\log\frac1{1-m}$ (Lemma~\ref{lem:dualfacts}(d)). Conversely, for
$\varepsilon\in(0,1-m)$, the law $Q_\varepsilon\coloneqq(1-m-\varepsilon)\delta_0+(m+\varepsilon)\delta_1$
belongs to $\mathcal C$, has mean $m+\varepsilon>m$, and $\KL(\delta_0,Q_\varepsilon)=\log\frac1{1-m-\varepsilon}$,
which tends to $\log\frac1{1-m}$. Hence $\Kinf(\delta_0,m;\mathcal C)=\log\frac1{1-m}$. The last
sentence of (b) is then Theorem~\ref{cite:GMS19-thm1} with $\mathcal D=\mathcal C$, applied to the
\term{w:instance}{instance} $(\Pq_{2,\varpi},\delta_0)\in\mathcal C^{\{1,2\}}$. The class $\bigcap_{r<2}\Sm_r$ contains
$\Pq_{2,\varpi}$ by Step 1 and every law on $\{0,1\}$.
\end{proof}

\section{Sources of the techniques}\label{sec:attribution}

For each technique of the proofs, where it is used and the earliest source we found.

\begin{center}
\small
\begin{tabular}{@{}p{0.5\textwidth}p{0.44\textwidth}@{}}
\toprule
Technique (where) & Earliest source found \\
\midrule
The model of semi-bounded laws, \imed{}, the dual formula of $\Kinf$ (Definition~\ref{def:imed}, Lemma~\ref{cite:HT15-thm2}) & \citet[Algorithm~1, Theorem~2]{HT15}; the bounded-support dual: \citet{HT10} \\
Counting lemma for the minimal index (Lemma~\ref{lem:sup}(a)) & \citet[Lemma~13]{HT15} \\
Decomposition of the count at the level $\ms-\delta$ (Lemma~\ref{lem:decomp}, \eqref{eq:decomp}) & \citet[proof of Theorem~3, (29)]{HT15} \\
Counting by a supremum at a random level (Lemmas~\ref{lem:sup}(b) and~\ref{lem:termA}) & new; \citet[(25)--(26)]{HT15} count by a sum \\
Chernoff bound at the dual maximiser for a \term{w:subopt}{suboptimal} arm (Lemmas~\ref{lem:subopt} and~\ref{lem:termC}) & \citet[Proposition~10 and Lemma~6]{HT15} \\
Level shift of $\Kinf$ (Lemmas~\ref{lem:shift} and~\ref{lem:termB}) & source not known; replaces \citet[Lemma~7]{HT15} \\
Grid for the dual variable, convexity in $\lambda$ of $\E(1-\lambda Z)^\gamma$ (Lemmas~\ref{lem:grid}, \ref{lem:moments}) & \citet[proof of Proposition~11, (19)--(24)]{HT15} \\
An exponent $\gamma>1$ with $\E((1-X)/(1-\mu))^\gamma\le1$ (Lemma~\ref{lem:moments}(b)) & \citet[Proposition~11(i)]{HT15}, there under an exponential moment \\
Decomposition into big jumps and a good part (Proposition~\ref{prop:key}) & the split of a sum into its large summands and a truncated rest, classical in the large deviations of heavy-tailed sums (see also \citealp{FN71,Nag79}); new: its use inside the dual product of $\Kinf$, where one big jump costs a linear factor \\
Comparison with the last exit time (the paragraph after Proposition~\ref{prop:key}) & see also \citet{CL75,Jan86} \\
Bennett's inequality (Lemma~\ref{lem:bennett}) & \citet{B62} \\
The optimal law $G^\star$ (Lemma~\ref{lem:class}) & \citet[proof of Theorem~2(ii)]{HT15} \\
Lower bound for \term{w:ufc}{uniformly fast convergent} \term{w:learner}{learners} (Corollaries~\ref{cor:optimal}, \ref{cor:pareto}) & \citet[Theorem~1]{GMS19}; originally \citet{LR85} and \citet{BK96} \\
Frozen index after a negative jump (Lemma~\ref{lem:jump}, Theorem~\ref{thm:counter}) & new \\
Maximum from pairs (Lemma~\ref{lem:pairs}) & elementary; source not known \\
The summed bound needs three moments (Lemma~\ref{lem:three-moments}) & necessity half of a one-sided Baum--Katz theorem; see also \citet{CL75}, cited for this point by \citet[Section~9]{HT15}, and \citet{BK65} for the two-sided theorem \\
Mixture over the dual variable and maximal inequality (Lemma~\ref{lem:pa-assumptions}(a)) & \citet[Proposition~5, Appendix~B.4.1]{AJK21}, for a known moment bound (two dual variables); \citet[Appendix~A.1]{PA26}, for bounded supports; new: the time-uniform bound on $n\Kinf(\hP_n,\mean(P))$ for semi-bounded laws with a finite mean only; the assumption it checks is \citet[Assumption~2.3]{PA26} \\
Uniform mixture against the best constant-rebalanced portfolio, factor $n+1$ for two assets (Lemma~\ref{lem:pa-assumptions}(a), Step~3) & \citet[Theorem~1]{BK99}; see also \citet{CO96}, to whom they credit it, and \citet{Cov91} \\
Betting martingales for observations unbounded on one side (Lemma~\ref{lem:pa-assumptions}(a), Steps~1--2) & see also \citet[Proposition~3 in arXiv v7]{WSR24} and \citet[Section~2]{Min22} \\
\bottomrule
\end{tabular}
\end{center}


\section{Notation}\label{sec:notation}

Every symbol is a link to its definition.

\begin{center}
\small
\begin{tabular}{@{}p{0.27\textwidth}p{0.47\textwidth}p{0.2\textwidth}@{}}
\toprule
Symbol & Meaning & Defined in \\
\midrule
$\Sm$, $\Sm_r$, $\Mm_r$ & laws on $(-\infty,1]$ with finite mean; with a finite moment of order $r$; the latter without $\delta_1$ & Definition~\ref{def:model} \\
$\mean(P)$ & mean of a law $P$ & Definition~\ref{def:model} \\
$\arms$, $\Karms$ & set of arms, its size & Definition~\ref{def:model} \\
$\nu$, $\marm_a$, $\ms$, $\gap_a$ & \term{w:instance}{instance}, mean of arm $a$, best mean, gap & Definition~\ref{def:model} \\
\term{w:optimal}{optimal}, \term{w:subopt}{suboptimal} & an arm with $\gap_a=0$; with $\gap_a>0$ & Definition~\ref{def:model} \\
$\KL$, $\Kinf(P,\mu;\mathcal D)$, $\Kinf(P,\mu)$ & Kullback--Leibler divergence; its infimum over the laws of $\mathcal D$ (of $\Sm$) with mean $>\mu$ & Definition~\ref{def:kinf} \\
$\hP[x_{1:s}]$, $\bx_s$ & empirical law and mean of $x_1,\dots,x_s$ & Definition~\ref{def:kinf} \\
$\Xs_{a,s}$, $\Useed$ & $s$-th sample of arm $a$; random seed of the \term{w:learner}{learner} & Definition~\ref{def:run} \\
\term{w:learner}{learner}, $\Ap_t$, $\Ob_t$ & sequence of maps; arm pulled and observation at round $t$ & Definition~\ref{def:run} \\
$\Np_a(t)$, $\Reg_T$ & pulls of $a$ before round $t$; regret over the rounds $0,\dots,T-1$ & Definition~\ref{def:run} \\
$\hnu_{a,s}$, $\hm_{a,s}$ & empirical law and mean of the first $s$ samples of arm $a$ & Definition~\ref{def:run} \\
\term{w:anytime}{anytime} & not knowing the horizon & Section~\ref{sec:setting} \\
\term{w:consistent}{consistent}, \term{w:ufc}{uniformly fast convergent} & classes of \term{w:learner}{learners} & Definition~\ref{def:classes} \\
$\hmt_a(t)$, $\hnut_a(t)$, $\hms(t)$ & current empirical mean and law of $a$; largest empirical mean & Definition~\ref{def:imed} \\
\term{w:leader}{empirical leader}, \term{w:tie}{tie rule} & an arm with $\hmt_a(t)=\hms(t)$; a measurable choice among minimisers & Definition~\ref{def:imed} \\
$\Dv_a(t)$, $\Ix_a(t)$, $\Is(t)$ & divergence term, index of \imed{}, minimal index & Definition~\ref{def:imed} \\
$\ws_s(x;\mu)$, $\Gam(P,\mu)$ & deviation functional of a sample; expectation of its supremum & Definition~\ref{def:dev} \\
$\lst(P,\mu)$, $\Lt(P,\mu,u)$ & a dual maximiser; Chernoff rate of a \term{w:subopt}{suboptimal} arm & Definition~\ref{def:rate} \\
$\Ki_i$, $\Lt_i(\delta)$ & $\Kinf(\nu_i,\ms)$; the rate at $u=\Ki_i-\delta/(1-\ms)$ & Definition~\ref{def:rate} \\
$\Pq_{q,\varpi}$ & the law equal to $1$ with probability $1-\varpi$ and to minus a Pareto($q$) variable otherwise & Corollary~\ref{cor:pareto} \\
$\kl$ & Bernoulli Kullback--Leibler divergence & Section~\ref{sec:examples} \\
$\lmax$ & $1/(1-\mu)$ & Lemma~\ref{lem:optimizer} \\
$\TA$, $\TB$, $\TC$ & the terms of the decomposition & Lemmas~\ref{lem:decomp}, \ref{lem:termB} \\
$\ui_i$, $\MT_T$ & level of the deviations of arm $i$; pulls paid at first order & \eqref{eq:ui}, \eqref{eq:MT} \\
\bottomrule
\end{tabular}
\end{center}

\bibliographystyle{plainnat}
% refs: the works cited; ../../papers/library: the library's papers (generated; keys INC-P<nnn>).
\bibliography{refs,../../papers/library}

\end{document}
